% concatenated from main_neurips.tex, contrast_background.tex
\documentclass[letterpaper]{article}
\usepackage[a-2b,mathxmp]{pdfx}
\usepackage[textheight=9in,
    % textwidth=5.5in,
    top=1in,
    left=1in,
    right=1in,
    headheight=12pt,
    headsep=25pt,
    footskip=30pt]{geometry}

% if you need to pass options to natbib, use, e.g.:
%     \PassOptionsToPackage{numbers, compress}{natbib}
% before loading neurips_2024


% ready for submission
% \usepackage{NeurIPS/neurips_2024}


% to compile a preprint version, e.g., for submission to arXiv, add add the
% [preprint] option:
%     \usepackage[preprint]{neurips_2024}


% to compile a camera-ready version, add the [final] option, e.g.:
% \usepackage[final]{NeurIPS/neurips_2024}


% to avoid loading the natbib package, add option nonatbib:
%    \usepackage[nonatbib]{neurips_2024}


\usepackage{enumitem}

% to compile a preprint version, e.g., for submission to arXiv, add add the
% [preprint] option:
%\usepackage[preprint]{neurips_2023}
%\usepackage[preprint]{neurips_2023}

% to compile a camera-ready version, add the [final] option, e.g.:
%     \usepackage[final]{neurips_2023}


% to avoid loading the natbib package, add option nonatbib:
%    \usepackage[nonatbib]{neurips_2023}

\usepackage{xr}

% \usepackage[utf8]{inputenc} % allow utf-8 input
\usepackage{hyperref,pifont} % hyperlinks
\usepackage{url}            % simple URL typesetting
\usepackage{booktabs}       % professional-quality tables
\usepackage{amsfonts}       % blackboard math symbols
\usepackage{nicefrac}       % compact symbols for 1/2, etc.
\usepackage{microtype}      % microtypography
\usepackage[rgb]{xcolor}         % colors
\selectcolormodel{natural}
\usepackage{ninecolors}
\selectcolormodel{rgb}
\usepackage{booktabs}
\usepackage{diagbox}

\usepackage{todonotes}
\usepackage{times}
\usepackage{amsmath,color}
\usepackage{subcaption}
\usepackage{amsthm}
\usepackage{array}
% \usepackage{natbib}
\usepackage{graphicx}
\usepackage{multirow}
\usepackage{multicol}
\usepackage{caption}
\usepackage{makecell}
\usepackage{bbm}
\usepackage{breqn}
\usepackage{bm}
\usepackage{hyperref}
\captionsetup[figure]{font=small,skip=0pt}
\usepackage{algorithm}
\usepackage{algorithmic}
\usepackage{authblk}
\usepackage{wrapfig}
\usepackage{tabularray}
\usepackage{lmodern}

\newtheorem{assumption}{Assumption}
\newtheorem{property}{Properties of Contrast Functions}
\newtheorem*{property*}{Properties of Contrast Functions}
\newtheorem{remark}{Remark}
\newtheorem{theorem}{Theorem}
\newtheorem{corollary}{Corollary}[theorem]
\newtheorem{lemma}[theorem]{Lemma}
\newtheorem{proposition}[theorem]{Proposition}

\newcommand{\ba}[1]{\begin{align}#1\end{align}}
\newcommand{\bas}[1]{\begin{align*}#1\end{align*}}
\newcommand{\bu}{\pmb{u}}
\newcommand{\balpha}{\pmb{\alpha}}
\newcommand{\bab}{\pmb{b}}
\newcommand{\bx}{\pmb{x}}
\newcommand{\by}{\pmb{y}}
\newcommand{\br}{\pmb{r}}
\newcommand{\bt}{\pmb{t}}
\newcommand{\be}{\pmb{e}}
\newcommand{\bg}{\pmb{g}}
\newcommand{\bz}{\pmb{z}}
\newcommand{\bw}{\pmb{w}}
\newcommand{\bbeta}{\pmb{\beta}}
\newcommand{\bv}{\pmb{v}}
\newcommand{\bq}{\pmb{q}}
\newcommand{\bxi}{\pmb{\xi}}
\newcommand{\bvar}{\pmb{a}}
\newcommand{\nt}{\|\bt\|}
\newcommand{\nts}{\|\bt\|^2}
\newcommand{\rd}{\color{red}}
\newcommand{\bk}{\color{black}}
\newcommand{\E}{\mathbb{E}}
\newcommand{\norm}[1]{\left\|#1\right\|}
\newcommand{\rank}{\operatorname{rank}}
\newcommand{\cov}{\text{cov}}
\newcommand{\var}{\text{var}}
\newcommand\bindent{%
  \begingroup
  \setlength{\itemindent}{\myindent}
  \addtolength{\algorithmicindent}{\myindent}
}
\newcommand\eindent{\endgroup}
\newcommand\independent{\protect\mathpalette{\protect\independenT}{\perp}}
\def\independenT#1#2{\mathrel{\rlap{$#1#2$}\mkern2mu{#1#2}}}

\newlength\myindent
\setlength\myindent{2em}
 \setlength{\belowcaptionskip}{-5pt}

\newcommand{\vp}{V_{\perp}}
\newcommand{\inner}[2]{\langle #1,#2 \rangle}
\newcommand{\m}{(\mathcal{M}+\lambda_1)}
\newcommand{\csqone}{10}
\newcommand{\etakub}{\frac{1}{12}}
\newcommand{\lref}[1]{\hyperref[{#1}]{Property~\ref*{#1}}}
\newcommand{\mult}[2]{\the\numexpr #1 * #2 \relax}
\newcommand{\add}[2]{\the\numexpr #1 + #2 \relax}
\newcommand{\mydiv}[2]{\the\numexpr #1/#2 \relax}
\newcommand{\csq}{\mult{\csqone}{2}}
\newcommand{\bb}[1]{\left(#1\right)}
\newcommand{\bbb}[1]{\left[#1\right]}
\newcommand{\babs}[1]{\left|#1\right|}

\newcommand{\F}{\mathcal{F}}
\newcommand{\tr}[1]{\left\langle#1\right\rangle}
\newcommand{\bl}{\color{blue}}
\newcommand{\orange}{\color{orange}}
\newcommand{\reverseB}[1]{C{#1}}
\newcommand{\newB}[2]{B_{#2,#1}}
\newcommand{\init}{u}
\newcommand{\sgn}{\text{sgn}}
\newcommand{\bi}{\begin{itemize}}
\newcommand{\ib}{\end{itemize}}
\newcommand{\p}{\item}
\newcommand{\dmix}{d_{\text{mix}}}
\newcommand{\indep}{\perp \!\!\! \perp}
\newcommand{\M}{\mathcal{M}}
\newcommand{\B}{B}
\newcommand{\eiid}{\E_{\text{IID}}}

\DeclareMathOperator{\Tr}{Tr}
\DeclareMathOperator{\Diag}{Diag}
\DeclareMathOperator*{\argmin}{arg\,min}
\DeclareMathOperator{\Var}{Var}
\DeclareMathOperator{\gap}{gap}
\DeclareMathOperator{\sym}{sym}
\DeclareMathOperator{\vect}{vec}
\DeclareMathOperator{\diag}{diag}
\DeclareMathOperator{\poisson}{Poisson}

\title{Nonparametric Evaluation of Noisy ICA Solutions}

% The \author macro works with any number of authors. There are two commands
% used to separate the names and addresses of multiple authors: \And and \AND.
%
% Using \And between authors leaves it to LaTeX to determine where to break the
% lines. Using \AND forces a line break at that point. So, if LaTeX puts 3 of 4
% authors names on the first line, and the last on the second line, try using
% \AND instead of \And before the third author name.

\usepackage{authblk,babel}

\makeatletter
\title{Nonparametric Evaluation of Noisy ICA Solutions}
\author[1]{Syamantak Kumar\thanks{syamantak@utexas.edu}}
\affil[1]{Department of Computer Science, University of Texas at Austin}
\author[2]{Purnamrita Sarkar \thanks{purna.sarkar@utexas.edu}}
\affil[2]{Department of Statistics and Data Sciences, University of Texas at Austin}
\author[3]{Peter Bickel \thanks{bickel@stat.berkeley.edu}}
\affil[3]{Department of Statistics, University of California, Berkeley}
\author[4]{Derek Bean \thanks{derekb@stat.wisc.edu}}
\affil[4]{Department of Statistics, University of Wisconsin, Madison}

\date{}

\begin{document}
\maketitle
% \vspace{-5pt}
\begin{abstract}
% \vspace{-5pt}
Independent Component Analysis (ICA) was introduced in the 1980's as a model for Blind Source Separation (BSS), which refers to the process of recovering the sources underlying a mixture of signals, with little knowledge about the source signals or the mixing process. While there are many sophisticated algorithms for estimation, different methods have different shortcomings. In this paper, we develop a nonparametric score to adaptively pick the right algorithm for ICA with arbitrary Gaussian noise. The novelty of this score stems from the fact that it just assumes a finite second moment of the data and uses the characteristic function to evaluate the quality of the estimated mixing matrix without any knowledge of the parameters of the noise distribution. In addition, we propose some new contrast functions and algorithms that enjoy the same fast computability as existing algorithms like FASTICA and JADE but work in domains where the former may fail. While these also may have weaknesses, our proposed diagnostic, as shown by our simulations, can remedy them. Finally, we propose a theoretical framework to analyze the local and global convergence properties of our algorithms. %This framework uses quasi-orthogonalization inherently, and our results extend the classical analysis of cumulant-based objective functions to noisy ICA. 
%We demonstrate the efficacy of our algorithms via experimental results on simulated datasets.
\end{abstract}
% \vspace{-5pt}
\section{Introduction}
% \vspace{-5pt}
\label{section:introduction}
Independent Component Analysis (ICA) was introduced in the 1980's as a model for Blind Source Separation (BSS)  \cite{comon2010handbook, comon1994independent}, which refers to the process of recovering the sources underlying a mixture of signals, with little knowledge about the source signals or the mixing process. It has become a powerful alternative to the Gaussian model, leading to PCA in factor analysis. A good account of the history, many results, and the role of dimensionality in noiseless ICA can be found in~\cite{hyvarinen2001independent,auddy2023large}. In one of its first forms, ICA is based on observing $n$ independent samples from a $k$-dimensional population:
% \begin{fleqn}[0pt]
% \setlength{\abovedisplayskip}{10pt}
% \setlength{\belowdisplayskip}{10pt}
% \begin{alignat}{7}
%     &&&  \bx = B\bz + \bg \label{eq:ICA_noise_model}
% \end{alignat}
% \end{fleqn}
\ba{ \bx = B\bz + \bg \label{eq:ICA_noise_model} }
where $B\in\mathbb{R}^{d\times k}$ ($d\geq k$) is an unknown mixing matrix, $\bz\in \mathbb{R}^{k}$ is a vector of statistically independent mean zero ``sources'' or ``factors'' and $\bg\sim N(\pmb{0},\Sigma)$ is a $d$ dimensional mean zero Gaussian noise vector with covariance matrix $\Sigma$, independent of $\bz$. We denote $\diag\bb{\bvar} := \cov\bb{\bz}$ and $S := \cov\bb{\bx} = B\diag\bb{\bvar}B^{T} + \Sigma$. The goal of the problem is to estimate $B$. 
\\ \\
ICA was initially developed in the context of $\bg=0$, now called noiseless ICA. In that form, it spawned an enormous literature  \cite{389881, 542183, 599941, 650095, cardoso1999high, hyvarinen1999fast, lee1999independent, oja2000independent, hastie2002independent, chen2005ICA,bach2003kernel,hao2009kernel} and at least two well-established algorithms FASTICA~\cite{hyvarienen1997fast} and JADE~\cite{JADE1993Cardoso} which are widely used. % in engineering, physical sciences, social and biological sciences~\cite{HIMBERG20041214,Daubechies2009Fmri}. 
%\rd Independently, a version in which $B$ is $d\times k$ dimensional arose in factor analysis  from an interest to interpret the factors.\bk 
%An interesting recent paper~\cite{rohe2020vintage} gives the background with some new contributions. These models [to which our proposals] are also frequently applied to biological and social sciences literature [CITE].
The general model with nonzero $\bg$ and unknown noise covariance matrix $\Sigma$, known as \textit{noisy} ICA, has developed more slowly, with its own literature ~\cite{arora2012Quasi,NIPS2013_4d2e7bd3,NIPS2015_89f03f7d,Ilmonen2011SemiparametricallyEI,ilmonen2011semiparametrically, rohe2020vintage}. In the following sections, we will show that various ICA and noisy ICA methods have distinct shortcomings. Some struggle with heavy-tailed distributions and outliers, while others require approximations of entropy-based objectives, which have their own challenges (see eg. ~\cite{issoglio2021estimation}). Although methods for noiseless cases may sometimes work in noisy settings~\cite{NIPS2015_89f03f7d}, they are not always reliable (e.g., see Figure 1 in~\cite{NIPS2015_89f03f7d} and Theorem 2 in~\cite{chen2005ICA}). Despite the plethora of algorithms for noisy and noiseless ICA, the literature has largely been missing a diagnostic to decide which algorithm to pick for a given dataset. This is our primary goal.

% Consequently, we will show in the coming sections that different methods for ICA and noisy ICA suffer from different shortcomings. Some work poorly for heavy-tailed source distributions and outliers; some require estimation or approximation of entropy-based measures, which have their own difficulties~\cite{issoglio2021estimation}. Methods for the noiseless case may work well in the noisy setting~\cite{NIPS2015_89f03f7d}, but not always (see for eg. Figure~1 in~\cite{NIPS2015_89f03f7d} and Theorem~2 in~\cite{chen2005ICA}). % from low to moderately large sample sizes. %What makes the problem more difficult is that in a real-world scenario, there is no way of knowing what the source distributions are like, if there is noise or not, are the distributions heavy-tailed or not, and so on.  
% %Furthermore, in a real data setting, it is impossible to know what the underlying model is, and typically, algorithms designed for the noiseless setting, work sufficiently well for small to moderately large sample sizes even in the noisy case (see Figure~1 in~\cite{NIPS2015_89f03f7d}). \bk
% Despite the plethora of algorithms for noisy and noiseless ICA, the literature has largely been missing a diagnostic to decide which algorithm to pick for a given dataset. This is our main goal. 

% In this paper, \textit{our contribution} is twofold. First, we propose a diagnostic that can be used to pick the best algorithm from a set of candidate algorithms. Second, we propose some methods for noisy ICA that only require bounded second moments. These, as we will show, have their own difficulties, possibly complementary to the problems of the existing methods, which our diagnostic can remedy.

The paper is organized as follows. Section~\ref{sec:background} contains background and notation. Section~\ref{subsection:nonparametric_score} contains the independence score, its properties, and a Meta algorithm~\ref{alg:meta} that utilizes this score. %which is used later  This is later used in Section~\ref{sec:algorithms} to devise a simple algorithm, referred to as the Meta algorithm (Algorithm \ref{alg:meta}) which uses the score to pick the optimal out of a range of multiple contrast functions, without knowledge of the true mixing matrix, B. 
Section \ref{section:chf_cgf_fn} introduces new contrast functions based on the natural logarithm of the characteristic function and moment generating function of $\bx^{T}\bu$. Section~\ref{sec:convergence} presents the global and local convergence results for noisy ICA. %It contains a standalone result of independent interest about the loss landscape of noisy ICA for a large class of contrast functions in Section \ref{sec:power}. This covers all higher-order cumulants and is a generalization of the results of Delfosse et al. \cite{conf/icassp/DelfosseL96} about the local optima of the kurtosis-based contrast function for noiseless ICA. It also provides a result for local convergence of power-method-based algorithms when initialized sufficiently close to the optima. 
Section~\ref{section:experiments} demonstrates the empirical performance of our new contrast functions as well as the Meta algorithm applied to a variety of methods for inferring $B$.
% \vspace{-8pt}
\section{Background and overview}
\label{sec:background}
% \vspace{-7pt}
This section contains a general overview of the main concepts of ICA. %We start with the important topic of identifiability.

 \textbf{Notation:} We denote vectors using bold font as $\bv \in \mathbb{R}^{d}$. The dataset is represented as $\left\{\bx^{\bb{i}}\right\}_{i \in [n]}$ for $\bx^{\bb{i}} \in \mathbb{R}^{d}$. 
Scalar random variables and matrices are represented using upper case letters, respectively. 
 $\bm{I}_k$ denotes the $k$- dimensional identity matrix. Entries of vector $\bv$ (matrix $M$) are represented as $v_{i}$ ($M_{ij}$). $M(:,i)$ ($M(i,:)$) represents the $i^{\text{th}}$ column (row) of matrix $M$. $M^{\dag}$ represents the Moore–Penrose inverse of $M$. $\|.\|$ represents the Euclidean $l_{2}$ norm for vectors and operator norm for matrices. $\|.\|_{F}$ similarly represents the Frobenius norm for matrices. $\E\bbb{.}$ ($\widehat{\E}\bbb{.}$) denotes the expectation (empirical average); $\bx \indep \by$ denotes statistical independence. $\|.\|_{\psi_{2}}$ represents the subgaussian Orlicz norm (see Def. 2.5.6 in \cite{verhsynin-high-dimprob}). $X_{n} = O_{P}\bb{a_n}$ implies that $\frac{X_n}{a_n}$ is stochastically bounded i.e, $\forall \epsilon > 0, \exists $ a finite $M$ such that $\sup_{n}\mathbb{P}(\left|\frac{X_n}{a_n}\right| > M) \leq \epsilon$ (See \cite{van2000asymptotic}).
 
 %Later, for suitable choices of functions $g$, we will introduce contrast functions, $f$, denoted by $f(\bu|P)\equiv f(\bu|\bx)=\E_{\bx\sim P}[g(\bu^T\bx)]$ for any random variable $\bx$ generated from distribution $P$. %This notation essentially Let $P$ be a probability Functions $\gamma\bb{.|P}$ over probability dist  
 
\textbf{Identifiability:} 
One key issue in Eq~\ref{eq:ICA_noise_model} is the identifiability of $A:=B^{-1}$\footnote{If $B$ is not invertible, this can be replaced by the Moore-Penrose inverse.}, which holds up to permutation and scaling of the coordinates of $\bz$ if, at most, one component of $\bz$ is Gaussian (see~\cite{darmois1953analyse,skitovitch1953property}). %A version of this result summarizing other previous research can be found in~\cite{Davies2004identifiability}.
%, and the characterization of Gaussian random variables~\cite{kagan1973characterization}).  We do not pursue this here.\bk
If $d=k$ and $\Sigma$ are unknown, it is possible, as we shall see later, to identify the canonical versions of $B$ and $\Sigma$. For $d>k$ and unknown $\Sigma$, it is possible (see~\cite{Davies2004identifiability}) to reduce to the $d=k$ case. In this work, we assume $d=k$. 
%to identify $k$ and $d-k$ dimensional subspaces corresponding to $B(\bz+\bg^{(1)})\in \mathbb{R}^k$ and $\bg^{(2)}\in \mathbb{R}^{d-k}$, where $\bg^{(1)},\bg^{(2)}$ are independent Gaussians. Thus, the $d>k$ case reduces to the $d=k$ case. In this work, we assume $d=k$. 
% We shall briefly discuss the case of unknown $k$ later.
%A version with $d>k$, and where $\bz$ could be Gaussian has long been present under the name ``Factor Analysis''. This arose in the context of psychometrics where $\bz$ represented ``important'' independent factors. An interesting recent paper~\cite{rohe2020vintage} gives this background and gives some new contributions.
 For classical factor models, when $\bz$ is Gaussian, identifiability can only be identified up to rotation and scale, leading to non-interpretability.  This suggests fitting strategies focusing on recovering $B^{-1}$ through nongaussianity as well as independence. In the noisy model (Eq~\ref{eq:ICA_noise_model}), an additional difficulty is that recovery of the source signals becomes impossible even if $B$ is known exactly. This is because, the ICA models $\mathbf{x} = B(\mathbf{z} + \mathbf{r}) + \mathbf{g}$ and $\mathbf{x} = B\mathbf{z} + (B\mathbf{r} + \mathbf{g})$ are indistinguishable 
(see~\cite{DBLP:journals/corr/VossBR15} pages 2-3). This is resolved in~\cite{DBLP:journals/corr/VossBR15} via maximizing the Signal-to-Interference-plus-Noise ratio (SINR). In classical work on ICA, a large class of algorithms (see \cite{oja2000independent}) choose to optimize a measure of non-gaussianity, referred to as contrast functions.
%The essential difficulty of ICA and the related factor analysis model~\cite{rohe2020vintage} is that the matrix $B$ is not unique. If $\bz$ is Gaussian $B$ can only be identified up to rotation and scale. If $\bz$ has no Gaussian component and $\Sigma$ is known, it can be identified up to scale changes and permutations of columns. 

\bk

\vspace{.1em}
\textbf{Contrast functions:}
 Methods for estimating $B$ typically optimize an empirical contrast function.
Formally, we define a contrast function $f(\bu|P)$ by $f:\mathcal{B}_k\times \mathcal{P}\rightarrow \mathbb{R}$ where $\mathcal{B}_k$ is the unit ball in $\mathbb{R}^k$ and $\mathcal{P}$ is essentially the class of mean zero distributions on $\mathbb{R}^d$. More concretely, for suitable choices of $g$, $f(\bu|P)\equiv f(\bu|\bx)=\E_{\bx\sim P}[g(\bu^T\bx)]$ for any random variable $\bx \sim P$.
The most popular example of a contrast function is the kurtosis, which is defined as $f(\bu|P)=\E_{\bx\sim P} \left.(\bu^T\bx)^4\right/\E\bbb{(\bu^T\bx)^2}^2-3$. 
The aim is to maximize an empirical estimate, $\hat{f}(\bu|P)$.  
%Optimizing over $\bu$ essentially gives us the ``demixing'' directions. Suitable choices of $f$, referred to as ``contrast functions,'' achieve their optimum when $\bu$ coincides (up to scale) with a row of the inverse of the mixing matrix $B$; a sequential procedure typically retrieves all rows, one at a time. 
Notable contrast functions are the negentropy, mutual information (used by INFOMAX~\cite{bell1995information}), the $\tanh$ function, etc. (See \cite{oja2000independent} for further details). 

\vspace{.1em}
\textbf{Fitting strategies: }
 There are two broad categories of the existing fitting strategies.
 (I) Find $A$ such that the components of $A\bx$ are both as nongaussian and as independent as possible. See, for example, the multivariate cumulant-based method JADE~\cite{JADE1993Cardoso} and characteristic function-based methods~\cite{ERIKSSON20032195,chen2005ICA}. The latter we shall call the PFICA method.
 (II) Find successively the $j^{th}$ row of $A$, denoted by $A(j,:)\in\mathbb{R}^d$, $j=1,\dots, k$ such that $A(j,:)^T\bx$ are independent and each as nongaussian as possible, that is, estimate $A(1,:)$, project, and then apply the method again on the residuals successively. The chief representative of these methods is FastICA~\cite{hyvarinen1999fast} based on univariate kurtosis. 

 \vspace{.1em}
 \textbf{From noiseless to noisy ICA: } In the noiseless case one can first prewhiten the dataset using the empirical variance-covariance matrix, $\hat{\Sigma}$, of $\bx^{(i)}$, using $\hat{\Sigma}^{-1/2}$. Therefore, WLOG, $B$ can be assumed to be orthogonal. Searching over orthogonal matrices not only simplifies algorithm design for fitting strategy (I) but also makes strategy (II) meaningful. It also greatly simplifies the analysis of kurtosis-based contrast functions (see ~\cite{conf/icassp/DelfosseL96}). For noisy ICA, prewhitening requires knowledge of the noise covariance, and therefore many elegant methods ~\cite{arora2012Quasi,DBLP:journals/corr/VossBR15,NIPS2015_89f03f7d} avoid this step. 

% For noisy ICA, it seems that this is only possible if the covariance matrix of the Gaussian noise ($\var(\bg)$) is known. Indeed, for asymptotically vanishing Gaussian noise, Theorem~1 of ~\cite{chen2005ICA} establishes that PFICA works even in the noisy case. Finally, JADE~\cite{JADE1993Cardoso} works when $\var(\bg)$ is a scaled multiple of the identity and $d>k$. This is possible because one can estimate the scale by looking at the trailing eigenvalues of the covariance matrix of the data.
% However, the idea of quasi-orthogonalization (see~\cite{arora2012Quasi}) basically makes it possible to prewhiten using any estimate of the form $BDB^T$ where $D$ is a positive diagonal matrix. It seems that the Hessian of cumulant-based contrast functions have this form, with the catch that the diagonal of $D$ can be a mix of positive and negative numbers. As we will show later, this can be remedied by the pseudo-euclidean iteration introduced by~\cite{NIPS2015_89f03f7d}. 


\vspace{-1pt}

% \vspace{1em}
\textbf{Individual shortcomings of different contrast functions and fitting strategies:} %Unfortunately, all contrast functions fail to obey some of these conditions. 
To our knowledge, no single ICA algorithm or contrast function works universally across all source distributions. Adaptively determining which algorithm is useful in a given setting would be an important tool in a practitioner's toolbox.
Key challenges include:
%\begin{itemize}[leftmargin=5pt,align=left,labelwidth=\parindent,labelsep=2pt]

     \textbf{Cumulants:} Even though contrast functions based on \textit{even cumulants} have no spurious local optima (Theorem~\ref{theorem:global_convergence},~\cite{NIPS2015_89f03f7d}), they can vanish for some non-Gaussian distributions. Our experiments (Section~\ref{section:experiments}) show that kurtosis-based contrast functions can perform poorly even when there are a few independent components with zero kurtosis. They also suffer under heavy-tailed source distributions~\cite{rademacher2017heavy,chen2005ICA,hyvarinen1997ica}.

     \textbf{PFICA:} PFICA can outperform cumulant-based methods in some situations (Section~\ref{section:experiments}). However, it is computationally much slower, and poorly performs for Bernoulli($p$) sources with small $p$.
     
    \textbf{FastICA: } \cite{issoglio2021estimation} show that FastICA's use of Negentropy approximations for computational efficiency may result in a contrast function where optimal directions do not correspond to directions of low entropy. $\tanh$ is another popular smooth contrast function used by FastICA. However, in~\cite{wei2017fastica}, the authors show that this tends to return spurious solutions for some distributions~\cite{wei2017fastica}.
%\end{itemize}

Contrast functions are usually nonconvex and may have spurious local optima. However, they may still, under suitable conditions, converge to maximal directions. We introduce such a contrast function which vanishes iff $\bx$ is Gaussian in Section~\ref{section:chf_cgf_fn} and does not require higher moments.
%It is important to note that contrast functions are typically nonconvex and may have spurious local optima.
% Thus, even though these algorithms work well in many situations, there are competitors that work better in specific cases. 
%\\ \\ \noindent

\textbf{Our contributions:}

1) Using Theorem~\ref{theorem:independence_score_correctness} we obtain a computationally efficient scoring method for evaluating the $B$ estimated by different algorithms. Our score can also be extended to evaluate each demixing direction in sequential settings, which we do not pursue here.

2) We propose new contrast functions that work under scenarios where cumulant-based methods fail. We propose a fast sequential algorithm as in~\cite{NIPS2015_89f03f7d} to optimize these and further provide theoretical guarantees for convergence.

% \vspace{-0.3cm}
% \begin{enumerate}
%     \itemsep0em 
%     \item Using Theorem~\ref{theorem:independence_score_correctness} to obtain a scoring method for choosing between different algorithms, without sacrificing the speed too much. Our score can also be applied to choose and evaluate each demixing direction in sequential settings. However, we do not pursue this here.
%     \item Proposing new contrast functions with well-defined domains that work under scenarios where cumulant-based methods fail. In doing so, we extend the one-at-a-time approach of~\cite{NIPS2015_89f03f7d}.
% \end{enumerate}
% \vspace{-0.3cm}
 %However, there are many others based on optimizing scaled contrast functions. 
%measure of non-gaussianity is the absolute \rd univariate or multivariate \bk kurtosis of $\bz$. The univariate kurtosis of a zero mean random variable $Y$ is defined as $\kappa_{4}\bb{Y} = \E\bbb{Y^{4}} - 3\E\bbb{Y^{2}}^{2}$
 %This is the starting point of ICA, to quote Hyvarinen's famous words~\cite{hyvärinen2004independent}, ``nongaussian is independent''. The Central Limit Theorem basically says that even a linear combination of \textit{two} independent non-Gaussian random variables is more Gaussian than either random variable.
%Thus, if we can find a direction $\bu$ that maximizes the non-Gaussianity of $\bu^T \bz$, that should coincide with one of the rows of $B^{-1}$ and should give us a random variable proportional to one component of the $\bz$. 
%The algorithms for fitting $B$ are based on this identifiability property.  
%The most popular measure of non-gaussianity is the absolute \rd univariate or multivariate \bk kurtosis of $\bz$. The univariate kurtosis of a zero mean random variable $Y$ is defined as $\kappa_{4}\bb{Y} = \E\bbb{Y^{4}} - 3\E\bbb{Y^{2}}^{2}$. This is useful because of the well-known property that Gaussianity implies that all univariate cumulants (third onwards) vanish, while independence implies all multivariate cumulants beyond the third one to vanish as well. These properties are based on the quadratic nature of $\log E[\exp(t^T\bz)]$ (under Gaussianity) and additivity (under the independence of coordinates of $\bz$). These properties make kurtosis and higher cumulants an attractive choice for noisy ICA~\cite{arora2012Quasi, pmlr-v30-Belkin13, NIPS2013_4d2e7bd3, NIPS2015_89f03f7d}. %However such algorithms necessarily perform poorly for near zero or zero kurtosis and possibly heavy tails [CITE].
%However, it is also well known that for heavy-tailed data~\cite{rademacher2017heavy,chen2005ICA,hyvarinen1997ica} it may be preferable to use other contrast functions which do not rely on higher moments. %In extreme cases, the kurtosis of such random variables may not exist. On the other hand
%Moreover, examples abound of non-Gaussian random variables with zero or small kurtosis, %separating which from the Gaussian noise 
%which could elude a kurtosis-based contrast function. Following~\cite{ERIKSSON20032195}, ~\cite{chen2005ICA} proposed a criterion and an algorithm PFICA that works without these restrictions in the noiseless case. 
%Our experiments suggest that even the presence of a few independent components with zero kurtosis can derail the behavior of kurtosis-based contrast functions. In practice, one does not have a way to find which components have zero kurtosis without finding the demixing directions.
% Algorithms involving entropy require estimation of the probability density functions and typically simpler approximations are used~\cite{oja2000independent}. Furthermore, in a real data setting, it is impossible to know what the underlying model is, and typically JADE, FASTICA, INFOMAX and 
%  PFICA, which are designed for the noiseless setting, work sufficiently well for small sample settings even in the noisy case (see Figure~1 in~\cite{NIPS2015_89f03f7d}). %FastICA and JADE typically use the kurtosis for extracting the independent signals. FastICA is a fixed-point algorithm based on an approximation of Newton's method and is based on a deflationary approach to recover one component at a time. JADE, on the other hand, estimates the entire mixing matrix at once. PFICA is based on the characteristic function and uses the \textit{uncorrected} independence score (see Section \ref{subsection:nonparametric_score}).
%Following~\cite{ERIKSSON20032195}, ~\cite{chen2005ICA} proposed a criterion which works without these restrictions in the noiseless case. 
% Our contributions are
% \begin{enumerate}
%     \item An extension of the criterion proposed by~\cite{chen2005ICA} with a novel use: selection among other competing algorithms.
%     \item Proposals and algorithms for different contrast functions that enjoy the same fast computability as FASTICA or JADE but which work in domains where the former may fail. These methods may have their own weaknesses, but our first contribution will, as we will show, remedy this.
% \end{enumerate}
%\rd [R] These properties enable such algorithms to work well in both noiseless and noisy settings. \bk
% \textbf{Overview of related methods:}
% For noiseless ICA, existing results use the remarkable fact (see ~\cite{conf/icassp/DelfosseL96}) that the local maxima and minima are achieved when the demixing directions $\bu$ are aligned with the rows of the inverse of $B$, i.e. directions that extract each independent component. Furthermore, when the data is whitened (i.e. has an identity covariance matrix), $B$ can be assumed an orthogonal matrix without loss of generality, reducing the complexity in the search for all directions. This naturally leads to a simple algorithm: prewhiten the data, and then maximize the absolute value of kurtosis of $\bx^T\bu$, i.e, $f\bb{\bu}:= \kappa_{4}\bb{\bx^T\bu}$, subject to a norm constraint plus orthogonality to previously-found directions. %{\rd I'm pretty sure kurtosis is optimized by the un-mixing matrix, regardless of whitening or not. The value of whitening is in orthogonalizing the mixing matrix, making the optimization space less complex. -Derek}
% However, prewhitening is challenging for noisy ICA because the covariance matrix of the data includes that of the Gaussian noise, which is typically unknown. Of many excellent approaches that attempt to overcome this problem, we highlight ones most related to our work, beginning with \cite{arora2012Quasi}. There, the authors show that it is enough to \textit{quasi} orthogonalize the data, i.e. multiply it by a matrix which makes the data have a diagonal covariance matrix (not necessarily the identity). Subsequent work~\cite{NIPS2015_89f03f7d,NIPS2013_4d2e7bd3} uses cumulants to create a $C$ matrix of the form $BDB^T$, and then use this matrix to achieve quasi-orthogonalization. However, a difficulty is that $D$ may not be positive semidefinite (PSD)--for example, when the signal mixes positive and negative kurtosis random variables. In an earlier paper ~\cite{NIPS2013_4d2e7bd3}, Voss et al. propose a method to learn a $C$ that is PSD. However, their later work ~\cite{NIPS2015_89f03f7d} foregoes this step via a novel iterative power method, where the updates use $C$ in such a way that implicitly, the optimization happens in a pseudo-Euclidean space. Their method extracts one column of $B$ at a time.
%The methodology developed by Voss et. al. still employs the kurtosis as a contrast function. kurtosis can be an advantageous choice, for its theoretical properties which we previously noted, as well as its ease of computation. However, it is also well known that in heavy-tailed data (\rd Cite the other paper by Rademacher, aiyou also talks about heavy-tailed stuff, maybe Ericsson does too? \bk) it may be more preferable to use other contrast functions which do not rely on higher moments. In extreme cases, the kurtosis of such random variables may not exist. On the other hand, it is not hard to find examples of non-Gaussian random variables with zero or small kurtosis; directions corresponding to such random variables could elude a kurtosis-based contrast function. Therefore it is natural to ask if there is a way to distinguish between two different contrast functions in a non-parametric way. {\rd I don't understand the question in the preceding sentence. Are we talking about distinguishing between contrast functions or directions? I also don't understand how this question follows from the previous sentence. Finally I don't understand what non-parametric has to do with it. I'm probably missing something subtle and important, but I think we want to make the more generic statement that it would be good to have methods for finding and validating non-Gaussian directions that don't require computing the kurtosis; hence, the contributions of our method. -Derek} Ideally, if we knew $B$, this would be trivial, but in reality, that is the final goal. With this in mind, we state our contributions.
% The paper is organized as follows. Section~\ref{sec:background} contains background and notation. Section~\ref{subsection:nonparametric_score} contains the independence score, its properties, and a Meta algorithm~\ref{alg:meta} that utilizes this score. %which is used later  This is later used in Section~\ref{sec:algorithms} to devise a simple algorithm, referred to as the Meta algorithm (Algorithm \ref{alg:meta}) which uses the score to pick the optimal out of a range of multiple contrast functions, without knowledge of the true mixing matrix, B. 
% Section~\ref{sec:convergence} presents the global and local convergence results for noisy ICA. %It contains a standalone result of independent interest about the loss landscape of noisy ICA for a large class of contrast functions in Section \ref{sec:power}. This covers all higher-order cumulants and is a generalization of the results of Delfosse et al. \cite{conf/icassp/DelfosseL96} about the local optima of the kurtosis-based contrast function for noiseless ICA. It also provides a result for local convergence of power-method-based algorithms when initialized sufficiently close to the optima. 
% Section \ref{section:chf_cgf_fn} introduces new contrast functions based on the natural logarithm of the characteristic function and moment generating function of $\bx^{T}\bu$. Section~\ref{section:experiments} demonstrates the empirical performance of our new contrast functions as well as the Meta algorithm applied to a variety of methods for inferring $B$.  %These new contrast functions outperform kurtosis in high noise settings and the presence of zero-kurtosis signals further establishes that the Meta algorithm can adaptively choose between the appropriate contrast function and achieve scores close to the best one.
% \begin{itemize}
%     \item Section \ref{section:main_results} provides our main results. 
%     \begin{itemize}
%         \item We start by describing the loss landscape of noisy ICA for a large class of contrast functions in Section \ref{section:loss_landscape}. This covers all higher-order cumulants and is a generalization of the results of Delfosse et al. \cite{conf/icassp/DelfosseL96} about the local optima of the kurtosis-based contrast function for noiseless ICA.
%         \item Noting the limitations of kurtosis-based (or in general cumulant-based) contrast functions for heavy-tailed or zero-kurtosis data, we describe a nonparametric score in Section \ref{subsection:nonparametric_score}, which when provided with a list of candidate columns generated by different contrast functions, picks the best one. 
%         \item Section \ref{section:independence_score_sequential} explains how to use this score when extracting columns of the mixing matrix, $B$, in a sequential manner in Algorithm \ref{alg:independence}. This is later used to devise a simple algorithm, referred to as the Meta algorithm (Algorithm \ref{alg:meta}) which uses the score to pick each column iteratively when provided with multiple contrast functions.
%         \item Section \ref{subsection:chf_fn} introduces one such novel candidate contrast function based on the natural logarithm of the characteristic function of $\bx^{T}\bu$, which is free from the limitations suffered by the cumulants.
%     \end{itemize}
%     \item Section \ref{section:experiments} provides experiments to support our claims. We compare the performance of the kurtosis-based fixed-point algorithm proposed in \cite{NIPS2015_89f03f7d} against Algorithm \ref{alg:meta} on various signal distributions to show that the Meta algorithm is able to obtain performance reasonably close to the best candidate.
% \end{itemize}
% \rd Idea: put in a section here called e.g. "Description of method." Inform the reader how the algorithm works before justifying it in the main results section. Suggested structure: start with the "meta" or "outer" algorithm: suppose you have a set of candidate non-Gaussian directions. How to find the "most" non-G one? Here's our non-parametric score that we can use to check it. Don't worry, we will justify it later (that's what's in the "Main results" section). Then move on to the inner algorithm: we have our own chf-based power method that can serve as an alternative to kurtosis and we can use the meta-algorithm to distinguish between them; here is that power method. Justification in the next section. -Derek
% \vspace{-7pt}
\section{Main contributions}
\label{section:main_results}
% \vspace{-7pt}
%\subsection{Independence score in a sequential algorithm}
In this section, we present the ``Independence score'' and state the key result that justifies it. We conclude with two new contrast functions. 
%\vspace{-6pt}
% \vspace{-5pt}
\subsection{Independence score}
% \vspace{-5pt}
\label{subsection:nonparametric_score}
While there are many tests for independence~\cite{gretton2007kernelindependence,kirshner2008ICA,hao2009kernel,bach2003kernel}, we choose the Characteristic function-based one because it is easy to compute and has been widely used in normality testing~\cite{ebner2021bahadur,ebner2023eigenvalues,HENZE19971} and ICA~\cite{chen2005ICA,ERIKSSON20032195,fabian2005chf}.
\bk
Let us start with noiseless ICA to build intuition. Say we have estimated the inverse of the mixing matrix, i.e. $B^{-1}$ using a matrix $F$. Ideally, if $F=B^{-1}$, then $F\bx=\bz$, where $\bz$ is a vector of independent random variables. Kac's theorem \cite{Kac1936SurLF} tells us that
$\E\bbb{\exp(it^T\bz)}=\prod_{j=1}^k \E\bbb{\exp(it_jz_j)}$ if and only if $z_i$ are independent. In~\cite{ERIKSSON20032195}, a novel characteristic function-based objective (CHFICA), is further analyzed and studied by~\cite{chen2005ICA} (PFICA). Here one minimizes $|\E\bbb{\exp(it^TF\bx)}-\prod_{j=1}^k\E\bbb{\exp(i t_j (F\bx)_j)}|$ over $F$. We refer to this objective as the \textit{uncorrected} independence score.
%  For the noiseless case, the correct quantity to look at (also mentioned in~\cite{ERIKSSON20032195})
%  \ba{
% \max_{F}\left|\E\exp(i\bt^TF\bx)\exp\bb{-\frac{\bt^T\diag(F\Sigma F^T )\bt}{2}}-\\
% \prod_{j=1}^k\E\bbb{\exp(i t_j (F\bx)_j)}\exp\bb{-\frac{\bt^TF\Sigma F^T \bt}{2}}\right|
%  }
% However, this requires the knowledge of the noise covariance matrix. Instead, 
We propose to adapt the CHFICA objective using \textit{estimable parameters}, $S$, to the noisy ICA setting. We will minimize:
{\small
\ba{\label{eq:score}
\Delta(\bt,F|P):=\left|\underbrace{\E\bbb{\exp(i\bt^TF\bx)}\exp\bb{\frac{-\bt^T\diag(F S F^T)\bt}{2}}}_{\textsc{Joint}}\right.
\left.-\underbrace{\prod_{j=1}^k\E\bbb{\exp(i t_j (F\bx)_j)}\exp\bb{\frac{-\bt^TFS F^T\bt}{2}}}_{\textsc{Product}}\right|
}
}
where $S=\E[\bx\bx^T]$ and can be estimated using the sample covariance matrix. Hence, we do not require knowledge of any model parameters. 
We refer to this score as the (\textit{corrected}) independence score. The second terms in JOINT and PRODUCT (Eq~\ref{eq:score}) \textit{corrects} the original score by canceling out the additional terms resulting from the Gaussian noise using the covariance matrix S of the data (estimated via the sample covariance). %Eq~\ref{eq:score} ingeniously uses the form of the characteristic function of a Gaussian random variable. In the proof sketch, we build intuition with a small toy example.  
See~\cite{ERIKSSON20032195} for a related score requiring knowledge of the unknown Gaussian covariance matrix $\Sigma$. 
We are now ready to present our first theoretical result for consistency of $\Delta(\bt,F|P)$.
%\vspace{-2pt}

\begin{theorem}
\label{theorem:independence_score_correctness}
    %Recall the independence score $\Delta(\bt,F|P)$ defined in Eq~\ref{eq:score}. 
    If $F \in \mathbb{R}^{k \times k}$ is invertible and the joint and marginal characteristic functions of all independent components, $\left\{z_{i}\right\}_{i\in[k]}$, are twice-differentiable, then $\forall \bt \in \mathbb{R}^{k}$, $\Delta(\bt,F|P)=0$ \textit{iff} $F=DB^{-1}$ where $D$ is a permutation of an arbitrary diagonal matrix.
    %\bas{
%\Delta(\bt,F|P):=\E\bbb{\exp(i\bt^TF\bx)}\exp\bb{-\frac{\bt^T\diag\bb{F S F^T}\bt}{2}}\\
%-\prod_{j=1}^k\E\bbb{\exp(i t_j (F\bx)_j)}\exp\bb{-\frac{\bt^TFS F^T \bt}{2}}=0
  %  }
\end{theorem}
%\vspace{-2pt}
Unfortunately, $F$ is not uniquely defined if $\Sigma$ is unknown as we have noted in Section~\ref{section:introduction}.
The proof of Theorem~\ref{theorem:independence_score_correctness} is provided in the Appendix Section~\ref{section:proofs}. When using this score in practice, $S$ is replaced with the sample covariance matrix of the data, and the expectations are replaced by the empirical estimates of the characteristic function. The convergence of this empirical score, $\Delta(\mathbf{t},F|\hat{P})$, to the population score, $\Delta(\mathbf{t},F|P)$, is given by the following theorem.

%\rd Discussion about time taken for computation of score \bk.
% \begin{theorem} For any given demixing matrix $F = D B^{-1}$ (D being any permutation of a diagonal matrix), under suitable distributional assumptions (conditions $\bb{C1}$, $\bb{C2}$ and $\bb{C3}$ in \cite{HESSE1990186} along with finite, bounded moments as in (1.6) in \cite{vershynin2010close}), we have, 
% \bas{
%     \sup_{\mathbf{t} \in \mathbb{S}^{d-1}}\left|\hat{\Delta}_{\mathbf{t}}\bb{F} - \Delta_{\mathbf{t}}\bb{F}\right| = O_{p}\bb{\max\left\{\left\|F\right\|_{2}^2, 1\right\}}
% }
% \end{theorem}
% \begin{theorem}
%     \label{theorem:uniform_convergence}
%     Let $\mathcal{F}:=\{F\in \mathbb{R}^{k\times k}:\|F\| \leq 1\}$. Assume that $\bx\sim\text{subgaussian}(\sigma)$, i.e, $\forall\; \bv \in \mathbb{R}^{k}, \|\inner{\bv}{\bx}\|_{\psi_{2}} \leq \sigma \|S^{\frac{1}{2}}\bv\|$. Then $\forall \bt \in \mathbb{R}^{k} \; \norm{\bt}=1$, we have
%     \bas{
%      &\sup_{F\in\mathcal{F}}|\Delta(\bt,F|P)-\Delta(\bt,F|\hat{P})=O_P\bb{\sqrt{\left.k\|S\|\max(k,\sigma^{4}\|S\|)\log n\right/n}}
%     }
% \end{theorem}
\begin{theorem}
    \label{theorem:uniform_convergence}
    Let $\mathcal{F}:=\{F\in \mathbb{R}^{k\times k}:\|F\|\leq 1\}$, $\bx\sim\text{subgaussian}(\sigma)$ and $C_{k} := \max(1,
k\log\bb{n}\Tr(S))$. Then, we have 
    {\normalsize
    \bas{
    \sup_{F\in\mathcal{F}}|\E_{\bt\in N(0,I_k)}\Delta(\bt,F|P)-\E_{\bt\in N(0,I_k)}\Delta(\bt,F|\hat{P})|=O_P\bb{\sqrt{\frac{k^2\|S\|\max(k,\sigma^{4}\|S\|)\log^2 (n C_k)}{n}}}
    }}
\end{theorem}
This bound shows that uniformly over $F$, the empirical average 
$\E_{\bt\in N(0,I_k)}\Delta(\bt,F|\hat{P})$, is close to the population score. This guarantees that as long as the difference between the population scores of the two candidate algorithms is not too small, the meta-algorithm can pick up the better score.
\begin{remark}[Subgaussianity assumption]
The subgaussianity assumption in Theorem~\ref{theorem:uniform_convergence} simply ensures the concentration of the sample covariance matrix since it is used in the score (see Theorem A.3, line 825). It can be relaxed if the concentration in the operator norm is ensured. Appendix Section~\ref{subsection:supergaussian} contains experiments with more super-Gaussian source signals.
\end{remark}
The proof of this result is deferred to the Appendix section \ref{subsection:ica_uniform_convergence}. It is important to note that $\Delta(\bt,F|\hat{P})$ is not easy to optimize. As we show in Section~\ref{subsection:properties_contrast}, objective functions that are computationally more amenable to optimize for ICA, e.g., cumulants, satisfy some properties (see Assumption~\ref{assump:third}). The independence score does not satisfy the first property (Assumption~\ref{assump:third} (a)).  In the \textit{noiseless} case, whitening reduces the problem to $B$ being orthogonal. This facilitates efficient gradient descent in the Stiefel manifold of orthogonal matrices~\cite{chen2005ICA}. However, in the noisy setting, the prewhitening matrix contains the unknown noise covariance $\Sigma$, making optimization hard. Furthermore, as noted in Remark~\ref{remark:average_t}, in practice, it is better to use $\E_{\bt\sim N(0,I_k)}\Delta(\bt,F|\hat{P})$~\cite{chen2005ICA,ERIKSSON20032195,HENZE19971} instead of using a fixed vector, $\bt$, to evaluate $\Delta(\bt,F|\hat{P})$. Any iterative optimization method requires repeatedly estimating this expectation at each gradient step. Therefore, instead of using this score directly as the optimization objective, we use it to choose between different contrast functions after extracting the demixing matrix with each. This process is referred to as the Meta algorithm (Algorithm \ref{alg:meta}). \bk
%\label{sec:algorithms}
% Given a set of contrast functions, the algorithm picks one at each stage of the sequential extraction algorithm. This is possible because of the sequential computation of the independence score. We propose the following simple algorithm (Algorithm \ref{alg:meta}) to choose the best contrast function amongst a list of candidate functions.

    % \vspace{-20pt}
    % \begin{algorithm}[H]
    %     \caption{\label{alg:meta} Meta-algorithm for choosing best candidate algorithm.}
    %     \begin{algorithmic}
    %         \STATE \textbf{Input}
    %         \bindent
    %         $\text{Algorithm list } \mathcal{L}$, $\text{Data } X \in \mathbb{R}^{n \times k}$
    %         \eindent
    %         \FOR{$j$ in range[1, $\text{size}\bb{\mathcal{L}}$]}
    %             \STATE{$B_{j} \leftarrow \mathcal{L}_{j}\bb{X}$}
    %             \COMMENT{\text{Extract mixing matrix} $B_{j}$ \text{using } $j^{th}$ \text{candidate algorithm}}
    %             \STATE{$\delta_{j} \leftarrow \widehat{\E}_{\bt \sim \mathcal{N}\bb{0,\bm{I}_{k}}}\bbb{\Delta\bb{\mathbf{t}, B_{j}^{-1}|\hat{P}}}$}
    %         \ENDFOR
    %         \STATE{$i_{*} \leftarrow \argmin_{j \in [\text{size}\bb{\mathcal{L}}]}\bbb{\delta_{j}}$}
    %         \RETURN $B_{i_{*}}$
    %     \end{algorithmic}
    % \end{algorithm}
    % \vspace{-20pt}

% \begin{wrapfigure}[r]{0.5\textwidth}
%     \centering
%     % \vspace{-\baselineskip}
%     \begin{algorithm}[H]
%         \caption{\label{alg:meta} Meta-algorithm for choosing best candidate algorithm.}
%         \begin{algorithmic}
%             \STATE \textbf{Input}
%             \bindent
%             $\text{Algorithm list } \mathcal{L}$, $\text{Data } X \in \mathbb{R}^{n \times k}$
%             \eindent
%             \FOR{$j$ in range[1, $\text{size}(\mathcal{L})$]}
%                 \STATE{$B_{j} \leftarrow \mathcal{L}_{j}(X)$}
%                 \COMMENT{\text{Extract mixing matrix} $B_{j}$ \text{using } $j^{\text{th}}$ \text{candidate algorithm}}
%                 \STATE{$\delta_{j} \leftarrow \widehat{\E}_{\bt \sim \mathcal{N}(0,\bm{I}_{k})}\left[\Delta(\mathbf{t}, B_{j}^{-1}|\hat{P})\right]$}
%             \ENDFOR
%             \STATE{$i_{*} \leftarrow \argmin_{j \in [\text{size}(\mathcal{L})]}\left[\delta_{j}\right]$}
%             \RETURN $B_{i_{*}}$
%         \end{algorithmic}
%     \end{algorithm}
% \end{wrapfigure}
%\rd Add belkin's algorithm \bk

\begin{figure}[h]
    \vspace{-10pt}
    \begin{minipage}{0.45\textwidth}
        \begin{algorithm}[H]
            \caption{\label{alg:meta} Meta-algorithm for choosing best candidate algorithm.}
            \begin{algorithmic}
                \STATE \textbf{Input}: Algorithm list $\mathcal{L}$, Data $X \in \mathbb{R}^{n \times k}$
                \FOR{$j$ in range[1, $\text{size}\bb{\mathcal{L}}$]}
                    \STATE{$B_{j} \leftarrow \mathcal{L}_{j}\bb{X}$}
                    \COMMENT{\text{Extract mixing matrix} $B_{j}$ \text{using } $j^{th}$ \text{candidate algorithm}}
                    \STATE{$\delta_{j} \leftarrow \widehat{\E}_{\bt \sim \mathcal{N}\bb{0,\bm{I}_{k}}}\bbb{\Delta\bb{\mathbf{t}, B_{j}^{-1}|\hat{P}}}$}
                \ENDFOR
                \STATE{$i_{*} \leftarrow \argmin_{j \in [\text{size}\bb{\mathcal{L}}]}\bbb{\delta_{j}}$}
                \RETURN $B_{i_{*}}$
            \end{algorithmic}
        \end{algorithm}
    \end{minipage}
    \hfill
    \begin{minipage}{0.45\textwidth}
        \centering
        \includegraphics[width=\textwidth]{experiments/correlation_ind_amari_final.png}
        \caption{Correlation of independence score (with std. dev.) with Amari error between $B' = \epsilon B + \bb{1-\epsilon}I$ and $B$, averaged over 10 random runs.}
        \label{fig:amari_ind}
    \end{minipage}
\end{figure}
% \begin{wrapfigure}{l}{7 cm} % 
% \begin{minipage}{.5\textwidth}
%     \begin{algorithm}[H]
% \caption{\label{alg:meta} Meta-algorithm for choosing best candidate algorithm.}
% \begin{algorithmic}
% \STATE \textbf{Input}
% \bindent 
% $\text{Algorithm list } \mathcal{L}$, $\text{Data } X \in \mathbb{R}^{n \times k}$
% \eindent
% \FOR{$j$ in range[1, $\text{size}\bb{\mathcal{L}}$]}
%     \STATE{$B_{j} \leftarrow \mathcal{L}_{j}\bb{X}$}
%     \COMMENT{\text{Extract mixing matrix} $B_{j}$ \text{using } $j^{th}$ \text{ candidate algorithm}}
%     \STATE{$\delta_{j} \leftarrow \widehat{\E}_{\bt \sim \mathcal{N}\bb{0,\bm{I}_{k}}}\bbb{\Delta\bb{\mathbf{t}, B_{j}^{-1}|\hat{P}}}$}
% \ENDFOR
% \STATE{$i_{*} \leftarrow \argmin_{j \in [\text{size}\bb{\mathcal{L}}]}\bbb{\delta_{j}}$}
% \RETURN $B_{i_{*}}$
% \end{algorithmic}
% \end{algorithm}
% \end{minipage}
% \end{wrapfigure}
% \begin{wrapfigure}{l}{5 cm} % 
%   \begin{figure}[H]
%     % \vspace{3pt}
%     \centering
%     % \begin{subfigure}{0.48\textwidth}
%         % \centering
%         \includegraphics[width=.5\textwidth]{experiments/correlation_ind_amari_final.png}
%         \vspace{0.5pt}
%         \caption{Correlation of independence score (with std. dev.) with Amari error between $B' = \epsilon B + \bb{1-\epsilon}I$ and $B$, averaged over 10 random runs.\rd can we align\bk}
%         % \vspace{14pt}
%         \label{fig:amari_ind}
%     % \end{subfigure}
% \end{figure}
% \end{wrapfigure}
% \hspace{3pt}
% \begin{minipage}{0.45\textwidth}
% \end{minipage}
\begin{remark}[Averaging over $\bt$]\label{remark:average_t}
Algorithm~\ref{alg:meta} averages the independence score over $\bt\sim \mathcal{N}\bb{0,\bm{I}_{k}}$. While Eq~\ref{eq:score} defines the independence score for one direction $\bt$, in practice, however, there may be some directions such that a non-gaussian signal has a small score in that direction. Hence, it is desirable to average over $\bt$, following the convention in~\cite{chen2005ICA,ERIKSSON20032195,HENZE19971}.
\end{remark}
To gain an intuition of the independence score, we conduct an experiment where we use the dataset mentioned in Section \ref{section:experiments} (Figure \ref{exp:3}) and compute the independence score for $B' = \epsilon B + \bb{1-\epsilon}I, \epsilon \in \bb{0.5, 1}$. As we increase $\epsilon$, $B'$ approaches $B$, and hence the Amari error $d_{B', B}$ (see Eq~\ref{eq:amari_error}) decreases. Figure \ref{fig:amari_ind} shows the independence score versus Amari error, indicating that the independence score accurately predicts solution quality even without knowing the true $B$.
% \vspace{-6pt}
%\section{New Contrast Functions (\rd too short a section? \bk)}
%\label{section:chf_cgf_fn}
%\subsection{Desirable properties of contrast functions:}
%\input{desirable}
% \vspace{-7pt}
\subsection{Desired properties of contrast functions}
\label{subsection:properties_contrast}
% \vspace{-5pt}
The following properties are often useful for designing provable optimization algorithms for ICA.
\begin{assumption}[Properties of contrast functions]\label{assump:third}
    %We assume that $g_i
    % \vspace{-15pt}
    Let $f$ be a contrast function defined for a mean zero random variable $X$. Then, 
    \begin{enumerate}[leftmargin=20pt,align=left,labelwidth=\parindent,labelsep=2pt]
    \vspace{-5pt}
      \itemsep0em 
      \item[(a)] $f(u|X+Y)=f(u|X)+f(u|Y)$ for $X \indep Y$ 
        \item[(b)]  $f(0|X)=0, f'(0|X)=0$, $f''(0|X)=0$
        \item[(c)]  $f(u|G)=0,\forall u$ for $G\sim\mathcal{N}\bb{0,1}$
      %  \item[(d)]  The third derivative of $f'''(u|X)$ does not change the sign in the half line $[0,\infty)$ for the non-Gaussian random variable considered in the ICA problem.
        \item[(d)] WLOG, $f(u|X)=f(-u|X)$ (symmetry)
        % \item[(a)] $h_{X+Y}(u)=h_X(u)+h_Y(u)$ for independent random variables $X,Y$
        % \item[(b)]  $h_{Z}(0)=0, h'_{Z}(0)=0$, $h''_{Z}(0)=0$, for any random variable $Z$
        % \item[(c)]  $h_{G}(u)=0,\forall u$ for a mean zero Gaussian random variable.
        % \item[(d)]  The third derivative of $h_X(u)$ does not change sign in the half line $[0,\infty)$ for the non-Gaussian random variable considered in the ICA problem.
        % \item[(e)] WLOG, we assume symmetry, i.e. $h_X(u)=h_X(-u)$.
    \end{enumerate}
    %, , and $h_{G}(u)=0,\forall u$ for a mean zero Gaussian random variable.   We assume that the third derivative of $h(u)$ does not change sign in the half line $[0,\infty)$ or $(-\infty,0]$.
\end{assumption}
Properties (a) and (c) ensure that $f(u|G+X)=f(u|X)$ for non-gaussian $X$ and independent non-gaussian $G$, which means that the additive independent Gaussian noise does not change $f$. Property (c) ensures that $|f(u|G)|$ is minimal for a Gaussian.  Property (d) holds without loss of generality because one can always  symmetrize the distribution.
\footnote{Note that $y_i := x_i-x_{\lfloor n/2 \rfloor +i}, \; i = 1\cdots \lfloor n/2 \rfloor$ has a symmetric distribution, and also remains a noisy version of a linear combination of source signals with the same mixing matrix.}
% \vspace{-7pt}
\subsection{New contrast functions}\label{section:chf_cgf_fn}
% \vspace{-5pt}
In Section~\ref{sec:background}, we provided a discussion of the individual shortcomings of different contrast functions for existing contrast functions.  Before we introduce new contrast functions in this section, we revisit the algorithmic issues posed by the added Gaussian noise with unknown $\Sigma$ in Eq~\ref{eq:ICA_noise_model}.

Prewhitening the data is challenging for noisy ICA because $\E[\bx\bx^T]$ includes the unknown Gaussian covariance matrix $\Sigma$. ~\cite{arora2012Quasi} show that it is enough to \textit{quasi} orthogonalize the data, i.e., multiply it by a matrix, which makes the data have a diagonal covariance matrix (not necessarily the identity). Subsequent work~\cite{NIPS2013_4d2e7bd3, NIPS2015_89f03f7d} uses cumulants and the Hessian of $f(\bu)$ to construct a matrix $C := BDB^T$ where $D$ is a diagonal matrix, and then use this matrix to achieve quasi-orthogonalization. In general, $D$ may not be positive semidefinite (e.g. when components of $\bz$ have kurtosis of different signs). %\rd In a series of papers~\cite{NIPS2013_4d2e7bd3,NIPS2013_4d2e7bd3}, this is remedied using a novel iterative power method for successive extraction of columns of $B$. Here the updates use $C$ so that implicitly, the optimization happens in a pseudo-Euclidean space. \bk
%Optimizing the contrast function $f(\bu|P)$ with a constraint of $\left\|\bu\right\|=1$, leads to a natural fixed point algorithm 
%$\bu_{\text{opt}}$, $\bu_{\text{opt}}=\nabla f(\bu_{\text{opt}}|P)/\|\nabla f(\bu_{\text{opt}}|P)\|$. This leads to a natural fixed-point algorithm for computing $\bu_{\text{opt}}$ 
%which, given an initialisation $\bu_{0}$, proceeds as 
%$\bu_{t} \leftarrow \nabla f(\bu_{t-1}|P)/\|\nabla f(\bu_{t-1}|P)\|$.
To remedy this, in ~\cite{ NIPS2015_89f03f7d}, the authors propose computing the $C$ matrix using the Hessian of $f(\bu)$ at some fixed unit vector and then perform an elegant power method in the pseudo-Euclidean space:
\ba{
\bu_{t} \leftarrow \left.\nabla f(C^{\dag}\bu_{t-1}|P)\right/\|\nabla f(C^{\dag}\bu_{t-1}|P)\| \label{eq:power}
}
A pseudo-Euclidean space is a generalization of the Euclidean space used by \cite{DBLP:journals/corr/VossBR15}. Here the produce between vectors $\bu,\bv$ is given as $u^{\top}Av$,
%  where $A$
%  does not need to be positive definite.
% We restate their algorithm as Algorithm~\ref{alg:belkin_algo} for completeness.
\begin{algorithm}[htb]
    \caption{\label{alg:belkin_algo}Pseudo-Euclidean Power Iteration for ICA (Algorithm 2 in \cite{DBLP:journals/corr/VossBR15}) $\tilde{B}$ is the recovered matrix for the ICA model in Eq~\ref{eq:ICA_noise_model}. $\tilde{A}$ is a running estimate of  $\tilde{B}^{\dagger}$.}
    \begin{algorithmic}
        \STATE \textbf{Input}: $C \in \mathbb{R}^{k \times k}$, $\nabla f$
        \STATE $\tilde{A} \leftarrow 0$, $\tilde{B} \leftarrow 0$
        \FOR{$j$ in range[1, $k$]}
            \STATE Draw $\mathbf{u}$ uniformly at random from $\mathcal{S}^{k-1}$
            \WHILE{Convergence (up to sign)}
                \STATE $\mathbf{u} \leftarrow \tilde{B}\tilde{A}\mathbf{u}$
                \STATE $\mathbf{u} \leftarrow \nabla f\bb{C^{\dagger}\mathbf{u}}/\norm{\nabla f\bb{C^{\dagger}\mathbf{u}}}_{2}$
            \ENDWHILE
            \STATE{$\tilde{B}_{j} \leftarrow \mathbf{u}$, $\tilde{A}_{j} \leftarrow \bbb{C^{\dagger}\tilde{B}_{j}/\bb{\bb{C^{\dagger}\tilde{B}_{j}}^{\top}\tilde{B}_{j}}}^{\top}$}
        \ENDFOR
        \RETURN $\tilde{B}, \tilde{A}$
    \end{algorithmic}
\end{algorithm}

In this section, we present two new contrast functions based on the logarithm of the characteristic function (CHF) and the cumulant generating function (CGF). These do not depend on a particular cumulant, and the first only requires a finite second moment, which makes it suitable for heavy-tailed source signals.
 %Section~\ref{sec:convergence}, we presented a general framework for analyzing loss functions with pseudo-Euclidean constraints that arise in power method type inference algorithms for ICA~\cite{NIPS2015_89f03f7d}.
% While it may seem like even cumulants are the only candidates satisfying Assumption~\ref{assump:third} a), b), c), and e), there are other natural candidates like the logarithm of the Characteristic function or Cumulant generating function which also satisfy a), b), c) and e). In what follows, we show that the log of the characteristic function and the moment generating function lead to Hessians of the form $BDB^T$, and we use this for our power iteration. As shown in the experiments section, contrast functions based on the characteristic function or the cumulant generating function outperform kurtosis-based contrast functions for certain distributions involving heavy tails, or zero-kurtosis. Therefore, including them as candidates in Algorithm \ref{alg:meta} increases the robustness of the algorithm against different signal distributions. 
\bk
% As shown in the experiments section, this contrast function, albeit not necessarily with good global optima properties like the kurtosis, achieves good error rates in heavy-tailed and high-noise settings. Contrast functions using the cumulant generating function(CGF) and the logarithm of the characteristic functions satisfy the conditions a)-e) for some distributions.
% But experimentally, we saw that CGF performs poorly for heavy-tailed data and hence only used the log of CHF-based contrast function.
%\end{remark}
%\subsection{A characteristic-function based contrast function}
%Typically for noisy ICA, there has been much emphasis on cumulants \cite{arora2012provable, pmlr-v30-Belkin13, NIPS2013_4d2e7bd3, NIPS2015_89f03f7d} because they are additive, i.e. for the $i^{th}$ cumulant $\kappa_i(X+Y)=\kappa_i(X)+\kappa_i(Y)$ for two independent random variables $X$ and $Y$. Furthermore, cumulants enjoy the property that Gaussian random variables are the only ones for which all cumulants beyond the second cumulant (the variance) are zero.
We will use $f(\bu|P)$ or $f(\bu|\bx)$ where $\bx\sim P$ interchangeably to represent the population contrast function. In constructing both of the following contrast functions, we use the fact that, like the cumulants, the CGF and CHF-based contrast functions satisfy Assumption~\ref{assump:third} (a). To satisfy Assumption~\ref{assump:third} (c), we subtract out the part resulting from the Gaussian noise, which leads to the additional terms involving $\bu^{\top}S\bu$.

% \vspace{.5em}
\textbf{CHF-based contrast function:} Recall that $S$ denotes the covariance matrix of $\bx$. We maximize the \textit{absolute value} of following:
\ba{
f(\bu|P) &= \log\E\exp(i\bu^T\bx)+\log \E\exp(-i\bu^T\bx) 
         + \bu^TS\bu\notag\\ 
 &=\log(\E\cos(\bu^T\bx))^2+\log (\E\sin(\bu^T\bx))^2+\bu^TS\bu\label{eq:logchf}
}
The intuition is that this is exactly zero for zero mean Gaussian data and maximizing this leads to extracting non-gaussian signal from the data. It also satisfies Assumption~\ref{assump:third} a), b) and c).
%Note that it trivially satisfies the properties sketched in Eq~\ref{eq:properties}. 

%\rd Syamantak, please cite the other second CHF paper\bk
% \vspace{-5pt}
% \begin{theorem}
%    \label{theorem:CHF_quasi_orth}
%    Consider the data generated from the noisy ICA model (Eq~\ref{eq:ICA_noise_model}). Define $f(\bu|P)$ as in Eq~\ref{eq:logchf}. We have $\nabla^2 f(\bu|P)=BD_uB^T$, for diagonal matrix $D_{u}$.
% \end{theorem}
% \vspace{-5pt}
% \begin{proof}
% First note that 
% \begin{align*}
% g(\bu;\bx)=    \log \E\exp(i\bu^TB\bz)+\log \E\exp(-i\bu^TB\bz)+\bu^TB\diag(\bvar)B^T\bu,
% \end{align*}
%   where $\diag(\bvar)$ is the covariance matrix of $\bz$. For simplicity of notation, define $\phi(\bz;\bu):=\E\bbb{\exp(i\bu^TB\bz)}$.
%     \begin{align*}
%         \nabla g(\bu;\bx)=iB\bb{\underbrace{\frac{\E\bbb{\exp(i\bu^TB\bz)\bz}}{\phi(\bz;\bu)}}_{\mu(\bu;\bz)}-i\frac{\E\bbb{\exp(-i\bu^T\bz)\bz}}{\phi(\bz;-\bu)}}+2B\diag(\bvar)B^T\bu
%     \end{align*}
% Define
% $$H(\bu;\bz):=\frac{\E\bbb{\exp(i\bu^TB\bz)\bz\bz^T}}{\phi(\bz;\bu)}-\mu(\bu;\bz)\mu(\bu;\bz)^T.$$

% Taking a second derivative, we have:
%     \begin{align*}
%         \nabla^2 g(\bu;\bx)=B\bb{-H(\bu;\bz)-H(\bz;-\bu)+2\diag(\bvar)}B^T
%     \end{align*}
    
%   %Clearly the third term is of the form $BDB^T$ for $D=\diag(\bvar)$. We will show that 
%   All we have to show at this point is that $H(\bu;\bz)$ is diagonal. We will evaluate the $k,\ell$ entry. Let $B_i$ denote the $i^{th}$ column of $B$. Let $Y_{\setminus k,\ell}:=\bu^TB\sum_{j\neq k,\ell}z_j$. Using independence of the components of $\bz$, we have, for $k\neq \ell$,
%   \ba{\label{eq:Huz}
%   \frac{\E\bbb{\exp(i\bu^TB_kz_k+i\bu^TB_{\ell}z_{\ell}+Y_{\setminus k\ell})z_kz_{\ell}}}{\phi(\bz;\bu)}=\frac{\E\bbb{\exp(i\bu^TB_kz_k)z_k}\E\bbb{\exp(i\bu^TB_{\ell}z_{\ell})z_{\ell}}}{\E\bbb{\exp(i\bu^TB_kz_k)}\E\bbb{\exp(i\bu^TB_\ell z_\ell)}}
%   }
%   Now we evaluate:
%   \ba{\label{eq:muuz}
%   \be_k^T\mu(\bu;\bz):=\frac{\E\bbb{\exp(i\bu^TB\bz)z_k}}{\phi(\bz;\bu)}=\frac{\E\bbb{\exp(i\bu^TB_k z_k)z_k}}{\E\bbb{\exp(i\bu^TB_k z_k)}}
%   }
%   Using Eqs~\ref{eq:Huz} and~\ref{eq:muuz}, we see that indeed $H(\bu;\bz)$ is diagonal.
% \end{proof} 
\textbf{CGF-based contrast function:} The cumulant generating function also has similar properties (see~\cite{YEREDOR2000897} for the noiseless case). In this case, we maximize the \textit{absolute value} of following:
\ba{
f(\bu|P) &= \log \E\exp(\bu^T\bx) - \frac{1}{2}\bu^TS\bu \label{eq:cgf}
}
Like CHF, this vanishes iff $\bx$ is mean zero Gaussian, and satisfies Assumption~\ref{assump:third} a), b) and c). However, it cannot be expected to behave well in the heavy-tailed case. In the Appendix (Section~\ref{section:proofs},  Theorem~\ref{theorem:CGF_CHF_quasi_orth}), we show that the Hessian of both functions obtained from Eq~\ref{eq:logchf} and Eq~\ref{eq:cgf} evaluated at some $\bu$ is, in fact, of the form $BDB^T$ where $D$ is some diagonal matrix. 
% Finally, we show that the Hessian of this function evaluated at some $\bu$ is, in fact, of the form $BDB^T$ where $D$ is some diagonal matrix.

% \begin{theorem}
%    \label{theorem:CGF_quasi_orth}
%    Consider the data generated from the noisy ICA model (Eq~\ref{eq:ICA_noise_model}). Define $f(\bu|P)$ as in Eq~\ref{eq:cgf}. We have $\nabla^2 f(\bu|P)=BD_uB^T$, for diagonal matrix $D_{u}$.
% \end{theorem}
% Proof of Theorem~\ref{theorem:CGF_CHF_quasi_orth} can be found in Appendix Section~\ref{section:proofs}. 



%To use these contrast functions as candidates in Algorithm \ref{alg:meta}, we follow Algorithm 2 in~\cite{NIPS2015_89f03f7d}. 
%We estimate $C_i=\nabla^2 g(\br_i|P)$ using its empirical version at some randomly picked unit vector $\br_i$ for $i=1,\dots,k$. Let $C_i^{\dag}$ denote the pseudo-inverse of $C_i$. 
 %To estimate $B$ using the contrast functions in Eq~\ref{eq:logchf} and~\ref{eq:cgf}, we first estimate a $C$ matrix, and then apply Algorithm~2 in~\cite{NIPS2015_89f03f7d}. 
% in~\cite{voss} We denote by $C^{\dag}$ the pseudo-inverse of $C$. Next, we update two matrices:  one that estimates the mixing matrix $B$ one column at a time ($U$) and one that estimates $A=B^{-1}$ \textit{one row} at a time ($V$). For estimating column $i$ (under some permutation) of $B$, we first choose an initial random unit vector $\bu^{(i)}_0$, and then do a power iteration using $C$ with $\bu\rightarrow C^{\dag}\bb{I-UV^T}\bu_t^{(i)}$. Let $M(:,i)$ denote the $i^{th}$ column of $M$ and $M(j,:)$ denote the $j^{th}$ row of $M$. We update $U$ as $U(:,i)=\bu^{(i)}$ and $V(j,:)$ as $\bv^{(j)}=\frac{C^{\dag}\bu^{(j)}}{{\bu^{(j)}}^T C^{\dag}\bu^{(j)}}$.
% \ba{
% \bu_{t+1}^{(i)}=\frac{\nabla g(C_i^{\dag}\bb{I-UV^T}\bu_t^{(i)};\bx)}{\|\nabla g(C_i^{\dag}\bu_t^{(i)},\bx)\|}
% }
% Upon convergence, as shown in Eqs~\ref{eq:projectionfirst}, we update $U$ and $V$ as follows:
% % Upon convergence, we update two matrices, one that estimates the mixing matrix $B$ one column at a time and one that estimates $B^{-1}$ \textit{one row} at a time. \rd Please check this.\bk
% \ba{\label{eq:projection}
% U(:,i)=\bu^{(i)}\qquad V(i,:)=\frac{C_i^{\dag}\bu^{(i)}}{{\bu^{(i)}}^TC_i^{\dag}\bu^{(i)}}
% }

% \vspace{-7pt}
\section{Global and local Convergence: loss landscape of noisy ICA} \label{sec:convergence}
% \vspace{-7pt}
Here we present sufficient conditions for global and local convergence of a broad class of contrast functions. This covers the cumulant-based contrast functions and our CHF and CGF-based proposals. \bk

% \vspace{-7pt}
\subsection{Global convergence}
\label{subsection:global_convergence}
% \vspace{-7pt}
%\section{Characterizating the Loss landscape of noisy ICA and a new contrast function}\label{sec:power}
%\rd Cumulants enjoy Properties ~\ref{item:2} to ~\ref{item:4}.\bk
% \begin{fleqn}[0pt]
% \setlength{\abovedisplayskip}{10pt}
% \setlength{\belowdisplayskip}{10pt}
% \begin{alignat}{2}
%     & f(\bu|\bx+\by) = f(\bu|\bx)+f(\bu|\by), \text{ when $\bx \indep \by$} \label{eq:properties} \\
%     & f(\bu|\bg) = 0, \text{ when $\bg$ is mean zero gaussian}
% \end{alignat}
% \end{fleqn}
% \ba{
%     & f(\bu|\bx+\by) = f(\bu|\bx)+f(\bu|\by), \text{when $\bx \indep \by$} \label{eq:properties} \\
%     & f(\bu|\bg) = 0, \text{when $\bg$ is mean zero gaussian}
% }
%where $\bx \indep \by$ denotes statistical independence. 
%The classical proof of the global optima of the kurtosis or cumulant-based loss function was done by~\cite{conf/icassp/DelfosseL96} in the non-noisy case. %Here, the objective function is $\max |\kappa_{4}\bb{\bx^{T}\bu}| \; s.t. \; \bu^T\bu=1$.
The classical argument for global optima of kurtosis for noiseless ICA in~\cite{conf/icassp/DelfosseL96} assumes WLOG that $B$ is orthogonal. This reduces the optimization over $\bu$ into one for $\bv=B\bu$. Since $B$ is orthogonal, $\|\bu\|=1$ translates into $\|\bv\|=1$. It may seem that this idea should extend to the noisy case due to property~\ref{assump:third} a) and c) by optimizing over $\bv=B\bu$ over potentially non-orthogonal mixing matrices $B$.

However, for non-orthogonal $B$, the norm constraint on $\bv$ is no longer $\bv^T\bv=1$, \textit{rendering the original argument invalid}. In what follows, we extend the original argument to the noisy ICA setting by introducing pseudo-euclidean constraints. Our framework includes a large family of contrast functions, including cumulants. 

To our knowledge, this is the first work to characterize the loss functions in the pseudo-euclidean space. In contrast,~\cite{NIPS2013_4d2e7bd3} provides global convergence of the cumulant-based methods by a convergence argument of the power method itself.

Consider the contrast function $f\bb{\bu|P} := \E_{\bx\sim P}\bbb{g\bb{\bx^{T}\bu}}$ and recall the definition of the quasi orthogonalization matrix $C = BDB^{T}$, where $D$ is a diagonal matrix. For simplicity, WLOG let us assume that $D$ is invertible. 
% \bas{
%     f\bb{C^{\dag}\bu} &= \kappa_{4}\bb{\bu^{T}C^{\dag}\bx} \\ 
%                       &= \kappa_{4}\bb{\bu^{T}\bb{B^{T}}^{-1}D^{\dag}\bz + \bu^{T}C^{\dag}\bg} \text{ using } \ref{eq:ICA_noise_model} \\ 
%                       &= \sum_{i \in [k]}\kappa_{4}\bb{\frac{\bb{B^{-1}\bu}_{i}}{D_{ii}}\bz_{i}} \\ 
%                       &= \sum_{i \in [k]}\kappa_{4}\bb{\frac{\alpha_i}{D_{ii}}\bz_{i}}
% }
%When applied to each independent component $X$, this we define $h_X(u):=\E_X\bbb{g(X;u)}$ for a function satisfying Eq~\ref{eq:properties}. For example, this could be the kurtosis of the random variable $uX$. Typically this is applied to a $k$ dimensional vector via the kurtosis of $\bx^T\bu$. %We will denote that by $f(\bu^T\bx)$.
We now aim to find the optimization objective that leads to the pseudo-Euclidean update in Eq~\ref{eq:power}.  For $f\bb{C^{-1}\bu|P} = \E_{\bx\sim P} \bbb{g\bb{\bu^{T}C^{-1}\bx}}$, consider the following:
\ba{
    &f\bb{C^{-1}\bu|P} = \sum_{i \in [k]}\E\bbb{g\bb{\bb{B^{-1}\bu}_{i}z_{i}/D_{ii}}} = \sum_{i \in [k]}\E \bbb{g\bb{\alpha_i\tilde{z}_i}} 
    = \sum_{i \in [k]}f\bb{\alpha_i/D_{ii}|z_i} \label{eq:powerobj}
                      %\sum_{i \in [k]}h_{z_i}\bb{\frac{\alpha_i}{D_{ii}}}  
}
% \begin{fleqn}[0pt]
% \setlength{\abovedisplayskip}{0pt}
% \setlength{\belowdisplayskip}{0pt}
% \begin{alignat}{2}
% f\bb{C^{-1}\bu|P} &= \sum_{i \in [k]}\E \bbb{g\bb{\frac{\bb{B^{-1}\bu}_{i}}{D_{ii}}\bz_{i}}} =
%                     \\&
%                     = \sum_{i \in [k]}\E \bbb{g\bb{\alpha_i\tilde{z}_i}}\nonumber = 
%                    \sum_{i \in [k]}f\bb{\frac{\alpha_i}{D_{ii}}|z_i}%\sum_{i \in [k]}h_{z_i}\bb{\frac{\alpha_i}{D_{ii}}}
%                    \label{eq:powerobj}    
% \end{alignat}
% \end{fleqn}
% \ba{
%     % &f\bb{C^{-1}\bu|P} = \E_{\bx\sim P} \bbb{g\bb{\bu^{T}C^{-1}\bx}} \notag\\
%     % &=
%     %                  \E \bbb{g\bb{\bu^{T}\bb{B^{T}}^{-1}D^{-1}\bz + \bu^{T}C^{-1}\bg}} \text{ using } \ref{eq:ICA_noise_model} \notag\\ 
%                       f\bb{C^{-1}\bu|P} &= \sum_{i \in [k]}\E \bbb{g\bb{\frac{\bb{B^{-1}\bu}_{i}}{D_{ii}}\bz_{i}}} =
%                        \sum_{i \in [k]}\E \bbb{g\bb{\alpha_i\tilde{z}_i}}\nonumber \\&=
%                        \sum_{i \in [k]}f\bb{\frac{\alpha_i}{D_{ii}}|z_i}%\sum_{i \in [k]}h_{z_i}\bb{\frac{\alpha_i}{D_{ii}}}
%                        \label{eq:powerobj}
% }
where we define $\balpha := B^{-1}\bu$ and $\tilde{z}_i=z_i/D_{ii}$. 
We now examine $f(C^{-1}\bu)$ subject to the ``pseudo'' norm constraint $\bu^{T}C^{-1}\bu=1$. The key point is that for suitably defined $f$, one can construct a matrix $C=BDB^T$ from the data even when $B$ is unknown. So our new objective is to optimize
\bas{
 f\bb{C^{-1}\bu|P} \text{ s.t. } \bu^TC^{-1}\bu=1
}
%Note that optimizing this with a Lagrange multiplier, and
%setting the gradient of this to zero naturally gives us the direction 
 Using the Lagrange multiplier method, optimizing the above simply gives the power method update (Eq~\ref{eq:power})
$\bu\propto\nabla f(C^{-1}\bu)$. Furthermore, optimizing Eq~\ref{eq:powerobj} leads to the following transformed objective: 
% \rd [removed first eqn and equivalent sign]
\ba{
%&\max_{\bu} |f\bb{C^{-1}\bu}|\quad\text{s.t.}\ \ \bu^TC^{-1}\bu=1\\
\max_{\balpha} \left|\sum_i f\bb{\left.\alpha_i/D_{ii}\right|z_i}\right|\quad\text{s.t.}\ \ \sum_i \alpha_i^2/D_{ii}=1 \label{eq:transformed_objective}
}
Now we provide our theorem about maximizing $f(C^{-1}\bu|P)$. Analogous results hold for minimization. We will denote the corresponding derivatives evaluated at $t_{0}$ as $f'(t_0|z_i)$ and $f''(t_0|z_i)$. 
\begin{theorem}
\label{theorem:global_convergence} %h_i''(u_i)
Let $C$ be a matrix of the form $BDB^T$, where $d_{i} := D_{ii} = f''(u_i|z_i)$ for some random  $u_i$. Let $S_+=\{i:d_i>0\}$. Consider a contrast function $f:\mathbb{R}\rightarrow \mathbb{R}$. Assume that for every non-Gaussian independent component, Assumption~\ref{assump:third} holds and the third derivative, $f'''(u|X)$, does not change the sign in the half-lines $[0,\infty)$, $(\infty,0]$. Then $f\bb{\left.C^{-1}\bu\right|\bx}$ with the constraint of $\langle\bu,\bu\rangle_{C^{-1}}=1$ has local maxima at $B^{-1}\bu=\be_i,i\in S_+$. %Solutions with $B^{-1}\bu=\be_i,i\in S_-$ are local minima. 
 All solutions satisfying $\be_i^T B^{-1}\bu\neq 0,\ \ \forall i \in [1,k]$ are minima, and all solutions with $|\left\{i : \be_i^TB^{-1}\bu\neq 0\right\}| < k$ are saddle points. These are the only fixed points.
\end{theorem}
Note that the above result immediately implies that for $\tilde{C}=-C$, the same optimization in Theorem~\ref{theorem:global_convergence}
will have maxima at all $\bu$ such that $B^{-1}\bu=\be_i,i\in S_-$.
% \begin{remark}
% Note that the above result immediately implies that for $\tilde{C}=-C$, the same optimization in Theorem~\ref{theorem:global_convergence}
% will have maxima at all $\bu$ such that $B^{-1}\bu=\be_i,i\in S_-$.
% \end{remark}
\begin{corollary}
    $2k^{th}$ cumulant-based contrast functions for $k>1$, have maxima and minima at the directions of the columns of $B$, provided $\left\{\bz_{i}\right\}_{i \in [k]}$ have a corresponding non-zero cumulant.
\end{corollary}
\vspace{-10pt}
\begin{proof}
This proof (see Appendix Section~\ref{section:proofs})  is immediate since cumulants satisfy (a), (c). For (c) and (d), note that $f(u|Z)$ is simply $u^{2j}\kappa_{2j}(Z)$, where $\kappa_{2j}(Z)$ is the $2j^{th}$ cumulant of $Z$.
\end{proof}
%\vspace{.5pt}
% \begin{proof}
% \end{proof}
%\begin{remark}
Theorem~\ref{theorem:global_convergence} can be easily extended to the case where $C$ is not invertible. The condition on the third derivative, $f'''\bb{u|X}$, may seem strong, and it is not hard to construct examples of random variables $Z$ where this fails for suitable contrast functions (see~\cite{doi:10.1080/00031305.1994.10476058}). However, for such settings, it is hard to separate $Z$ from a Gaussian. See the Appendix~\ref{subsection:necessary_condition} for more details.
% \vspace{-5pt}
\subsection{Local convergence}
% \vspace{-5pt}
In this section, we establish a local convergence result for the power iteration-based method described in Eq~\ref{eq:power}. Define $\forall t \in \mathbb{R}, \forall i \in [d]$, the functions $q_{i}\bb{t} := f'\bb{\left.\frac{\alpha_i}{D_{ii}}\right|z_i}$.
% Define $\forall \; x \in \mathbb{R}, \forall \; t = \bb{t_1, t_2, \cdots t_d} \in \mathbb{R}^{d}$,
% \bas{
%     & q_{i}\bb{x} := h_{z_{i}}'\bb{\frac{\alpha_i}{D_{ii}}} \in \mathbb{R}, \\
%     & w\bb{t} := [q_{1}\bb{t_1}, q_{2}\bb{t_2} \cdots q_{d}\bb{t_d}]^{T} \in \mathbb{R}^{d}
% }
% Then, Eq~\ref{eq:power} can be rewritten in terms of $\balpha := B^{-1}\bu$ as 
% \bas{
%     & B\balpha_{t+1} = \frac{Bw\bb{\balpha_{t}}}{\left\|Bw\bb{\balpha_{t}}\right\|}  \Rightarrow \forall i \in [d],  \bb{\balpha_{t+1}}_{i} = \frac{q_{i}\bb{\bb{\balpha_{t}}_{i}}}{\left\|Bw\bb{\balpha_{t}}\right\|}
% }
% WLOG let the fixed point be denoted as, $\balpha^{*} = \frac{\be_{1}}{\|B\be_{1}\|}$. Then, $\frac{Bw\bb{\balpha^{*}}}{\left\|Bw\bb{\balpha^{*}}\right\|_{2}} = B\be_{1}$. Therefore, 
% \bas{
%     & \forall i \in [2,d], \bb{\balpha_{t+1}}_{i} - \alpha^{*}_{i} = \bb{\balpha_{t+1}}_{i} = \frac{q_{i}\bb{\bb{\balpha_{t}}_{i}}}{\left\|Bw\bb{\balpha_{t}}\right\|} \\ 
%     & \text{ for } i = 1, \bb{\balpha_{t+1}}_{1} - \alpha^{*}_{1} = \frac{q_{1}\bb{\bb{\balpha_{t}}_{i}}}{\left\|Bw\bb{\balpha_{t}}\right\|} - \frac{1}{\|B\be_{1}\|}
% }
% \rd Define $\balpha_*$ \bk
Then, without loss of generality, we have the following result for the first corner:
\begin{theorem}
    \label{theorem:local_convergence}
    Let $\alpha_i=B^{-1}\bu_i$ and $\balpha^*:=\be_1$. Let the contrast function $f(.|X)$, satisfy assumptions \ref{assump:third} (a), (b), (c), and (d). Let $\mathcal{B} := [-\|B^{-1}\|_{2},\|B^{-1}\|_{2}]$ and define $c_{1}, c_{2}, c_{3} > 0$ such that $\forall i \in [d]$, $\sup_{x \in \mathcal{B}}\left\{\frac{|q_{i}\bb{x}|}{c_{1}}, \frac{|q_{i}'\bb{x}|}{c_{2}}, \frac{|q_{i}''\bb{x}|}{c_{3}}\right\} \leq 1$.
    % \begin{enumerate}
    %     % \vspace{-0.4cm}
    %     \itemsep0em 
    %     \item $$ 
    %     % \item $\sup_{x \in [-\|B^{-1}\|_{2},\|B^{-1}\|_{2}]}|q_{i}'\bb{x}| \leq c_{2} \label{assumption:4}$
    %     % \item $\sup_{x \in [-\|B^{-1}\|_{2},\|B^{-1}\|_{2}]}|q_{i}''\bb{x}| \leq c_{3} \label{assumption:5}$
    %     % \vspace{-0.4cm}
    % \end{enumerate}
    Define $\epsilon := \frac{\norm{B}_{F}}{\norm{B\mathbf{e}_{1}}^{2}}\max\left\{\frac{c_{3}+c_{2}\norm{B\mathbf{e}_{1}}}{\left|q_{1}\bb{\alpha^{*}_{1}}\right|}, \frac{c_{2}^{2}}{\left|q_{1}\bb{\alpha^{*}_{1}}\right|^{2}}, \frac{c_{3}\norm{B\mathbf{e}_{1}}}{\left|q_{1}'\bb{\alpha^{*}_{1}}\right|}\right\}$.
    Let $\|\balpha_{0}-\balpha^{*}\|_{2} \leq R \label{assumption:2}$, $R \leq \max\left\{\frac{c_{2}}{c_{3}}, \frac{1}{5\epsilon}\right\}$. Then, we have $\forall t \geq 1, \; \dfrac{\left\|\balpha_{t+1}-\balpha^{*}\right\|}{\left\|\balpha_{t}-\balpha^{*}\right\|} \leq \frac{1}{2}$.
\end{theorem}
% \vspace{.5pt}
Therefore, we can establish linear convergence without the third derivative constraint required for global convergence (see Theorem~\ref{theorem:global_convergence}), by assuming some mild regularity conditions on the contrast functional and a sufficiently close initialization. The proof is included in Appendix Section \ref{subsection:ica_local_convergence}.
\begin{remark}
   Let $\balpha=B^{-1}\bu$ denote the demixed direction, $\bu$.  Theorem~\ref{theorem:local_convergence} shows that if the initial $\bu_0$ is such that $\balpha_0$ is close to the ground truth $\balpha^*$ (which is WLOG taken to be $\be_1$), then there is geometric convergence to $\balpha^*$. $\left|q_1(\alpha_1^*)\right|$ and $\left|q_1'(\alpha_1^*)\right|$ essentially quantify how non-gaussian the corresponding independent component is since they are zero for a mean zero Gaussian. When these are large quantities, $\epsilon$ is small, and the initialization radius $\|\balpha_0-\balpha^*\|_2$ can be relatively larger. 
\end{remark}
% \vspace{-10pt}
\section{Experiments}
% \vspace{-5pt}
\label{section:experiments}
In this section, we provide experiments to compare the fixed-point algorithms based on the characteristic function (CHF), the cumulant generating function (CGF) (Eqs~\ref{eq:logchf} and~\ref{eq:cgf}) with the kurtosis-based algorithm (PEGI-$\kappa_{4}$ \cite{NIPS2015_89f03f7d}). We also compare against noiseless ICA\footnote{MATLAB implementations (under the GNU General Public License) can be found at - \href{https://research.ics.aalto.fi/ica/fastica/}{FastICA} and \href{https://github.com/ruohoruotsi/Riddim/blob/master/MATLAB/jade.m}{JADE}. The code for PFICA was provided on request by the authors.} algorithms - FastICA, JADE, and PFICA. These six algorithms are included as candidates for the Meta algorithm (see Algorithm \ref{alg:meta}). We also present the \textit{Uncorrected Meta} algorithm (\textit{Unc.\ Meta}) to denote the Meta algorithm with the uncorrected independence score proposed in CHFICA ~\cite{ERIKSSON20032195}. Table~\ref{table:kurtosis} and Figure~\ref{fig:plots_exp} provide experiments on simulated data, and Figure~\ref{figure:image_demixing_experiment} provides an application for image demixing \cite{deVito_ICA_for_demixing_images}. We provide additional experiments on denoising MNIST images in the Appendix Section~\ref{section:ica_additional_experiments}.

\textbf{Experimental setup:}
Similar to~\cite{NIPS2015_89f03f7d}, the mixing matrix $B$ is constructed as $B = U\Lambda V^{T}$, where $U,V$ are $k$ dimensional random orthonormal matrices, and $\Lambda$ is a diagonal matrix with $\Lambda_{ii}\in[1,3]$. The covariance matrix $\Sigma$ of the noise $\bg$ follows the Wishart distribution and is chosen to be $ \frac{\rho}{k} RR^{T}$, where $k$ is the number of sources, and $R$ is a random Gaussian matrix. Higher values of the \textit{noise power} $\rho$ can make the noisy ICA problem harder. Keeping $B$ fixed, we report the median of 100 random runs of data generated from a given distribution (different for different experiments). The quasi-orthogonalization matrix for CHF and CGF is initialized as $\hat{B}\hat{B}^T$ using the mixing matrix, $\hat{B}$, estimated via PFICA. The performance of CHF and CGF is based on a single random initialization of the vector in the power method (see Algorithm~\ref{alg:belkin_algo}). Our experiments were performed on a Macbook Pro M2 2022 CPU with 8 GB RAM.

\textbf{Error Metrics:} Due to the inherent ambiguity in signal recovery discussed in Section~\ref{sec:background}, we measure error using the accuracy of estimating $B$. %Several metrics are available for evaluating the goodness of the estimated $\hat{B}$. %Examples include mean absolute deviation ~\cite{hyvarinen1999fast} uses the mean absolute deviation of the obtained signal with the true signal.~\cite{650095, chen2005ICA} use the Amari error.~\cite{rademacher2017heavy} uses the Frobenius norm of the matrix difference, along with perfect matching, to align the columns of the estimated and true mixing matrices. \cite{NIPS2015_89f03f7d} use the SINR metric and show that it is equivalent to the expected mean-squared loss of the estimated and true signal values. 
We report the widely used Amari error~\cite{amari1995new,hastie2002ICA,chen2005ICA,chen2006ICA} for our results. For estimated demixing matrix $\widehat{B}$, define $W = \widehat{B}^{-1}B$, after normalizing the rows of $\widehat{B}^{-1}$ and $B^{-1}$. Then we measure $d_{\widehat{B}, B}$ as - 
\ba{\label{eq:amari_error}
    %d\bb{\widehat{B}, B}
    d_{\widehat{B}, B}&:= \frac{1}{k}\bb{\sum_{i=1}^{k}\frac{\sum_{j=1}^{k}|W_{ij}|}{\max_{j}|W_{ij}|} + \sum_{j=1}^{k}\frac{\sum_{i=1}^{k}|W_{ij}|}{\max_{i}|W_{ij}|}} - 2 
}
 % \vspace{.1em}
 % \\
\textbf{Variance reduction using Meta:} We first show that the independence score can also be used to pick the best solution from many random initializations. In Figure~\ref{fig:plots_exp} (a), the top panel has a histogram of 40 runs of CHF, each with one random initialization. The bottom panel shows the histogram of 40 experiments, where, for each experiment, the \textit{best out of 30 random initializations} are picked using the independence score. The top and bottom panels have (mean, standard deviation) (0.51, 0.51) and (0.39,0.34) respectively.  This shows a reduction in variance and overall better performance. 
% \begin{figure}
%     \centering
%     % \raggedleft
%     \includegraphics[width=0.6\textwidth]{histogram_one_vs_best.png}
%     % \vspace{2pt}
%     \caption{Histograms of Amari error for the experimental setup used in Figure~\ref{fig:plots_exp} with $n=10K$ and noise power $\rho=0.2$. The top panel contains a histogram of 40 runs with one random initialization. The bottom panel contains a histogram of 40 runs, each of which is the best independence score out of 30 random initializations.}
%     \label{fig:enter-label2}
%      % \vspace{1em}
% \end{figure}
% \vspace{.5em}
% \\ \noindent


\textbf{Effect of varying kurtosis:}
For our second experiment (see Table~\ref{table:kurtosis}), we use $k=5$ independent components, sample size $n=10^5$ and noise power $\rho=0.2$ from a Bernoulli($p$) distribution, where we vary $p$ from 0.001 to $0.5-1/\sqrt{12}$. The last parameter makes kurtosis zero.
\begin{table*}[!hbt]
    % \centering
    \hspace{-1cm}
 %    \begin{tabular}{|c|c|c|c|c|c|c|c|c|}\hline
 %    \vspace{1pt}
 %   % \textbf{Algorithm} & \textbf{p = 1e-3} & \textbf{p = 5e-3} & \textbf{p = 1e-2} & \textbf{p = 5e-2} & \textbf{p = 1e-1} & \textbf{p = 1.5e-1} & \textbf{p} = $\mathbf{\frac{1}{2} - \sqrt{\frac{1}{12}}}$ \\
 %    \textbf{Algorithm} & \textbf{$\kappa_4$ = 994} & \textbf{194} & \textbf{95.01} & \textbf{15.05} & \textbf{5.11} & \textbf{1.84} & \textbf{0} \\
 %    \hline
 %    \hline
 %    \vspace{1pt}
 %    \textbf{Meta} & \textbf{\rd 0.00678 \bk} & \textbf{\rd 0.01035 \bk } & \textbf{\rd 0.0106 \bk} & \textbf{\rd 0.01006 \bk} & \textbf{\rd 0.01057 \bk} & \textbf{\rd 0.01124 \bk} & \textbf{\rd 0.02335 \bk} \\\hline
 %    \textbf{CHF} & 1.52357 & 0.33621 & \textbf{0.01067} & \textbf{0.01014} & \textbf{0.01069} & \textbf{0.01124} & \textbf{0.02942} \\\hline
 %    \textbf{CGF} & \textbf{0.00678} & \textbf{0.01064} & \textbf{0.01083} & \textbf{0.01613} & 0.02887 & 0.04439 & 0.07121 \\\hline
 %    \textbf{Kurtosis} & \textbf{0.00709} & \textbf{0.01049} & \textbf{0.01077} & \textbf{0.01042} & \textbf{0.01213} & \textbf{0.0165} & 1.80196 \\\hline
 %    \textbf{PFICA} & 1.52466 & 0.8851 & 0.54019 & 0.0241 & 0.02326 & 0.02281 & 0.02363 \\\hline
 %    \textbf{JADE} & 0.02132 & 0.02207 & 0.02128 & 0.0217 & \text{0.02213} & 0.02317 & 1.90863 \\\hline
 %    \textbf{FastICA} & 0.02407 & 0.02721 & 0.02609 & 0.0263 & 0.026 & 0.02658 & 5 \\\hline
 % %   \textbf{Uncorrected Meta} & 1.5185 & 0.04969 & 0.04112 & 0.04082 & 0.04133 & 0.04143 & 0.04186 \\\hline
 %    \end{tabular}
   \begin{tabular}{|c|c|c|c|c|c|c|c|c|c|c|}\hline
    \vspace{1pt}
   % \textbf{Algorithm} & \textbf{p = 1e-3} & \textbf{p = 5e-3} & \textbf{p = 1e-2} & \textbf{p = 5e-2} & \textbf{p = 1e-1} & \textbf{p = 1.5e-1} & \textbf{p} = $\mathbf{\frac{1}{2} - \sqrt{\frac{1}{12}}}$ \\
   \diagbox[width=9em]{\textbf{Algorithm}}{\textbf{Scaled} $\kappa_4$}
     & \textbf{994} & \textbf{194} & \textbf{95} & \textbf{15} & \textbf{5} & \textbf{2} & \textbf{0.8} & \textbf{0.13} & \textbf{0} \\
    \hline
    \hline
    \vspace{1pt}
    \textbf{Meta} & \textbf{\rd 0.007 \bk} & \textbf{\rd 0.010 \bk } & \textbf{\rd 0.011 \bk} & \textbf{\rd 0.010 \bk} & \textbf{\rd 0.011 \bk} & \textbf{\rd 0.011 \bk} & \textbf{\rd 0.0128 \bk} & \textbf{\rd 0.01981 \bk} & \textbf{\rd 0.023 \bk} \\\hline
    \textbf{CHF} & 1.524 & 0.336 & \textbf{0.011} & \textbf{0.010} & \textbf{0.011} & \textbf{0.011} & \textbf{0.0129} & \textbf{0.0213} & 0.029 \\\hline
    \textbf{CGF} & \textbf{0.007} & 0.011 & \textbf{0.011} & 0.016 & 0.029 & 0.044 & 0.05779 & 0.06521 & 0.071 \\\hline
    \textbf{PEGI} & \textbf{0.007} & \textbf{0.010} & \textbf{0.011} & \textbf{0.010} & 0.012 & 0.017 & 0.02795 & 0.13097 & 1.802 \\\hline
    \textbf{PFICA} & 1.525 & 0.885 & 0.540 & 0.024 & 0.023 & 0.023 & 0.0212 & 0.0224 & \textbf{0.024} \\\hline
    \textbf{JADE} & 0.021 & 0.022 & 0.021 & 0.022 & \text{0.022} & 0.023 & 0.029 & 0.089 & 1.909 \\\hline
    \textbf{FastICA} & 0.024 & 0.027 & 0.026 & 0.026 & 0.026 & 0.027 & 0.02874 & 0.0703 & - \\\hline
   \textbf{Unc. Meta} & 1.52 & 0.049 & 0.041 & 0.0408 & 0.0413 & 0.0414 & 0.0413 & 0.0416 & 0.0419 \\\hline
    \end{tabular}
    \caption{Median Amari error with varying $p$ in Bernoulli$(p)$ data. The scaled kurtosis is given as $\kappa_{4} := (1-6p\bb{1-p})/(p\bb{1-p})$. Observe that the Meta algorithm (shaded in red) performs at par or better than the best candidate algorithms. FastICA did not converge for zero-kurtosis data. }
    \vspace{-0.5pt}
    % Furthermore, using the Meta algorithm with the uncorrected independence score leads to a significant increase in the Amari error. }
    \label{table:kurtosis}
    % \vspace{1em}
\end{table*}
Different algorithms perform differently (best candidate algorithm is highlighted in bold font) for different $p$. In particular, characteristic function-based methods like PFICA and CHF perform poorly for small values of $p$. We attribute this to the fact that the characteristic function is close to one for small $p$. Kurtosis-based algorithms like PEGI, JADE, and FASTICA perform poorly for kurtosis close to zero. Furthermore, the Uncorrected Meta algorithm performs worse than the Meta algorithm since it shadows PFICA.
% \vspace{.5em}
% \\ \\ \noindent

\textbf{Effect of varying noise power:}
For our next two experiments, we use $k=9$, out of which 3 are from $\text{Uniform}\bb{-\sqrt{3},\sqrt{3}}$, 3 are from Exponential$(5)$ and 3 are from $\bb{\text{Bernoulli}\bb{\frac{1}{2} - \sqrt{\frac{1}{12}}}}$ and hence have zero kurtosis. In these experiments, we vary the sample size (Figure \ref{exp:3}), $n$ fixing $\rho=0.2$, and the noise power (Figure \ref{exp:2}) , $\rho$ fixing $n=10^5$, respectively. Note that with most mixture distributions, it is easily possible to have low or zero kurtosis. We include such signals in our data to highlight some limitations of PEGI-$\kappa_{4}$ and show that Algorithm \ref{alg:meta} can choose adaptively to obtain better results.
% \begin{table*}[!hbt]
%     \centering
%     \begin{tabular}{|c|c|c|c|c|c|c|c|c|}\hline
%     \vspace{1pt}
%     \textbf{Algorithm} & \textbf{p = 1e-3} & \textbf{p = 5e-3} & \textbf{p = 1e-2} & \textbf{p = 5e-2} & \textbf{p = 1e-1} & \textbf{p = 1.5e-1} & \textbf{p} = $\mathbf{\frac{1}{2} - \sqrt{\frac{1}{12}}}$ \\
%     \hline
%     \hline
%     \vspace{1pt}
%     \textbf{CHF} & 1.474 & 0.318 & 0.0106 & 0.0105 & \textbf{0.011} & \textbf{0.0124} & \textbf{0.0427} \\\hline
%     \textbf{CGF} & 0.00703 & \textbf{0.0108} & 0.01107 & 0.01704 & 0.03148 & 0.04905 & 0.08087 \\\hline
%     \textbf{Kurtosis} & 0.00717 & 0.0109 & 0.01074 & 0.01087 & 0.01265 & 0.01777 & 1.95977 \\\hline
%     \textbf{PFICA} & 1.51392 & 0.97486 & 0.55423 & 0.04905 & 0.04634 & 0.04527 & 0.04405 \\\hline
%     \textbf{JADE} & 0.0432 & 0.04309 & 0.04285 & 0.04326 & 0.04343 & 0.04422 & 1.9026 \\\hline
%     \textbf{FastICA} & 0.04699 & 0.04918 & 0.04854 & 0.04826 & 0.0482 & 0.04889 & 3.48953 \\\hline
%     \textbf{Uncorrected Meta} & 1.51392 & 0.40536 & 0.09164 & 0.04509 & 0.04479 & 0.0444 & 0.04405 \\\hline
%     \textbf{Meta} & \textbf{0.00687} & 0.01233 & \textbf{0.01055} & \textbf{0.01035} & \textbf{0.011} & \textbf{0.0124} & 0.04405 \\\hline
%     \end{tabular}
%     \caption{Average Amari error with $p$-parameter for Bernoulli distributions. Observe that the Meta algorithm performs at par or better than the best candidate algorithms. Furthermore, using the Meta algorithm with the uncorrected independence score leads to a significant increase in the Amari error.}
%     \label{exp:1}
% \end{table*}
\begin{figure}[!hbt]
    \centering
    \begin{subfigure}[t]{0.43\textwidth}
        \centering
        \includegraphics[width=\textwidth]{images/noise_variation_UM.png}
        \caption{Variation of performance with noise power}
        \label{exp:2}
    \end{subfigure}
    \begin{subfigure}[t]{0.43\textwidth}
        \centering
        \includegraphics[width=\textwidth]{images/samplesize_UM.png}
        \caption{Variation of performance with sample size}
        \label{exp:3}
    \end{subfigure}
    \\[1em] % Adds a vertical gap
    \begin{subfigure}[t]{0.48\textwidth}
        \centering
        \includegraphics[width=\textwidth]{histogram_one_vs_best.png}
        \caption{Histograms of Amari error with $n=10^4$ and noise power $\rho=0.2$}
        \label{exp:1}
    \end{subfigure}
    \vspace{5pt}
    \caption{\label{fig:plots_exp} Amari error in the $\log$-scale on $y$ axis and varying noise powers (for $n=10^5$) and varying sample sizes (for $\rho=0.2$) on $x$ axis for figures \ref{exp:2} and \ref{exp:3} respectively. For figure~\ref{exp:1}, the top panel contains a histogram of 40 runs with one random initialization. The bottom panel contains a histogram of 40 runs, each of which is the best independence score out of 30 random initializations.}
\end{figure}
Figure~\ref{exp:2} shows that large noise power leads to worse performance for all methods. PEGI, JADE, and FastICA perform poorly consistently, because of some zero kurtosis components. CGF suffers because of heavy-tailed components. However, CHF and PFICA perform well consistently. The Meta algorithm mostly picks the best algorithm except for a few points, where the difference between the two leading algorithms is small. The Uncorrected Meta algorithm (which uses the independence score without the additional correction term introduced for noise) performs identically to PFICA.
 % \vspace{1em}
%\\ \\

\textbf{Effect of varying sample size:} Figure~\ref{exp:3} shows that the noiseless ICA methods have good performance at smaller sample sizes. However, they suffer a bias in performance compared to their noiseless counterparts as the sample size increases. The Meta algorithm can consistently pick the best option amongst the candidates, irrespective of the distribution, leading to significant improvements in performance. What is interesting is that up to a sample size of $10^4$, PFICA dominates other algorithms and Meta performs like PFICA. However, after that CHF dominates, and one can see that Meta starts to have a similar trend as CHF. We also see that the Uncorrected Meta algorithm (which uses the independence score without the additional correction term introduced for noise) has a near identical error as PFICA and has a bias for large $n$.
% \section{Experiments}
% \label{section:experiments}
% In this section, we provide experiments to compare the fixed-point algorithms based on the Characteristic function (CHF) (Section \ref{subsection:chf_fn}) with the kurtosis-based algorithm (PEGI-$\kappa_{4}$) (\cite{NIPS2015_89f03f7d}) and the Meta algorithm (Algorithm \ref{alg:meta}) using both of these as candidates. We avoid comparisons with other similar algorithms, such as FastICA or JADE, since these have been shown in \cite{NIPS2015_89f03f7d} to have been outperformed by PEGI-$\kappa_{4}$ and have been designed for noiseless ICA. We generate simulated data comprising of $k=11$ different source signals. 8 of these signals follow standard distributions - Exponential(1), Gamma(4,2), Poisson(1), Logistic(2,4), Bernoulli(0.5), Uniform(-$\sqrt{3}, \sqrt{3}$), Lognormal(0,1), Weibull(1,1). The remaining three distributions are constructed to have zero-kurtosis. We use a Bernoulli with $p = \frac{1}{2} + \sqrt{\frac{1}{12}}$, a 50:50 of $\sqrt{X_{1}}$ and $\sqrt{X_{2}}$, both sampled from $\poisson\bb{\frac{1}{2}}$ and a 50:50 mixture of $\mathcal{U}\bb{-1,1}$ and $\mathcal{U}\bb{-a,a}$ for $a = \sqrt{5 + \sqrt{24}}$. 

% Note that with most mixtures distributions, it is easily possible to have low or zero kurtosis. We include such signals in our data to highlight some of the limitations of PEGI-$\kappa_{4}$ and show that Algorithm \ref{alg:meta} can choose adaptively to obtain better results. The mixing matrix is constructed, similar to \cite{NIPS2015_89f03f7d}, using a reverse singular matrix decomposition, $B = U\Lambda V^{T}$. The matrices $U$ and $V$ are random orthonormal matrices, and $\Lambda$ is a diagonal matrix with minimum one and maximum three as its singular value. The covariance matrix $\Sigma$ of the noise $\eta$ follows the Wishart distribution and is chosen to be $ \frac{\rho}{k} UU^{T}$, where $k$ is the number of sources, and $U$ is a Gaussian random matrix. Here $\rho$ controls the noise level, and higher values of $\rho$ would typically make the noisy ICA problem harder. The quasi orthogonalization matrix, $C = BDB^{T}$, for PEGI-$\kappa_{4}$ is computed at $\{\be_{i}\}_{i=1}^{d}$ deterministically since they can obtain an interpretable closed-form expression for the Hessian, and guarantee that the diagonal elements $D_{ii}$ are non-zero if the signals have non-zero kurtosis. For the case of the CHF, since we do not have such a form (see Section \ref{subsection:chf_fn}), we instead compute the Hessian at a random unit vector before extracting each column. Furthermore, for CHF, we perform 10 random runs and select the one with the lowest independence score. Each data-point in our experiments is averaged over 100 randomly generated $B$ matrices and datasets.

% There are several metrics available for the noisy ICA problem.~\cite{hyvarinen1999fast} uses the mean absolute deviation of the obtained signal with the true signal. \cite{650095, chen2005ICA} use the Amari error. \cite{rademacher2017heavy} uses the Frobenius norm of the matrix difference, along with perfect matching, to align the columns of the estimated and true mixing matrices. \cite{NIPS2015_89f03f7d} use the SINR metric and show that it is equivalent to the expected mean-squared loss of the estimated and true signal values. We use a modified version of the metric proposed in \cite{777510}. They measure the quality of a single estimated column of the mixing matrix, $\bu$, as -
% $\begin{aligned}    \mathcal{S}(\bu) = \max_{i \in [d]}\left|\left(\frac{B^{-1}\bu}{\|B^{-1}\bu\|_{2}}\right)_{i}\right| \end{aligned}$.
% If $\bu$ coincides with a column of B, then the score, $\mathcal{S}$, will be 1, and otherwise, it will always be between 0 and 1. However, they only measure the loss for a single column. We instead propose to match each estimated column $\bu$ with the true column using $\mathcal{S}(\bu)$ and find a maximum perfect matching to get a comprehensive measure of the performance across all columns. We report the mean of the matching score for each column as our final metric and use the Hungarian algorithm for perfect matching. Additionally, to make our metric robust to noise, we first perform the computation for a random unit vector and then use that as a threshold to only consider high-quality matches. We will denote this by the \textit{Corner score}.

% Our first experiment in table \ref{table:varying_samplesize} varies the sample size for a fixed-noise level and reports the Corner score. We observe that as $n$ increases, the performance of the CHF-based algorithm increases. The Meta algorithm also follows this general trend and performs significantly better than PEGI-$\kappa_{4}$, which seems to stay roughly at the same level of performance.
% % \vspace{-1em}
% \begin{table}[H]
%     \centering
%     \begin{tabular}{|c|c|c|c|c|c|}\hline
%     \vspace{1pt}
%     \textbf{Algorithm} & \textbf{n = 1K} & \textbf{n = 5K} & \textbf{n = 10K} & \textbf{n = 15K} & \textbf{n = 20K}\\
%     \hline
%     \hline
%     \vspace{1pt}
%         \textbf{CHF} &.32$\pm$ .07&.36$\pm$.08&.38$\pm$.07&.39$\pm$.09&.4$\pm$.06  \\\hline
%         \vspace{1pt}
%         \textbf{Meta} &.32$\pm$.1&.31$\pm$.1&.35$\pm$.1&$.31\pm$.06&.36$\pm$.08\\\hline
%         \vspace{1pt}
%         \textbf{PEGI}-\bm{$\kappa_{4}$} &.24$\pm$.08&.22$\pm$.08&.24$\pm$.08&.21$\pm$.09&.21$\pm$.07\\\hline
%     \end{tabular}
%     \caption{Mean and std. error of the Corner score as $n$ grows for fixed noise level $\rho=.1$.}
%     \label{table:varying_samplesize}
% \end{table}
% \vspace{-1em}
%  Table \ref{table:varying_noise} varies the noise levels for a fixed sample size and reports the Corner score. We observe that, in general, the performance of the algorithms either stays the same or decreases with increasing noise, keeping in mind the standard errors. Again, the performance of the Meta algorithm appears to be \textit{very close to the performance of CHF} and \textit{significantly better} than PEGI-$\kappa_{4}$. The deterioration of the Meta algorithm with increasing noise is not surprising. Figure~\ref{fig:online_pca_comparison1} shows that the independence score worsens as we provide bad estimates of the correct column. As noise increases, the quality of all candidate contrast functions deteriorates, making the Meta algorithm suffer more.
% \begin{table}[H]
%     \centering
%     \begin{tabular}{|c|c|c|c|c|c|}
%     \hline
%     \vspace{1pt}
%     \textbf{Algorithm} & \bm{$\rho$} = \textbf{.01} & \bm{$\rho$} = \textbf{.05} & \bm{$\rho$} = \textbf{.1} & \bm{$\rho$} = \textbf{.3} & \bm{$\rho$} = \textbf{.5}\\\hline \hline
%     \vspace{1pt}
%         \textbf{CHF} &.32$\pm$ .07&.36$\pm$.08&.38$\pm$.07&.39$\pm$.09&.4$\pm$.06  \\\hline
%         \vspace{1pt}
%         \textbf{Meta} &.33$\pm$.1&.37$\pm$.09&.32$\pm$.12&.29$\pm$.1&.3$\pm$.11\\\hline
%         \vspace{1pt} \textbf{PEGI}-\bm{$\kappa_{4}$}&.22$\pm$.07&.23$\pm$.07&.25$\pm$.07&.19$\pm$.06&.22$\pm$.08\\\hline
%     \end{tabular}
%     \caption{Mean and std. error of the Corner score for fixed $n=20000$, and growing noise levels.}
%     \label{table:varying_noise}
% \end{table}
% \vspace{-2em}
% Although the CHF-based method performs well for the above two cases, it need not always be the case. Our next experiment considers six signals and compares three different settings - (i) six signals sampled from $\mathcal{U}\bb{-\sqrt{3},\sqrt{3}}$, (ii) A mixture of three zero-kurtosis Bernoulli signals and three $\mathcal{U}\bb{-\sqrt{3},\sqrt{3}}$ signals and finally (iii) six zero-kurtosis Bernoulli signals. For $X \sim \mathcal{U}\bb{-\sqrt{3},\sqrt{3}}$, the symmetrized characteristic function is given as $f\bb{t} := \frac{\sin^{2}\bb{\sqrt{3}t}}{3t^{2}}$, which is ill-defined at infinitely many values of $t$. Therefore, in this case, the CHF-based algorithm is prone to failure. Similarly, for the case of the zero-kurtosis signals, PEGI-$\kappa_{4}$ does not work. These observations are evident in Table \ref{table:zk_uniform_mix}.
% Remarkably, the Meta algorithm always has corner scores close to the best algorithm.
% \begin{table}[H]
%     \centering
%     \begin{tabular}{|c|c|c|c|}
%     \hline
%     \vspace{1pt}
%     \textbf{Algorithm} & \textbf{6 uniform} & \textbf{3 Uniform and 3 Zero-kurtosis} & \textbf{6 Zero-kurtosis} \\\hline \hline
%     \vspace{1pt}
%         \textbf{CHF} &.56$\pm$ .12 & .54 $\pm$ .13 & .52$\pm$ 0.16  \\\hline
%         \vspace{1pt}
%         \textbf{Meta} &.98$\pm$.06 & .52 $\pm$ .13 & .47$\pm$ 0.13 \\\hline
%         \vspace{1pt} \textbf{PEGI}-\bm{$\kappa_{4}$}&.99$\pm$.01&  .17$\pm$0.11 &.21$\pm$.13\\\hline
%     \end{tabular}
%     \caption{Mean and std. error of the Corner score for fixed $n=20000$, and noise level $\rho = 0.1$ with varying source signal distributions.}
%     \label{table:zk_uniform_mix}
% \end{table}

\begin{figure}[h!]
    \centering
    \begin{subfigure}[b]{\textwidth}
        \centering
        \begin{tabular}{cccc}
            \includegraphics[width=0.2\textwidth]{demixing_exp/Original_Images/source1.png} &
            \includegraphics[width=0.2\textwidth]{demixing_exp/Original_Images/source2} &
            \includegraphics[width=0.2\textwidth]{demixing_exp/Original_Images/source3} &
            \includegraphics[width=0.2\textwidth]{demixing_exp/Original_Images/source4} \\
        \end{tabular}
        \caption{Source Images}
        \vspace{5pt}
    \end{subfigure}
    
    \begin{subfigure}[b]{\textwidth}
        \centering
        \begin{tabular}{cccc}
            \includegraphics[width=0.2\textwidth]{demixing_exp/Mixed_Images/mixed_image_1.png} &
            \includegraphics[width=0.2\textwidth]{demixing_exp/Mixed_Images/mixed_image_2} &
            \includegraphics[width=0.2\textwidth]{demixing_exp/Mixed_Images/mixed_image_3} &
            \includegraphics[width=0.2\textwidth]{demixing_exp/Mixed_Images/mixed_image_4} \\
        \end{tabular}
        \caption{Mixed Images}
        \vspace{5pt}
    \end{subfigure}
    
    \begin{subfigure}[b]{\textwidth}
        \centering
        \begin{tabular}{cccc}
            \includegraphics[width=0.2\textwidth]{demixing_exp/CHF/source1_chf.png} &
            \includegraphics[width=0.2\textwidth]{demixing_exp/CHF/source2_chf} &
            \includegraphics[width=0.2\textwidth]{demixing_exp/CHF/source3_chf} &
            \includegraphics[width=0.2\textwidth]{demixing_exp/CHF/source4_chf} \\
        \end{tabular}
        \caption{Demixed Images using CHF-based contrast function ($\Delta(\bt,F|P) = 5.26 \times 10^{-3}$).}
        \vspace{5pt}
    \end{subfigure}
    
    \begin{subfigure}[b]{\textwidth}
        \centering
        \begin{tabular}{cccc}
            \includegraphics[width=0.2\textwidth]{demixing_exp/Kurtosis/source1_kurtosis.png} &
            \includegraphics[width=0.2\textwidth]{demixing_exp/Kurtosis/source2_kurtosis} &
            \includegraphics[width=0.2\textwidth]{demixing_exp/Kurtosis/source3_kurtosis} &
            \includegraphics[width=0.2\textwidth]{demixing_exp/Kurtosis/source4_kurtosis} \\
        \end{tabular}
        \caption{Demixed Images using Kurtosis-based contrast function ($\Delta(\bt,F|P) = 2.48 \times 10^{-2}$)}
    \end{subfigure}
    \vspace{0.1pt}
    \caption{\label{figure:image_demixing_experiment} We demix images using ICA by flattening and linearly mixing them with a $4 \times 4$ matrix $B$ (i.i.d entries $\sim \mathcal{N}(0,1)$) and Wishart noise ($\rho = 0.001$). The CHF-based method (c) recovers the original sources well, upto sign. The Kurtosis-based method (d) fails to recover the second source. This is consistent with its higher independence score. The Meta algorithm selects CHF from candidates CHF, CGF, Kurtosis, FastICA, and JADE. Appendix Section~\ref{section:ica_additional_experiments} provides results for other contrast functions and their independence scores.}
\end{figure}

% \vspace{-10pt}
\section{Conclusion}
% \vspace{-7pt}
ICA is a classical problem that aims to extract independent non-Gaussian components. The vast literature on both noiseless and noisy ICA introduces many different inference methods based on a variety of contrast functions for separating the non-Gaussian signal from the Gaussian noise. Each has its own set of shortcomings. We aim to identify the best method for a given dataset in a data-driven fashion. In this paper, we propose a nonparametric score, which is used to evaluate the quality of the solution of any inference method for the noisy ICA model. Using this score, we design a Meta algorithm, which chooses among a set of candidate solutions of the demixing matrix. We also provide new contrast functions and a computationally efficient optimization framework. While they also have shortcomings, we show that our diagnostic can remedy them. We provide uniform convergence properties of our score and theoretical results for local and global convergence of our methods. Simulated and real-world experiments show that our Meta-algorithm matches the accuracy of the best candidate across various distributional settings.

% \vspace{-7pt}
\section*{Acknowledgments}
% \vspace{-7pt}
The authors thank Aiyou Chen for many important discussions that helped shape this paper and for sharing his code of PFICA. The authors also thank Joe Neeman for sharing his valuable insights and the anonymous reviewers for their helpful suggestions in improving the exposition of the paper. SK and PS were partially supported by NSF grants 2217069, 2019844, and DMS 2109155.

% \rd Complete the checklist at the end of the paper. \bk

%\clearpage
% Acknowledgements should only appear in the accepted version.
% \section*{Acknowledgements}
% The authors thank Aiyou Chen for many important discussions that helped shape this paper and for sharing his code of PFICA. The authors also thank Joe Neeman for sharing his valuable insights. SK and PS were partially supported by NSF grants 2217069, 2019844, and DMS 2109155.


% In the unusual situation where you want a paper to appear in the
% references without citing it in the main text, use \nocite
% \nocite{langley00}
% \newpage
% \section{Impact Statement}
% This paper presents work whose goal is to advance the field of Machine Learning. There are many potential societal consequences of our work, none of which we feel must be specifically highlighted here.


%%%%%%%%%%%%%%%%%%%%%%%%%%%%%%%%%%%%%%%%%%%%%%%%%%%%%%%%%%%%
\newpage 
\bibliography{refs}
\bibliographystyle{plain}

\newpage 
\appendix

\section{Appendix}

\renewcommand{\thetheorem}{A.\arabic{theorem}}
\renewcommand{\thesection}{A.\arabic{section}}
\renewcommand{\thefigure}{A.\arabic{figure}}
\renewcommand{\theequation}{A.\arabic{equation}}
\newcounter{lemma}
\newcounter{proposition}

\setcounter{theorem}{0}
\setcounter{section}{0}
\setcounter{figure}{0}
\setcounter{lemma}{0}
\setcounter{proposition}{0}

The Appendix is organized as follows - 

\begin{itemize}
    \item Section \ref{section:proofs} proves Theorems \ref{theorem:independence_score_correctness}, \ref{theorem:uniform_convergence}, \ref{theorem:CGF_CHF_quasi_orth},   \ref{theorem:global_convergence} and \ref{theorem:local_convergence} 
    \item Section~\ref{section:ica_additional_experiments} provides additional experiments for noisy ICA
    \item Section \ref{section:sequential_algorithm} provides the algorithm to compute independence scores via a sequential procedure
    \item Section \ref{section:Assumption_1d} explores the third derivative constraint in Theorem~\ref{theorem:global_convergence} and provides examples where it holds
    \item Section \ref{section:surface_plots} contains surface plots of the loss landscape for noisy ICA using the CHF-based contrast functions
\end{itemize}

\section{Proofs}
\label{section:proofs}
\vspace{5pt}
% \begin{theorem}
%     % \label{theorem:independence_score_correctness}
%     If $F$ is invertible then, the following holds \textit{iff} $F=DB^{-1}$ where $D$ is a permutation of an arbitrary diagonal matrix.
%     \bas{
%     \Delta_F(\bt):=\E\bbb{\exp(i\bt^TF\bx)}\exp\bb{-\frac{\bt^T\diag\bb{F S F^T}\bt}{2}}-\prod_{j=1}^k\E\bbb{\exp(i t_j (F\bx)_j)}\exp\bb{-\frac{\bt^TFS F^T \bt}{2}}=0
%     }
% \end{theorem}

\begin{proof}[Proof of Theorem \ref{theorem:independence_score_correctness}]

We will give an intuitive argument for the easy direction for $k=2$. 
We will show that if $F=DB^{-1}$ for some permutation of a diagonal matrix $D$, then after projecting using that matrix, the data is of the form $\bz+\bg'$ for some Gaussian vector $\bg'$.  Let $k=2$.
Let the entries of this vector be $z_1+g_1'$ and $z_2+g_2'$, where $z_i$, $i\in\{1,2\}$ are mean zero \textit{independent} non-Gaussian random variables and $g'_i$, $i\in\{1,2\}$ are mean zero \textit{possibly dependent} Gaussian variables. The Gaussian variables are independent of the non-Gaussian random variables. Let $t_i$, $i\in\{1,2\}$ be arbitrary but fixed real numbers.  Let $\bt=(t_1, t_2)$ and let $\Sigma_{g'}$ denote the covariance matrix of the Gaussian $g'$. Assume for simplicity $\var(z_i)=1$ for $i\in\{1,2\}$. Denote by $\Lambda:=\cov(F\bx)=FSF^T=I+\Sigma_{g'}$. The \textsc{Joint} part of the score is given by:
\bas{
\text{\textsc{Joint}}%&=\E\bbb{\exp\bb{i\bb{t_1(z_1+g_1)+t_2(z_2+g_2)}}}\exp\bb{-\frac{\bt^T\diag(S)\bt}{2}}\\
&=\E\bbb{it_1z_1}\E\bbb{it_2z_2}\exp\bb{-\frac{\bt^T(\Sigma_{g'}+\diag(\Lambda))\bt}{2}}
}
The \textsc{Product} part of the score is given by
\bas{
\text{\textsc{Product}}%&=\E\bbb{\exp\bb{i\bb{t_1(z_1+g_1)}}}\E\bbb{\exp\bb{t_2(z_2+g_2)}}\exp\bb{-\frac{\bt^TS\bt}{2}}\\
&=\E\bbb{it_1z_1}\E\bbb{it_2z_2}\exp\bb{-\frac{\bt^T(\diag(\Sigma_{g'})+\Lambda)\bt}{2}}
}
% Note that $\Sigma_g+\diag(S)=\Sigma_g+I+\diag(\Sigma_g)=\diag(\Sigma_g)+S$
% \bas{
% \Sigma_g+\diag(S)=\Sigma_g+I+\diag(\Sigma_g)=\diag(\Sigma_g)+S
% }
It is not hard to see that \textsc{Joint} equals \textsc{Product}. The same argument generalizes to $k>2$. We next provide proof for the harder direction here. Suppose that $\Delta_F(\bt|P) = 0$, i.e, 
\ba{
\E\bbb{\exp(i\bt^TF\bx)}\exp\bb{-\frac{\bt^T\diag\bb{F S F^T}\bt}{2}}-\prod_{j=1}^k\E\bbb{\exp(i t_j (F\bx)_j)}\exp\bb{-\frac{\bt^TFS F^T \bt}{2}} = 0 \label{eq:delta_score_set_zero}
}
We then prove that $F$ must be of the form $DB^{-1}$ for $D$ being a permutation of a diagonal matrix. Then, taking the logarithm, we have
\ba{
    \ln\bb{\E\bbb{\exp(i\bt^TF\bx)}} - \sum_{j=1}^{k}\ln\bb{\E\bbb{\exp(i t_j (F\bx)_j)}} &= \frac{1}{2}\bt^T\bb{\diag\bb{F S F^T} - FS F^T}\bt \notag \\ 
    &= \frac{1}{2}\bt^T\bb{\diag\bb{F \Sigma F^T} - F \Sigma F^T}\bt \notag
}
Under the ICA model we have, $\bx = B\bz + \bg$. Therefore, from Eq~\ref{eq:delta_score_set_zero}, using the definition of the characteristic function of a Gaussian random variable and noting that $\cov\bb{\bg} = \Sigma$, we have
\bas{
    \ln\bb{\E\bbb{\exp(i\bt^TF\bx)}} &= \ln\bb{\E\bbb{\exp(i\bt^TFB\bz)}} + \ln\bb{\E\bbb{\exp(i\bt^TF\bg)}} \\ 
    &= \ln\bb{\E\bbb{\exp(i\bt^TFB\bz)}} - \frac{1}{2}\bt^T F\Sigma F^T\bt
}
and similarly,
\bas{
    \sum_{j=1}^{k}\ln\bb{\E\bbb{\exp(i t_j (F\bx)_j)}} &= \sum_{j=1}^{k}\ln\bb{\E\bbb{\exp(i t_j (FB\bz)_j)}} - \frac{1}{2}\bt^T\diag\bb{F \Sigma F^T}\bt
}
Therefore, from Eq~\ref{eq:delta_score_set_zero}, 
\bas{
    g_{F}\bb{\bt} := \ln\bb{\E\bbb{\exp(i\bt^TFB\bz)}} - \sum_{j=1}^{k}\ln\bb{\E\bbb{\exp(i t_j (FB\bz)_j)}}
}
must be a quadratic function of $\bt$. 
It is important to note that the second term is an additive function w.r.t $t_1, t_2, \cdots t_k$. For simplicity, consider the case of $k=2$ and assume that the joint and marginal characteristic functions of all signals are twice differentiable. Consider $\frac{\partial^{2}\bb{g_{F}\bb{\bt}}}{\partial t_{1}\partial t_{2}}$, then
\begin{align}
\sum_{k=1}^{2}\frac{\partial^{2}\bb{g_{F}\bb{\bt}}}{\partial t_{1}\partial t_{2}} & \equiv const.\label{eq:2nd-derivative} 
\end{align}
To simplify the notation, let $\psi_{j}(t) := \E\bbb{e^{itz_{j}}}$, $f_{j}\equiv\log\psi_{j}$, $M := FB$. The above is then equivalent to
\begin{align}
\sum_{k=1}^{2}M_{1k}M_{2k}f_{k}''(M_{1k}t_{1}+M_{2k}t_{2}) & \equiv const. \label{eq:f''} 
\end{align}
Note that $M$ is invertible since $F,B$ are invertible. Therefore, let ${\bf s}\equiv M^{T}{\bf t}$, then
\begin{align*}
\sum_{k=1}^{2}M_{1k}M_{2k}f_{k}''(s_{k}) & \equiv const.
\end{align*}
which implies that either $M_{1k}M_{2k}=0$ or $f_{k}''(s_{k})\equiv const$,
for $k=1,2$. For any $k\in\{1,2\}$, since $z_{k}$ is non-Gaussian, its
logarithmic characteristic function, i.e. $f_{k}$ cannot be a quadratic
function, so 
\begin{align*}
M_{1k}M_{2k} &= 0
\end{align*}
which implies that each column of $M$ has a zero. Since $M$ is invertible, thus $M$ is either a diagonal matrix or a permutation of a diagonal matrix. Now to consider $k>2$, we just need to apply the above $k=2$ argument to pairwise entries of $\{t_{1},\cdots,t_{k}\}$ with other entries fixed. Hence proved.
\end{proof}

% \begin{theorem}
% % \label{theorem:global_convergence}
% Let $C$ be a matrix of the form $BDB^T$, where $d_{i} := D_{ii} = h_i''(u_i)$ for some random  $u_i$. Let $S_+=\{i:d_i>0\}$. Consider a contrast function $f:\mathbb{R}\rightarrow \mathbb{R}$. Assume that when applied to every non-Gaussian independent component Assumption~\ref{assump:third} holds. Then $f\bb{\bu^TC^{-1}\bx}$ with the constraint of $\langle\bu,\bu\rangle_{C^{-1}}=1$ has local maxima at $B^{-1}\bu=\be_i,i\in S_+$. %Solutions with $B^{-1}\bu=\be_i,i\in S_-$ are local minima. 
%     All solutions with $\be_i^T B^{-1}\bu\neq 0,\ \ \forall i$ (all nonzero elements) are minima, and all solutions with $\be_i^TB^{-1}\bu\neq 0$ for $1<i<k$ (less than $k$, but at least one nonzero element) are saddle points. These are the only fixed points.
% \end{theorem}

\begin{proof}[Proof of Theorem \ref{theorem:global_convergence}] We start with considering the following constrained optimization problem - 
\bas{
    \sup_{\bu^TC^{-1}\bu=1}f\bb{C^{-1}\bu|P}%f\bb{\langle\bu,\bx\rangle_{C^{-1}}}
}
By an application of lagrange multipliers, for the optima $\bu_{\text{opt}}$, we have
\bas{
    \nu C^{-1}\bu_{opt} = C^{-1}\nabla f\bb{(C^{-1})^T\bu_{\text{opt}}|P}, \ \ \ \bu_{\text{opt}}^TC^{-1}\bu_{\text{opt}} = 1
}
with multiplier $\nu \in \mathbb{R}$. This leads to a fixed-point iteration very similar to the one in \cite{NIPS2015_89f03f7d}, which is the motivation for considering such an optimization framework. We now introduce some notations to simplify the problem.
\\ \\ 
Let $\bu = B\balpha$, $z_i' := \frac{z_i}{d_i}$ and $a'_i := \var(z_i') = \frac{a_i}{d_i^2}$, where $a_i := \var(z_i)$. Then,
%Let us re-evaluate the optimization objective in Eq~\ref{eq:obj}.
%Denote by $g_i(u)=\log E[\exp(u z_i')]$. 
% We note,
% \begin{align*}
%     f(\bu)&=\log E \exp(\bu^T \btx)-\bu^T S\bu/2\\
%     &=\log E \exp(\bu^T B\btz)+\log E \exp(\bu^T \bte)-\bu^T (B\diag(\bvar)B^T+\Sigma)\bu/2\\
%     &=\log E \exp(\bu^T B\btz)+\bu^T\Sigma \bu/2-\bu^T (B\diag(\bvar)B^T+\Sigma)\bu/2\\
%     &=\log E \exp(\bu^T B\btz)-\bu^T B\diag(\bvar)B^T\bu/2
% \end{align*}
% Now, we have,
% \begin{align}
%     f(C^{-1}\bu)&=\log E(\exp( \bu^T (C^{-1})^TB\btz))- \bu^T (C^{-1})^T\ba{B\diag(\bvar)B^T}(C^{-1})^T \bu/2 \notag \\
%   %  &=\log E(\exp( \bu^T (C^{-1})^TB\btz)+\bu^T (C^{-1})^TB\bte))- \bu^T (C^{-1})^TS(C^{-1})^T \bu/2\\
%    % &=\log E(\exp( \bu^T (C^{-1})^TB\btz))+\log E(\bu^T (C^{-1})^TB\bte))- \bu^T (C^{-1})^TS(C^{-1})^T \bu/2\\
%     &=\log E(\exp( \bu^T (B^{-1})^T\btz'))-\bu^T (B^{-1})^T\ba{D_0^{-1}\diag(\bvar)D_0^{-1}}B^{-1} \bu/2 \notag \\
%     &=\log E(\exp( \balpha^T\btz'))-\balpha^T\diag(\bvar')\balpha/2=:h(\balpha), \label{eq:alpha-space}
% \end{align}
% where $\balpha = B^{-1}\bu$.
% \\
% \\
% \noindent
% The power method optimizes the following objective.
\begin{align*}
    \sup_{\bu^TC^{-1}\bu=1}f\bb{C^{-1}\bu|P} = 
\sup_{\balpha^TD^{-1}\balpha=1}f\bb{D^{-1}\balpha|\bz}
\end{align*}
%Let $g_i(x) = \ln{\mathbb{E}[\exp\left(xz_{i}\right)]}$.
We will use $f(\bu|P)$ or $f(\bu|\bx)$ where $\bx\sim P$ interchangeably. By Assumption~\ref{assump:third}-(a), we see that $f\bb{D^{-1}\balpha|\bz}=\sum_{i=1}^{k} f\bb{\frac{\alpha_i}{d_i}|z_i}$. For simplicity of notation, we will use $h_i(\alpha_i/d_i)=f\bb{\alpha_i/d_i|z_i}$, since the functional form can be different for different random variables $z_i$. 
Therefore, we seek to characterize the fixed points of the following optimization problem -
\begin{align}
    &\sup_{\balpha}\sum_{i=1}^{k} h_i(\alpha_i/d_i) \label{original_optimization_problem} \\
    &\;\;\; s.t \;\;\; \sum_{i=1}^{k}\frac{\alpha_{i}^{2}}{d_{i}} = 1, \notag \\
    &\;\;\;\;\;\;\;\;\;\;\; d_{i} \neq 0 \;\; \forall \; i \in [k] \notag
\end{align}
where $\balpha, \bm{d} \in \mathbb{R}^{k}$. We have essentially moved from optimizing in the $\langle .,.\rangle_{C^{-1}}$ space to the $\langle .,.\rangle_{D^{-1}}$ space. We find stationary points of the Lagrangian
\begin{align}
\mathcal{L}(\balpha, \lambda) := \sum_{i=1}^{k}h_i\bb{\frac{\alpha_i}{d_i}} - \lambda\left(\sum_{i=1}^{k}\frac{\alpha_{i}^{2}}{d_{i}} - 1\right) \label{eq:surface}
\end{align}
The components of the gradient of $\mathcal{L}(\balpha, \lambda)$ w.r.t $\balpha$ are given as 
\begin{align*}
    \frac{\partial}{\partial \alpha_j} \mathcal{L}(\balpha, \lambda) =\frac{1}{d_j}h_j'\left(\frac{\alpha_j}{d_j}\right)-\lambda \frac{\alpha_j}{d_j},
\end{align*}
%where $s(\balpha,\bd)=\sgn\ba{\sum_{i=1}^{k}h_i\ba{\frac{\alpha_i}{d_i}}}$.
At the fixed point $(\balpha, \lambda)$ we have, 
\begin{align*}
    \forall j \in [k], h_j'\left(\frac{\alpha_j}{d_j}\right)-\lambda\alpha_j = 0
\end{align*}
By Assumption~\ref{assump:third}-(b), for $\alpha_i=0$, the above is automatically zero. However, for $\{j:\alpha_j\neq 0\}$, we have, 
\begin{align}
    \lambda=%\ba{\frac{E\baa{z'_j\exp(\alpha_j z'_j)}}{\alpha_j}-a'_j} = 
    \frac{h_j'\left(\frac{\alpha_j}{d_j}\right)}{\alpha_j} \label{lagrange_condition1}
\end{align}
So, for all $\{j:\alpha_j\neq 0\}$, we must have the same value and sign of $\frac{1}{\alpha_j}h_j'\left(\frac{\alpha_j}{d_j}\right)$. By assumption in the theorem statement, $h_i'''(x)$ does not change sign on the half-lines $x>0$ and $x<0$. Along with Assumption~\ref{assump:third}-(b), this implies that
\begin{align}
    & \; \forall \;x \in [0,\infty), \; \sgn(h_i(x))\;= \sgn(h_i'(x)) \; = \; \sgn(h_{i}''(x )) \; = \sgn(h_{i}'''(x)) = \kappa_{1} \label{sign_condition1} \\
    & \; \forall \;x \in (-\infty, 0],  \; \sgn(h_i(x))\;= -\sgn(h_{i}'(x)) \; = \; \sgn(h_{i}''(x)) \; = -\sgn(h_{i}'''(x)) = \kappa_{2} \label{sign_condition2}
\end{align}
where $\kappa_{1}, \kappa_{2} \in \{-1,1\}$ are constants. Furthermore, we note that Assumption~\ref{assump:third}-(d) ensures that $h_i(x)$ is a symmetric function. Therefore, $\forall \; x \in \mathbb{R}, \sgn(h_i(x)) = \sgn(h_{i}''(x)) = \kappa, \kappa \in \{-1,1\}$. Then, since $d_i = h_i''(u_i)$, 
\begin{align}
    \sgn\bb{d_i} = \sgn\bb{h_i(x)}, \forall x \in \mathbb{R} \label{eq:d_i_h_i_same_sign}
\end{align}
Now, using \ref{lagrange_condition1},
\begin{align}
    \sgn(\lambda)&=\sgn\left(h_j'\left(\frac{\alpha_j}{d_j}\right)\right)\times \sgn(\alpha_j) \notag \\
                 &= \sgn\left(\frac{\alpha_j}{d_j}\right) \times \sgn\left(h_j\left(\frac{\alpha_j}{d_j}\right)\right)\times \sgn(\alpha_j) ,\text{using Eq}~\ref{sign_condition1}\text{ and}~\ref{sign_condition2} \notag \\
                 &= \sgn\left(d_{j}\right) \times \sgn\left(h_j\left(\frac{\alpha_j}{d_j}\right)\right) \notag \\
                 &= 1 ,\text{using Eq}~\ref{eq:d_i_h_i_same_sign} \label{eq:sgn_lambda}
\end{align}
Keeping this in mind, we now compute the Hessian, $H \in \mathbb{R}^{k \times k}$, of the lagrangian, $\mathcal{L}(\balpha, \lambda)$ at the fixed point, $\bb{\balpha, \lambda}$. Recall that, we have $h'_i(0)=0$ and $h''_i(0)=0.$ Thus, for $\{i:\alpha_i=0\}$, 
\begin{align}\label{eq:vhv}
    H_{ii}=-\lambda/ d_i
\end{align}
This implies that $\sgn\bb{d_iH_{ii}} = \sgn\bb{\lambda} = -1$ for $\{i:\alpha_i=0\}$, using Eq~\ref{eq:sgn_lambda}.
\\ \\
For $\{i:\alpha_i \neq 0\}$, we have
\begin{align*}
     H_{ij} = \frac{\partial^2}{\partial \alpha_i\partial \alpha_j}  \mathcal{L}(\balpha, \lambda)\bigg|_{\balpha} &= \mathbbm{1}(i=j)\left[\frac{h_i''\left(\frac{\alpha_i}{d_i}\right)}{d_i^2}- \frac{\lambda}{d_i}\right]\\
     &=\mathbbm{1}(i=j)\bbb{\frac{h_i''\left(\frac{\alpha_i}{d_i}\right)}{d_i^2}-\frac{h_i'\left(\frac{\alpha_i}{d_i}\right)}{\alpha_id_i}}, \qquad \text{for $\alpha_i\neq 0$, using } \ref{lagrange_condition1}\\
     &=\mathbbm{1}(i=j)\frac{1}{d_i\alpha_i}\left[\frac{\alpha_i}{d_i}h_i''\left(\frac{\alpha_i}{d_i}\right)-h_i'\left(\frac{\alpha_i}{d_i}\right)\right], \qquad \text{for $\alpha_i\neq 0$}
\end{align*}
%$i:\alpha_i=0$, $H_{ii}=-\lambda/ d_i$.
\noindent
We consider the pseudo inner product space $\langle .,.\rangle_{D^{-1}}$ for optimizing $\balpha$. Furthermore, since we are in this pseudo-inner product space, we have
\begin{align*}
    \langle\bv,H\bv\rangle_{D^{-1}} = \bv D^{-1}H\bv
\end{align*}
So we will consider the positive definite-ness of the matrix $\tilde{H}:=D^{-1}H$ to characterize the fixed points. Recall that for a differentiable convex function $f$, we have $\forall \; x, y \in dom(f) \subseteq \mathbb{R}^{n}$
\begin{align*}
    f(y) \geq f(x) + \nabla f(x)^{T}\left(y - x\right)
\end{align*}
Therefore, for $\{i:\alpha_i \neq 0\}$, we have
\begin{align}\label{eq:dihii}
    \sgn\left(d_{i}H_{ii}\right) &= \sgn\left(d_{i}\right) \times \sgn\left(d_{i}\alpha_i\right) \times \sgn\left(\frac{\alpha_i}{d_i}h_i''\left(\frac{\alpha_i}{d_i}\right)-h_i'\left(\frac{\alpha_i}{d_i}\right)\right) \notag \\
    &= \sgn\left(\alpha_i\right) \times \sgn\left(h_i'''\left(\frac{\alpha_i}{d_i}\right)\right) ,\text{using convexity/concavity of } h'(.) \notag \\
    &= \sgn\left(\alpha_i\right) \times \sgn\left(h_i'\left(\frac{\alpha_i}{d_i}\right)\right) ,\text{using } \ref{sign_condition1} \text{ and } \ref{sign_condition2} \notag \\
    &= \sgn\left(\alpha_i\right) \times \sgn\left(\lambda\right) \times \sgn\left(\alpha_i\right) ,\text{using } \ref{lagrange_condition1} \notag \\ 
    &= \sgn\left(\lambda\right) \notag \\
    &= 1 ,\text{using Eq}~\ref{eq:d_i_h_i_same_sign}
\end{align}
%Therefore, we move forward under this Assumption and analyze the case of $\sgn\left(\lambda\right) = 1$.
%Without loss of generality, we consider the case of $\lambda$, and hence $s(\balpha,\bd)$ being positive and divide the analysis into following cases.
Let $S := \{i : \alpha_i \neq 0\}$. We then have the following cases - 
\begin{enumerate}
    \item Assume that only one $\alpha_i$ is nonzero, then $\forall \bv$ orthogonal to this direction, $\langle\bv,H\bv\rangle_{D^{-1}}<0$ by Eq~\ref{eq:vhv}. Thus this gives a local maxima.
    \item Assume more than one $\alpha_i$ are nonzero, but $|S| < k$. Then we have for $i\not\in S$, $\tilde{H}_{jj}<0$ using \ref{eq:vhv}. For $i\in S$, $\tilde{H}_{ii}>0$  from Eq~\ref{eq:dihii}. Hence these are saddle points.
    % \begin{itemize}
    % \item   Since $\lambda>0$, we have $sign(h_i'''(\alpha_i/d_i))=sign(\alpha_i)$. We will further divide these into two cases, $\alpha_i>0$ and $\alpha_i<0$. Let $u=\alpha_i/d_i$.
    % \begin{itemize}
    %     \item For $\alpha_i>0$, $h_i'''(u)>0$, $h_i'(u)$ is convex on $[0,\infty)$, or $(-\infty,0]$ and then by convexity, $uh_i''(u)>h_i'(u)$. So $\forall i\in S$, we have $sign(D_{ii}H_{ii})=sign(D_{ii}D_{ii})=1$.
    %   %  \item For $d_i<0,\alpha_i/d_i<0$. In this case, 
    %   \item For $\alpha_i<0$, $h_i'''(u)<0$ and $h_i'(u)$ is concave, and then by concavity, $uh_i''(u)<h_i'(u)$. Thus $\tilde{H}_{ii}>0$. %Similarly, for $g'''(u)<0$, $g_i'(u)$ is concave on $[-\infty,0]$. And the same conclusion holds, i.e. $g_i''(u)>g_i'(u)/u$. 
    % \end{itemize} 
    % \item Thus, $\forall i\in S$, we have $\tilde{H}_{ii}>0$, and $\forall i\not\in S$, we have $\tilde{H}_{ii}<0$. These are saddle points.
    % \end{itemize}
    \item Assume we have $S=[k]$, i.e. $\forall i$, $\alpha_i\neq 0$. In this case, $\tilde{H}\succ 0$. So, we have a local minima. 
\end{enumerate}
This completes our proof.
%A similar argument holds for $\lambda$ (and $s(\balpha,\bd)$) are negative. In this case all the good solutions are minima, solutions with all nonzero components are minima, and solutions with more than one and less than $k$ components nonzero are minima. %same proof technique can be used to find fixed points of the minimization problem. 

% We next prove that the value of the objective function for minima and maxima where $\forall \; i, \alpha_i \neq 0$, are positive and negative, respectively. This shows that all "bad" minimas are positive and all "bad" maximas are negative. 
% \\
% \\
% Without loss of generality, let's assume that $\lambda > 0$. We divide the analysis into the following cases -
% \begin{enumerate}
%     \item $\alpha_i > 0 $:
%     \begin{enumerate}
%         \item $d_{i} > 0$: In this case using Eq~\ref{lagrange_condition1} we have a point  $x = \frac{\alpha_i}{d_{i}} > 0$  where  $h_{i}'(x)  > 0$. Using Eq~\ref{sign_condition1} we have  $\forall \;x \in [0,\infty)$,  $h_{i}(x) > 0$. This implies that, $h_i(x) $ is a convex function. 
%  %  & \;\;\;\;\;\;\;\;\;\;\;\;\;\; 
%  Therefore, using convexity, we have  %$$\forall \;x \in [0,\infty), \; g_{i}(x) - \frac{a_ix^{2}}{2} > x\left(g_{i}'(x) - a_ix\right).$$ \\
%   % & \;\;\;\;\;\;\;\;\;\;\;\;\;\; 
%   $$\forall \;x \in [0,\infty), \; h_{i}(x)  > xh_{i}'(x) .$$ \\
%   Now, observe that, using Eq~\ref{lagrange_condition1}, 
%   $$h_{i}\left(\frac{\alpha_i}{d_{i}}\right)  > \frac{\alpha_i}{d_{i}}h_{i}'\left(\frac{\alpha_i}{d_{i}}\right) = \lambda\frac{\alpha_{i}^{2}}{d_{i}}$$%$$g_{i}\left(\frac{\alpha_i}{d_{i}}\right) - \frac{a_i}{2}\left(\frac{\alpha_i}{d_{i}}\right)^{2} > \frac{\alpha_i}{d_{i}}\left(g_{i}'\left(\frac{\alpha_i}{d_{i}}\right) - a_i\left(\frac{\alpha_i}{d_{i}}\right)\right) = \lambda\frac{\alpha_{i}^{2}}{d_{i}}$$
%   \item $d_{i} < 0$: In this case, using Eq~\ref{lagrange_condition1}  we have a point  $x = \frac{\alpha_i}{d_{i}} < 0$,  where  $h_{i}'(x)  > 0$. 
%   Using Eq~\ref{sign_condition2} we have  $\forall \;x \in (-\infty, 0], h_{i}^{''}(x) < 0$. This implies,  $h_{i}(x) $  is a concave function. 
%   % & \;\;\;\;\;\;\;\;\;\;\;\;\;\; 
%   Therefore, using concavity, we have $$\forall \;x \in (-\infty, 0], \; g_{i}(x) - \frac{a_ix^{2}}{2} > x\left(g_{i}'(x) - a_ix\right)$$ 
%   % & \;\;\;\;\;\;\;\;\;\;\;\;\;\; 
%   Using Eq~\ref{lagrange_condition1}, we have, 
%   $$h\ba{\frac{\alpha_i}{d_i}}>\frac{\alpha_i}{d_i}$$
%   %$$g_{i}\left(\frac{\alpha_i}{d_{i}}\right) - \frac{a_i}{2}\left(\frac{\alpha_i}{d_{i}}\right)^{2} > \frac{\alpha_i}{d_{i}}\left(g_{i}'\left(\frac{\alpha_i}{d_{i}}\right) - a_i\left(\frac{\alpha_i}{d_{i}}\right)\right) = \lambda\frac{\alpha_{i}^{2}}{d_{i}}$$
%     \end{enumerate}
% \end{enumerate}
% % \begin{align*}
% %    & \textbf{Case 1} : \alpha_i > 0 \\
% %    & \;\;\;\;\;\;\; \textbf{Case 1(a)} : d_{i} > 0 \\
% %    & \;\;\;\;\;\;\;\;\;\;\;\;\;\; \text{In this case using } \ref{lagrange_condition1} \text{ we have a point } x = \frac{\alpha_i}{d_{i}} > 0 \text{ where } g_{i}'(x) - x > 0. \\
% %    & \;\;\;\;\;\;\;\;\;\;\;\;\;\; \text{Using } \ref{sign_condition1} \text{ we have } \; \forall \;x \in [0,\infty), \; g_{i}''(x) - a_i > 0 \Rightarrow g_{i}(x) - \frac{a_ix^{2}}{2} \text{ is a convex function}. \\
% %    & \;\;\;\;\;\;\;\;\;\;\;\;\;\; \text{Therefore, using convexity, we have } \forall \;x \in [0,\infty), \; g_{i}(x) - \frac{a_ix^{2}}{2} > x\left(g_{i}'(x) - a_ix\right) \\
% %    & \;\;\;\;\;\;\;\;\;\;\;\;\;\; \text{Using } \ref{lagrange_condition1}, \; g_{i}\left(\frac{\alpha_i}{d_{i}}\right) - \frac{a_i}{2}\left(\frac{\alpha_i}{d_{i}}\right)^{2} > \frac{\alpha_i}{d_{i}}\left(g_{i}'\left(\frac{\alpha_i}{d_{i}}\right) - a_i\left(\frac{\alpha_i}{d_{i}}\right)\right) = \lambda\frac{\alpha_{i}^{2}}{d_{i}}
% % \end{align*}
% % \begin{align*}
% %    & \;\;\;\;\;\;\; \textbf{Case 1(b)} : d_{i} < 0 \\
% %    & \;\;\;\;\;\;\;\;\;\;\;\;\;\; \text{In this case, using  \ref{lagrange_condition1} \text{ we have a point } x = \frac{\alpha_i}{d_{i}} < 0 \text{ where } g_{i}'(x) - a_ix > 0. \\
% %    & \;\;\;\;\;\;\;\;\;\;\;\;\;\; \text{Using } \ref{sign_condition2} \text{ we have } \; \forall \;x \in (-\infty, 0], \; g_{i}^{''}(x) - a_i < 0 \Rightarrow g_{i}(x) - \frac{a_ix^{2}}{2} \text{ is a concave function}. \\
% %    & \;\;\;\;\;\;\;\;\;\;\;\;\;\; \text{Therefore, using concavity, we have } \forall \;x \in (-\infty, 0], \; g_{i}(x) - \frac{a_ix^{2}}{2} > x\left(g_{i}'(x) - a_ix\right) \\
% %    & \;\;\;\;\;\;\;\;\;\;\;\;\;\; \text{Using } \ref{lagrange_condition1}, \; g_{i}\left(\frac{\alpha_i}{d_{i}}\right) - \frac{a_i}{2}\left(\frac{\alpha_i}{d_{i}}\right)^{2} > \frac{\alpha_i}{d_{i}}\left(g_{i}'\left(\frac{\alpha_i}{d_{i}}\right) - a_i\left(\frac{\alpha_i}{d_{i}}\right)\right) = \lambda\frac{\alpha_{i}^{2}}{d_{i}}
% % \end{align*}
% The case for $\alpha_i < 0$ follows similarly.
% Therefore, in all cases, we find that $g_{i}\left(\frac{\alpha_i}{d_{i}}\right) - \frac{a_i}{2}\left(\frac{\alpha_i}{d_{i}}\right)^{2} > \lambda\frac{\alpha_{i}^{2}}{d_{i}} $. Summing over $i$, we obtain $$\sum_{i=1}^{n}g_i\left(\frac{\alpha_i}{d_{i}}\right) - \frac{a_i}{2}\left(\frac{\alpha_i}{d_{i}}\right)^{2} > \lambda\left(\sum_{i}\frac{\alpha_{i}^{2}}{d_{i}}\right) = \lambda > 0.$$ The analysis for $\lambda < 0$ follows analogously.

% % Assume that the initialization satisfies $\|\bu_k-\bu_0\|\leq 1/\|B^{-1}\|$.

% % Thus, we have:
% % \begin{align*}
% %     \|\bu_{k+1}-\bu_0\|\leq \|C^{-1}\|\|B^{-1}\|^{1/2}\left(4c_1^3+2c_1c_2^2\right)^{1/2}\|\bu_k-\bu_0\|^{3/2}
% % \end{align*}
\end{proof}

% \begin{theorem}
%    % \label{theorem:CHF_quasi_orth}
%    Consider the data generated from the noisy ICA model (Eq~\ref{eq:ICA_noise_model}). Define 
%    \bas{
%     g(\bu;\bx) := \log \E\exp(i\bu^T\bx)+\log \E\exp(-i\bu^T\bx)+\bu^T\cov(\bx)\bu
%     }
%    Then, we have $\nabla^2 g(\bu;\bx) = BD_uB^T$, for diagonal matrix $D_{u}$.
% \end{theorem}

\begin{theorem}
   \label{theorem:CGF_CHF_quasi_orth}
   Consider the data generated from the noisy ICA model (Eq~\ref{eq:ICA_noise_model}). If $f(\bu|P)$ is defined as in Eq~\ref{eq:logchf} or Eq~\ref{eq:cgf}, we have $\nabla^2 f(\bu|P)=BD_uB^T$, for some diagonal matrix $D_{u}$, which can be different for the differerent contrast functions.
\end{theorem}
\begin{remark}
    The above theorem is useful because it shows that the Hessian of the contrast functions based on the CHF (eq~\ref{eq:logchf} or CGF (eq~\ref{eq:cgf}) is of the form $BDB^T$, where $D$ is some diagonal matrix. This matrix can be precomputed at some vector $\bu$ and used as the matrix $C$ in the power iteration update (see eq~\ref{eq:power}) in Algorithm~2 of~\cite{NIPS2015_89f03f7d}. 
\end{remark}

\begin{proof}[Proof of Theorem \ref{theorem:CGF_CHF_quasi_orth}]
First, we show that the CHF-based contrast function (see Eq~\ref{eq:logchf}), the Hessian, is of the correct form.

\textbf{CHF-based contrast function}

\begin{align*}
    f(\bu|P) = \log \E\exp(i\bu^TB\bz)+\log \E\exp(-i\bu^TB\bz)+\bu^TB\diag(\bvar)B^T\bu,
\end{align*}
where $\diag(\bvar)$ is the covariance matrix of $\bz$. For simplicity of notation, define $\phi(\bz;\bu):=\E\bbb{\exp(i\bu^TB\bz)}$.
\begin{align*}
    \nabla f(\bu|P)=iB\bb{\underbrace{\frac{\E\bbb{\exp(i\bu^TB\bz)\bz}}{\phi(\bz;\bu)}}_{\mu(\bu;\bz)}-i\frac{\E\bbb{\exp(-i\bu^T\bz)\bz}}{\phi(\bz;-\bu)}}+2B\diag(\bvar)B^T\bu
\end{align*}
Define
$$H(\bu;\bz):=\frac{\E\bbb{\exp(i\bu^TB\bz)\bz\bz^T}}{\phi(\bz;\bu)}-\mu(\bu;\bz)\mu(\bu;\bz)^T.$$
Taking a second derivative, we have:
\begin{align*}
    \nabla^2 f(\bu|P)=B\bb{-H(\bu;\bz)-H(\bz;-\bu)+2\diag(\bvar)}B^T
\end{align*} 
%Clearly the third term is of the form $BDB^T$ for $D=\diag(\bvar)$. We will show that 
All we have to show at this point is that $H(\bu;\bz)$ is diagonal. We will evaluate the $k,\ell$ entry. Let $B_i$ denote the $i^{th}$ column of $B$. Let $Y_{\setminus k,\ell}:=\bu^TB\sum_{j\neq k,\ell}z_j$. Using independence of the components of $\bz$, we have, for $k\neq \ell$,
\ba{
    \label{eq:Huz}
    \frac{\E\bbb{\exp(i\bu^TB_kz_k+i\bu^TB_{\ell}z_{\ell}+Y_{\setminus k\ell})z_kz_{\ell}}}{\phi(\bz;\bu)}=\frac{\E\bbb{\exp(i\bu^TB_kz_k)z_k}\E\bbb{\exp(i\bu^TB_{\ell}z_{\ell})z_{\ell}}}{\E\bbb{\exp(i\bu^TB_kz_k)}\E\bbb{\exp(i\bu^TB_\ell z_\ell)}}
}
Now we evaluate:
\ba{
    \label{eq:muuz}
    \be_k^T\mu(\bu;\bz):=\frac{\E\bbb{\exp(i\bu^TB\bz)z_k}}{\phi(\bz;\bu)}=\frac{\E\bbb{\exp(i\bu^TB_k z_k)z_k}}{\E\bbb{\exp(i\bu^TB_k z_k)}}
}
Using Eqs~\ref{eq:Huz} and~\ref{eq:muuz}, we see that indeed $H(\bu;\bz)$ is diagonal.


Now, we prove the statement about the CGF-based contrast function. 

\textbf{CGF-based contrast function}

We have, 
\bas{
    \nabla f\bb{\bu|P} &= \frac{\mathbb{E}\bbb{\exp\bb{\bu^{T}B\bz}B\bz}}{\mathbb{E}\bbb{\exp\bb{\bu^{T}B\bz}}} - B\diag\bb{\bvar}B^{T}\bu, \\ 
    \nabla^{2} f\bb{\bu|P} &= \frac{\mathbb{E}\bbb{\exp\bb{\bu^{T}B\bz}B\bz \bz^{T}B^{T}}}{\mathbb{E}\bbb{\exp\bb{\bu^{T}B\bz}}} - \frac{\mathbb{E}\bbb{\exp\bb{\bu^{T}B\bz}B\bz}\mathbb{E}\bbb{\exp\bb{\bu^{T}B\bz}\bz^{T}B^{T}}}{\mathbb{E}\bbb{\exp\bb{\bu^{T}B\bz}}^{2}} - B\diag\bb{\bvar}B^{T} \\ 
    &= B\bbb{\underbrace{\frac{\mathbb{E}\bbb{\exp\bb{\bu^{T}B\bz}\bz\bz^{T}}}{\mathbb{E}\bbb{\exp\bb{\bu^{T}B\bz}}} - \frac{\mathbb{E}\bbb{\exp\bb{\bu^{T}B\bz}\bz}\mathbb{E}\bbb{\exp\bb{\bu^{T}B\bz}\bz^{T}}}{\mathbb{E}\bbb{\exp\bb{\bu^{T}B\bz}}^{2}}}_{\text{Covariance of } \bz \; \sim \; \mathbb{P}\bb{\bz}.\dfrac{\exp\bb{\bu^{T}B\bz}}{\mathbb{E}\bbb{\exp\bb{\bu^{T}B\bz}}}} - \diag\bb{\bvar}}B^{T}
}
The new probability density $\mathbb{P}\bb{\bz}.\frac{\exp\bb{\bu^{T}B\bz}}{\mathbb{E}\bbb{\exp\bb{\bu^{T}B\bz}}}$ is an exponential tilt of the original pdf, and since $\{z_{i}\}_{i=1}^{d}$ are independent, and the new tilted density also factorizes over the $z_{i}$'s, therefore, the covariance under this tilted density is also diagonal.
\end{proof}

\subsection{Uniform convergence}
\label{subsection:ica_uniform_convergence}
In this section, we provide the proof for Theorem \ref{theorem:uniform_convergence}. First, we state some preliminary results about the uniform convergence of smooth function classes which will be useful for our proofs.

\subsubsection{Preliminaries}
Let $D := \sup_{\theta, \tilde{\theta} \in \mathbb{T}}\rho_{X}\bb{\theta, \tilde{\theta}}$ denote the diameter of set $\mathbb{T}$, and let $\mathcal{N}_{X}\bb{\delta; \mathbb{T}}$ denote the $\delta$-covering number of $\mathbb{T}$ in the $\rho_{X}$ metric. Then, we have the following standard result:
\begin{proposition} \label{proposition:one_step_discretization} (Proposition 5.17 from~\cite{wainwright2019high}) Let $\left\{X_{\theta}, \theta \in \mathbb{T}\right\}$ be a zero-mean subgaussian process with respect to the metric $\rho_{X}$. Then for any $\delta \in \bbb{0,D}$ such that $\mathcal{N}_{X}\bb{\delta; \mathbb{T}} \geq 10$, we have
\bas{
    \E\bbb{\sup_{\theta, \tilde{\theta} \in \mathbb{T}}\bb{X_{\theta}-X_{\tilde{\theta}}}} \leq 2\E\bbb{\sup_{\gamma, \gamma' \in \mathbb{T}; \rho_{X}\bb{\gamma,\gamma'} \leq \delta}\bb{X_{\theta}-X_{\tilde{\theta}}}} + 4\sqrt{D^{2}\log\bb{\mathcal{N}_{X}\bb{\delta; \mathbb{T}}}}
}
\end{proposition}

\begin{remark}
    \label{remark:subgaussian_suprema}
    For zero-mean subgaussian processes, Proposition~\ref{proposition:one_step_discretization} implies, $\E\bbb{\sup_{\theta \in \mathbb{T}}X_{\theta}} \leq \E\bbb{\sup_{\theta, \tilde{\theta} \in \mathbb{T}}\bb{X_{\theta}-X_{\tilde{\theta}}}}$
\end{remark}

\begin{proposition} \label{proposition:rademacher_complexity_bound} (Theorem 4.10 from \cite{wainwright2019high}) For any $b$-uniformly bounded class of functions $\mathcal{C}$, any positive integer $n \geq 1$, any scalar $\delta \geq 0$, and a set of i.i.d datapoints $\left\{X^{\bb{i}}\right\}_{i \in [n]}$ we have
\bas{
    \sup_{f \in \mathcal{C}}\left|\frac{1}{n}\sum_{i=1}^{n}f\bb{X^{\bb{i}}} - \E\bbb{f\bb{X}}\right|  \leq 2\frac{1}{n}\E_{X,\epsilon}\bbb{\sup_{f \in \mathcal{C}}\left|\sum_{i=1}^{n}\epsilon_{i}f\bb{X^{\bb{i}}}\right|} + \delta 
}
with probability at least $1 - \exp\bb{-\frac{n\delta^{2}}{2b^{2}}}$. Here $\left\{\epsilon_{i}\right\}_{i \in [n]}$ are i.i.d rademacher random variables.
\end{proposition}

\subsubsection{Proof of theorem~\ref{theorem:uniform_convergence}}
For dataset $\left\{\bx^{\bb{i}}\right\}_{i \in [n]}, \bx^{\bb{i}} \in \mathbb{R}^{k}$, consider the following definitions:
\ba{
    &\phi(\bt,F|P) := \E_{\bx}\bbb{\exp(i \bt^T F\bx)}, \;\;\;\; \phi(\bt,F|\hat{P}) := \frac{1}{n}\sum_{j=1}^{n}\exp(i\bt^TF\bx^{\bb{j}}) \\
    &\psi(\bt,F|P) := \prod_{j=1}^{k}\E_{\bx}\bbb{\exp(i t_j (F\bx)_j)}, \;\;\;\; \psi(\bt,F|\hat{P}) := \frac{1}{n}\prod_{j=1}^k\sum_{r=1}^{n}\exp(i t_j (F\bx^{\bb{r}})_j)
}
Let $\Delta(\bt, F|\hat{P})$ be the empirical version where the expectation is replaced by sample averages.
Now define:
\bas{
\Delta(\bt,F|P)&=\babs{\phi(\bt,F|P)\exp\bb{-\bt^T\diag(FSF^T)\bt}-\psi(\bt,F|P)\exp\bb{-\bt^TFSF^T\bt}}\\
\Delta(\bt,F|\hat{P})&=\babs{\phi(\bt,F|\hat{P})\exp\bb{-\bt^T\diag(F\widehat{S}F^T)\bt}-\psi(\bt,F|\hat{P})\exp\bb{-\bt^TF\widehat{S}F^T\bt}}
}
%We have shown that $s(B^{-1})=0$.
\begin{theorem}
    Let $\mathcal{F}:=\{F\in \mathbb{R}^{k\times k}:\|F\|\leq 1\}$. Assume that $\bx\sim\text{subgaussian}(\sigma)$. 
    We have:
    \bas{
    \sup_{F\in\mathcal{F}}|\E_{\bt\in N(0,I_k)}\Delta(\bt,F|P)-\E_{\bt\in N(0,I_k)}\Delta(\bt,F|\hat{P})|=O_P\bb{\sqrt{\frac{k^2\|S\|\max(k,\sigma^{4}\|S\|)\log^2 (n C_k)}{n}}}
    }
    where $C_{k} := \max(1,
k\log\bb{n}\Tr(S))$.
\end{theorem}
\begin{proof}
Define $\mathcal{E}=\{\bt:\|\bt\|\leq \sqrt{2k\log n}\}$. We will use the fact that
$P(\mathcal{E}^c)\leq 2/n$ (Example 2.12,~\cite{wainwright2019high}). Also note that by construction, $\Delta(\bt,F|P)\leq 1$. Observe that:
    \bas{
    &\sup_{F\in\mathcal{F}}|\E_{\bt\in N(0,I_k)}\Delta(\bt,F|P)-\E_{\bt\in N(0,I_k)}\Delta(\bt,F|\hat{P})|\\
    &\leq \sup_{F\in\mathcal{F}}\E_{\bt\in N(0,I_k)}|\Delta(\bt,F|P)-\Delta(\bt,F|\hat{P})|\\
    &\leq \E_{\bt\in N(0,I_k)}\sup_{F\in\mathcal{F}}|\Delta(\bt,F|P)-\Delta(\bt,F|\hat{P})|\\
    &\leq \E_{\bt\in N(0,I_k)}\bbb{\left.\sup_{F\in\mathcal{F}}\left|\Delta(\bt,F|P)-\Delta(\bt,F|\hat{P})\right|\right|\mathcal{E}}+\E_{\bt\in N(0,I_k)}\bbb{\left.\sup_{F\in\mathcal{F}}|\Delta(\bt,F|P)-\Delta(\bt,F|\hat{P})\right||\mathcal{E}^c}P(\mathcal{E}^c)\\
    &\leq \E_{\bt\in N(0,I_k)}\bbb{\left.\sup_{F\in\mathcal{F}}\left|\Delta(\bt,F|P)-\Delta(\bt,F|\hat{P})\right|\right|\mathcal{E}}+2/n\\
%&=O_P\bb{\nt\sqrt{\frac{k\|S\|\max(k,\sigma^{4}\|S\|)\log (n C_t)}{n}}}+2/n\\
&=O_P\bb{\sqrt{\frac{k^2\|S\|\max(k,\sigma^{4}\|S\|)\log^2 (n C_k)}{n}}}+2/n , \text{ using Theorem~}\ref{thm:indscore}
    }
    The second inequality follows from Jensen's inequality, the fourth follows from $\E[\sup(.)]\leq \sup(\E[.])$.
\end{proof}


\begin{theorem}\label{thm:indscore}
    Let $\mathcal{F}:=\{F\in \mathbb{R}^{k\times k}:\|F\|\leq 1\}$. Assume that $\bx\sim\text{subgaussian}(\sigma)$. 
    We have:
    \bas{
    \sup_{F\in\mathcal{F}}|\Delta(\bt,F|P)-\Delta(\bt,F|\hat{P})|=O_P\bb{\nt\sqrt{\frac{k\|S\|\max(k,\sigma^{4}\|S\|)\log (nC_t)}{n}}}
    }
    where $C_{t} := \max(1,
    \nts\Tr(S))$.
\end{theorem}
%\rd Why is there a $\sigma^2$ and a $\|S\|$?\bk
\begin{proof}
\bas{
\Delta(\bt,F|P)-\Delta(\bt,F|\hat{P})\leq &\babs{\phi(\bt,F|P)-      \phi(\bt,F|\hat{P})}\exp\bb{-\bt^T\diag(FSF^T)\bt}\\
                                    &+\phi(\bt,F|\hat{P})\babs{\exp\bb{-\bt^T\diag(F\widehat{S}F^T)\bt}-\exp\bb{-\bt^T\diag(FSF^T)\bt}}\\
                                    &+\babs{\psi(\bt,F|P)-\psi(\bt,F|\hat{P})}\exp\bb{-\bt^TFSF^T\bt}\\
                                    &+\psi(\bt,F|\hat{P})\babs{\exp\bb{-\bt^TF\widehat{S}F^T\bt}-\exp\bb{-\bt^TFSF^T\bt}}
}
% Now note that, for $\sqrt{k/n}=o(\lambda_{min}(S))$,
% \ba{
% \lambda_{\min}(S')&=\lambda_{\min}(\lambda S+(1-\lambda)\hat{S})\notag\\
% &=\lambda_{\min}(S+(1-\lambda)(S-\hat{S}))\notag\\
% &\geq \lambda_{\min}(S)-O_P\bb{\sqrt{\frac{k}{n}}}\notag\\
% &\geq \frac{1}{2}\lambda_{\min}(S)\label{eq:exp}
% }
Finally, for some $S'=\lambda S+(1-\lambda)\hat{S}$,  %using eq~\ref{eq:exp},
\bas{
|\exp(-\bt^T FSF^T\bt)-\exp(-\bt^T F\hat{S}F^T\bt)|&=\left|\left\langle\left.\partial_{S}\exp(-\bt^T FSF^T\bt)\right|_{S'},\hat{S}-S\right\rangle\right|\\
&\leq \exp(-\bt^T FS'F^T\bt)\left|\bt^TF(S-\hat{S})F^T\bt\right| \\
&\leq \|\bt\|^2\|S-\hat{S}\|= \|\bt\|^2 O_P\bb{ \sigma^{2}\norm{S}\sqrt{\frac{k}{n}}}
}
where the last result follows from Theorem 4.7.1 in \cite{vershynin2018high}. Now, for the second term, observe that:
\ba{
&\left|\phi(\bt,F|\hat{P})\bb{\exp\bb{-\bt^T\diag(F\widehat{S}F^T)\bt}-\exp\bb{-\bt^T\diag(FSF^T)\bt}}\right|\notag\\
&\leq \exp\bb{-\bt^T\diag(FSF^T)\bt}\left|\exp\bb{-\bt^T\diag(F(\widehat{S}-S)F^T)\bt}-1\right|\label{eq:t2}
%&\leq 2|\bt^T\diag(F(\widehat{S}-S)F^T)\bt|
}
%The last inequality uses $|1-e^{x}| \leq 2|x|$, $\forall x \in \bbb{-1,1}$.
We next have that, with probability at least $1-1/n$,
\bas{   
    \left\|\diag\bb{F (S-\widehat{S}) F^T}\right\|_{2} 
    &\leq \left\| F (S-\widehat{S}) F^T \right\|_{2} 
    %&\leq \left\|F\right\|_{2}^{2}\left\|S-\widehat{S}\right\|_{2} \\
    \leq C\bb{ \sigma^{2}\norm{S}\sqrt{\frac{k\log n}{n}}}
}
Thus with probability at least $1-1/n$, using the inequality $|1-e^{x}| \leq 2|x|$, $\forall x \in \bbb{-1,1}$, Eq~\ref{eq:t2} leads to:
\bas{
\left|\phi(\bt,F|\hat{P})\bb{\exp\bb{-\bt^T\diag(F\widehat{S}F^T)\bt}-\exp\bb{-\bt^T\diag(FSF^T)\bt}}\right|
\leq 2C\|\bt\|^2\bk\bb{\sigma^{2}\norm{S}\sqrt{\frac{k\log n}{n}}}
%&\leq 2|\bt^T\diag(F(\widehat{S}-S)F^T)\bt|
}

Note that the matrices $\diag\bb{F (S-\widehat{S}) F^T}$ and $FSF^{T}$ are positive semi-definite and $\bt$ is unit-norm. Observe that, using Lemmas~\ref{lem:scorelhs} and~\ref{lem:scorerhs}, we have:
\bas{
    &\sup_{F\in \mathcal{F}}|\Delta(\bt,F|P)-\Delta(\bt,F|\hat{P})|
    \\ &\leq O_P\bb{\nt\sqrt{\frac{k\Tr(S)\log\bb{nC_{t}}}{n}}}
    +O_P\bb{ \nt\sigma^{2}\norm{S}\sqrt{\frac{k\log n}{n}}}+O_P\bb{\nt\sqrt{\frac{k^2\|S\|\log n}{n}}}\\
&=O_P\bb{\nt\sqrt{\frac{k\|S\|\max(k,\sigma^{4}\|S\|)\log (nC_t)}{n}}}
}
% Now we need the covering number of $\mathcal{F}$.
% Define
\end{proof}
 
\begin{lemma}\label{lem:scorelhs}
    Define $\mathcal{F}:=\{F\in \mathbb{R}^{k\times k}:\|F\|\leq 1\}$. Let $\left\{\bx^{(i)}\right\}_{i \in [n]}$ be i.i.d samples from a subgaussian$\bb{\sigma}$ distribution. We have: %Then, for $\forall \bt \in \mathbb{R}^{d}, \norm{\bt} = 1$, we have
    \bas{
        \sup_{F\in \mathcal{F}}\left|\frac{1}{n}\sum_{j=1}^{n}\exp(i\bt^TF\bx^{\bb{j}})-\E_{\bx}\bbb{\exp(i \bt^T F\bx)}\right| = O_P\bb{\nt\sqrt{\frac{k\Tr(S)\log\bb{nC_{t}}}{n}}}
    }
    where $C_{t} := \max(1,
    \nts\Tr(S))$.
\end{lemma}
\begin{proof}
%     First, we show that $|\phi(\bt,F|\hat{P})-\phi(\bt,F|P)|$ is uniformly small over a ball of unit operator norm.
% \bas{
% \sup_{\|F\|_F\leq 1}|\phi(\bt,F|\hat{P})-\phi(\bt,F|P)|
% }
Let $\phi(\bt,F;\bx^{(i)})=\exp(i \bt^TF\bx^{(i)})$.
Next, we note that $\norm{F}_{2} \leq 1$ imply $\norm{F^{T}\bt} \leq \nt$. Therefore,
\bas{
    \sup_{F\in \mathcal{F}}\left|\frac{1}{n}\sum_{j=1}^{n}\exp(i\bt^TF\bx^{\bb{j}})-\E_{\bx}\bbb{\exp(i \bt^T F\bx)}\right| \leq \sup_{\bu \in \mathbb{R}^{k}, \norm{\bu}\leq1}\left|\frac{1}{n}\sum_{j=1}^{n}\exp(i\bu^{T}\bx^{\bb{j}})-\E_{\bx}\bbb{\exp(i \bu^{T}\bx)}\right|
}
Let $f_X(\bu)=\sum_i\epsilon_i \xi(\bu;\bx^{(i)})$.
We will argue for unit vectors $\bu$, and later scale them by $\nt$.
Define $\mathcal{B}_{k} := \left\{\bu \in \mathbb{R}^{k}, \norm{\bu}\leq1\right\}$, $\xi(\bu;\bx) := \exp(i \bu^T\bx)$ and let 
\bas{
    d(\bu,\bu')^2:=\sum_i (\xi(\bu;\bx^{(i)})-\xi(\bu';\bx^{(i)}))^2
}
Then we have,
\bas{
\|\nabla_{\bu} \exp(i \bu^T\bx)\|_{2} \; = \; \|\bx\||-\sin(\bu^T\bx)+i\cos(\bu^T\bx)| \; \leq \; \|\bx\|
}
Then, 
\bas{
    \left|\xi(\bu;\bx^{\bb{i}})-\xi(\bt;\bx^{\bb{i}})\right| &\leq \min\bb{\|\nabla_{\bu''} \xi(\bu;\bx^{\bb{i}})\|_{2}\left\|\bu-\bu'\right\|, 2} \\
    &\leq \min\bb{\|\bx^{\bb{i}}\|_{2}\left\|\bu-\bu'\right\| ,2}
}
Let $\widehat{S} := \sum_{i}\bx^{(i)}\bb{\bx^{(i)}}^T/n$ and $\tau_n := n\Tr\bb{\widehat{S}}$.
We next have, 
\bas{
D^2 &:= \max_{\bu,\bu'}d(\bu,\bu')^2 \leq \min\left\{4n, \max_{\bu,\bu'}\sum_i\|\bx^{\bb{i}}\|_{2}^{2}\left\|\bu-\bu'\right\|^{2}\right\}\\
&\leq 4\min\left\{n, \tau_n\right\}\\
}

Next, we bound the covering number $N(\delta;\mathcal{B}_{k},d_\bu)$. Note that
$N(\delta;\mathcal{B}_{k},d_\bu)\leq N(\delta/\sqrt{\tau_n};\mathcal{B}_{k},\|.\|_2)$ since,
\bas{
    d\bb{\bu,\bu'}^{2} &=  \sum_i (\xi(\bu;\bx^{(i)})-\xi(\bu';\bx^{(i)}))^2 \\
    &\leq \sum_i\|\bx^{(i)}\|_{2}^{2}\left\|\bu-\bu'\right\|^{2} \\
    &= \sum_{i}\Tr\bb{\bx^{(i)}\bb{\bx^{(i)}}^T}\left\|\bu-\bu'\right\|^{2}
}

Using Cauchy-Schwarz, we have:
\bas{
|f_X(\bu)-f_X(\bu')|\leq \sum_{i=1}^n  \epsilon_i |\xi(\bu;\bx^{(i)})-\xi(\bu;\bx^{(i)})|\leq \sqrt{n d(\bu,\bu')^2}
}

Therefore, we have $N(\delta;\mathcal{B}_{k},d_\bu) \leq \bb{\frac{3\sqrt{\tau_n}}{\delta}}^{k}$. Therefore, using Proposition~\ref{proposition:one_step_discretization}, 
% \bas{
%     \E\bbb{\sup_{F\in\mathcal{F}}|f_X(F)-f_X(F')|}
%     %&\leq 2\E \sup_{\|F\|\leq 1}|f_X(F)-f_X(F')|+\exp(-c_1 n\delta^2)\\
%     &\leq 2\E\bbb{\sup_{d(F,F')\leq \epsilon}|f_X(F)-f_X(F')|}\;+\;4D\E\sqrt{\log N(\epsilon;\mathcal{F},d_F)}\\
%     &\leq \|S\|^{1/2}\bb{2\epsilon\sqrt{n} + 4c_2 k\sqrt{n\log (1/\epsilon)}}\\
%     &= O\bb{\sqrt{\|S\|nk\log n}}
% }
\bas{
    \E_{\bx, \epsilon}\bbb{\sup_{F\in\mathcal{F}}|f_X(\bu)-f_X(\bu')|}
    &\leq 2\E_{\bx}\bbb{\sup_{d(\bu,\bu')\leq \delta}|f_X(\bu)-f_X(\bu')|}\;+\;4\E_{\bx}\bbb{D\sqrt{\log N(\delta;\mathcal{B}_k,d_{\bu})}}\\
    &\leq 2\E_{\bx}\bbb{\inf_{\delta > 0}\left\{\delta\sqrt{n} + 8\sqrt{\min\left\{n, \tau_n\right\}}\sqrt{2k\log\bb{\frac{3\sqrt{\tau_n}}{\delta}}}\right\}} \\
    &= O\bb{\E_{\bx} \sqrt{(k\min(n,\tau_n)\log\bb{\max(\tau_n,n)}} }, \text{ setting } \delta := \sqrt{\min(1,\tau_n/n)} \\
     &= \sqrt{kn}O\bb{\sqrt{\min(1,\Tr(S))\log\bb{n\max(1,\Tr(S))} }}, \text{ Cauchy-Schwarz inequality}
   % &\leq O\bbb{\sqrt{\Tr\bb{S}} + 8\sqrt{n}\sqrt{\min\left\{1, \Tr(S)\right\}}\sqrt{2k\log\bb{3n}}}, \text{ using } \E\bbb{\bx\bx^{T}} = S, \text{ Cauchy-Schwarz inequality}
}
Now we recall that $\|F^T\bt\| \leq \nt$, since all vectors $\bt$ appear as $\bt^T F^T \bx^{(i)}$, this simply leads to scaling $\tau$ and $\|S\|$ by $\nts$.
Dividing both sides by $n$, we have
\bas{
    \frac{1}{n}\E_{\bx, \epsilon}\bbb{\sup_{F\in\mathcal{F}}|f_X(\bt)-f_X(\bt')|} &\leq \sqrt{\frac{k}{n}}O\bb{\sqrt{\min(1,\nts\Tr(S))\log\bb{n\max(1,
    \nts\Tr(S))} }}%\bbb{\sqrt{\frac{\Tr\bb{S}}{n}} + 4\sqrt{\frac{2k\log\bb{3n\nts}}{n}}}
}
% Now, we have w.p. at least $1-\exp(-c_1 n\delta^2)$,
% \bas{
% \sup_{F\in\mathcal{F}}|\phi(\bt,F|\hat{P})-\phi(\bt,F|P)| &\leq 
% \E \sup_{F\in\mathcal{F}}|\phi(\bt,F|\hat{P})-\phi(\bt,F|P)|+\delta\\
% &\leq 2 \E_{\epsilon,\bx}\sup_{F\in \mathcal{F}}f_X(F)+\delta
% }
Thus, using Proposition~\ref{proposition:rademacher_complexity_bound} and using the definition of $C_{t}$, we have,
\bas{
    \sup_{F\in \mathcal{F}}|\phi(\bt,F|\hat{P})-\phi(\bt,F|P)|=\sqrt{\frac{k}{n}}O\bb{\sqrt{\nts\Tr(S)\log\bb{nC_{t}} }} = \nt\sqrt{\frac{k}{n}}O\bb{\sqrt{\Tr(S)\log\bb{nC_{t}}}}
}
\end{proof}
\begin{lemma}\label{lem:scorerhs}
Let $\mathcal{F}=\{F\in \mathbb{R}^{k\times k}:\|F\|\leq 1\}$. We have:
\bas{
    \sup_{F\in \mathcal{F}}\left|\psi(\bt,F|\hat{P})-\psi(\bt,F|P)\right|\leq  O_P\bb{\nt\sqrt{\frac{k^2\|S\|\log n}{n}}}
}
\end{lemma}
\begin{proof}
%Let $F_j$ denote the $j^{th}$ row of $F$. $\mathcal{F}_j:=\{\bu\in \mathbb{R}^k:\|\bu\|\leq 1\}$. 
Let $\mathcal{B}_{k} := \left\{\bu \in \mathbb{R}^{k}, \norm{\bu}\leq1\right\}$
Now define $\psi_j(\bt,F;\bx)=\E[\exp(it_j F_j^T\bx)]$. Also define $f^{(j)}_X(F)=\frac{1}{n}\sum_i\epsilon_i \psi_j(\bt,F;\bx^{\bb{i}})$. 
% The only things that change are the distance, diameter, and covering number. Define 
% \bas{
%     d(F_j,F_j')^2:=\sum_i (\psi_j(\bt,F;\bx^{(i)})-\psi_j(\bt,F';\bx^{(i)}))^2
% }
%Also define $\mathcal{F}:=\{F:\|F\|_F\leq 1\}$.
% First, note that $\psi_j(\bt,F;\bx)$ is Lipschitz over $F$. Let $\be_j$ denote a vector with all zeros except one at the $j^{th}$ position.
% \bas{
%     \|\nabla \psi_j(\bt,F;\bx)\|_F = \|t_j \be_j \bx^T\|_F|-\sin(t_j (F\bx)_j)+i\cos(t_j (F\bx)_j)|\leq \|\be_{j}\bx^{T}\|_{F}|t_j|
% }
% Thus,
% \bas{
% D:=\max_{F,F'}d(F_j,F_j')\leq \sqrt{n}{\max_i\|\bx^{(i)}\|}|t_j|
% }
% Next, observe that, with probability $1-\delta$,
% \bas{
% &n\E\sup_{F\in\mathcal{F}^{(j)}}|f^{(j)}_X(F)-f^{(j)}_X(F')|\\
% %&\leq 2\E \sup_{\|F\|\leq 1}|f_X(F)-f_X(F')|+\exp(-c_1 n\delta^2)\\
% &\leq 2\E\sup_{d(F_j,F'_j)\leq \epsilon}|f^{(j)}_X(F)-f^{(j)}_X(F')|+2\E[D]\sqrt{\log N(\epsilon;\mathcal{F}_j,d_F)}\\
% &\leq |t_j|\E\max_i\|\bx^{(i)}\|\bb{\epsilon\sqrt{n} + 4c_2 \sqrt{nk\log (1/\epsilon)}}\\
% &= O\bb{\sqrt{nk^2\log n}}
% }
% The last step follows because $\bx^{(i)}\sim \text{subgaussian}(\sigma)$.
Thus, using the same argument as in Lemma~\ref{lem:scorelhs} for $\bt \leftarrow \nt\be_{j}$ and $\bx^{\bb{i}} \leftarrow t_{j}\bx^{\bb{i}}$, 
\bas{\sup_{F\in \mathcal{F}}|\psi_j(\bt,F|\hat{P})-\psi_j(\bt,F|P)|=O_P\bb{|t_{j}|\sqrt{\frac{k\|S\|\log n}{n}}}
}
Finally, we see that:
% \rd Syamantak, please check: \bk
\bas{
\sup_{F\in \mathcal{F}}\left|\psi(\bt,F|\hat{P})-\psi(\bt,F|P)\right| &\leq \sum_j  O_P\bb{|t_j|\sqrt{\frac{k\|S\|\log n}{n}}}\\
&= O_P\bb{\sqrt{\frac{k\|S\|\log n}{n}}}\bb{\sum_{j=1}^{k}|t_j|} 
&= \nt O_P\bb{\sqrt{\frac{k^2\|S\|\log n}{n}}}
}
The above is true because, for $|a_i|,|b_i|\leq 1$, $i=1,\dots, k$,
\bas{
\left|\prod_{i=1}^k a_i-\prod_{i=1}^k b_i\right|=\left|\sum_{j=0}^{k-1} \prod_{i\leq j}
b_i (a_{j+1}-b_{j+1})\prod_{i=j+2}^k
a_i\right| \leq \sum_{j=0}^{k-1}|a_{j+1}-b_{j+1}| %-\prod_{i=2}^k a_i+b_1(a_2-b_2)\prod_{i=3}^k a_i+\dots+b_1\dots (a_k-b_k)
}
\end{proof}

\subsection{Local convergence}
\label{subsection:ica_local_convergence}
In this section, we provide the proof of Theorem \ref{theorem:local_convergence}. Recall the ICA model from Eq~\ref{eq:ICA_noise_model}, 
\bas{
    & \bx = B\bz + \bg, \\
    & \diag\bb{\bvar} := \cov\bb{\bz}, \;\; \Sigma := \cov\bb{\bg} = B\diag\bb{\bvar}B^{T} + \Sigma
}
We provide the proof for the CGF-based contrast function. The proof for the CHF-based contrast function follows similarly.
% Let the CGF contrast function be 
% \bas{
%     f\bb{\bu|P} &:= \ln\bb{\mathbb{E}\bbb{\exp\bb{u^{T}x}}} - \frac{u^{T}Su}{2} \\ 
%     &= \ln\bb{\mathbb{E}\bbb{\exp\bb{u^{T}\bb{B\bz+e}}}} - \frac{u^{T}\bb{B\diag\bb{a}B^{T} + \Sigma}u}{2} \\ 
%     &= \ln\bb{\mathbb{E}\bbb{\exp\bb{\bu^{T}B\bz}}} -\frac{u^{T}B\diag\bb{a}B^{T}u}{2} +  \ln\bb{\mathbb{E}\bbb{\exp\bb{u^{T}e}}} - \frac{u^{T}\Sigma u}{2} \\ 
%     &= \ln\bb{\mathbb{E}\bbb{\exp\bb{\bu^{T}B\bz}}} -\frac{u^{T}B\diag\bb{a}B^{T}u}{2} \text{, using the definition of the MGF of a multivariate gaussian}
% } 
% Now, we talk about the first and second derivatives. We have, 
From the Proof of Theorem \ref{theorem:CGF_CHF_quasi_orth} we have,
\bas{
    \nabla f\bb{\bu|P} &= \frac{\mathbb{E}\bbb{\exp\bb{\bu^{T}B\bz}B\bz}}{\mathbb{E}\bbb{\exp\bb{\bu^{T}B\bz}}} - B\diag\bb{\bvar}B^{T}\bu, \\ 
    \nabla^{2} f\bb{\bu|P} &= \frac{\mathbb{E}\bbb{\exp\bb{\bu^{T}B\bz}B\bz \bz^{T}B^{T}}}{\mathbb{E}\bbb{\exp\bb{\bu^{T}B\bz}}} - \frac{\mathbb{E}\bbb{\exp\bb{\bu^{T}B\bz}B\bz}\mathbb{E}\bbb{\exp\bb{\bu^{T}B\bz}z^{T}B^{T}}}{\mathbb{E}\bbb{\exp\bb{\bu^{T}B\bz}}^{2}} - B\diag\bb{\bvar}B^{T} \\ 
    &= B\bbb{\underbrace{\frac{\mathbb{E}\bbb{\exp\bb{\bu^{T}B\bz}\bz\bz^{T}}}{\mathbb{E}\bbb{\exp\bb{\bu^{T}B\bz}}} - \frac{\mathbb{E}\bbb{\exp\bb{\bu^{T}B\bz}\bz}\mathbb{E}\bbb{\exp\bb{\bu^{T}B\bz}\bz^{T}}}{\mathbb{E}\bbb{\exp\bb{\bu^{T}B\bz}}^{2}}}_{\text{Covariance of } \bz \; \sim \; \dfrac{\mathbb{P}\bb{\bz}\exp\bb{\bu^{T}B\bz}}{\mathbb{E}\bbb{\exp\bb{\bu^{T}B\bz}}}} - \diag\bb{\bvar}}B^{T}
}
The new probability density $\dfrac{\mathbb{P}\bb{\bz}\exp\bb{\bu^{T}B\bz}}{\mathbb{E}\bbb{\exp\bb{\bu^{T}B\bz}}}$ is an exponential tilt of the original pdf, and since $\{z_{i}\}_{i=1}^{d}$ are independent, and the new tilted density also factorizes over the $z_{i}$'s, therefore, the covariance under this tilted density is also diagonal. 
\\ \\ 
Let $C := BD_{0}B^{T}$. We denote the $i^{th}$ column of B as $B_{.i}$. Define functions 
\bas{
    r_{i}\bb{\bu} &:= \ln\bb{\mathbb{E}\bbb{\exp\bb{\bu^{T}B_{.i}z_{i}}}} -\Var\bb{z_{i}}\frac{\bu^{T}B_{.i}B_{.i}^{T}\bu}{2} \\ 
    g_{i}\bb{\bu} &:= \nabla r_{i}\bb{\bu} %h_{i}\bb{\bu} = \nabla^{2} r_{i}\bb{\bu}
}
Then $f\bb{\bu|P} = \sum_{i=1}^{d}r_{i}\bb{\bu}$, $\nabla f\bb{\bu|P} = \sum_{i=1}^{d} g_{i}\bb{\bu}$. %and $\nabla^{2} f\bb{\bu|P} = \sum_{i=1}^{d} h_{i}\bb{\bu}$.
\noindent
For fixed point iteration, $\bu_{k+1} = \frac{\nabla f\bb{C^{-1}\bu_{k}|P}}{\left\|\nabla f\bb{C^{-1}\bu_{k}|P}\right\|}$. Consider the function $\nabla f\bb{C^{-1}\bu|P}$. Let $\bm{e}_{i}$ denote the $i^{th}$ axis-aligned unit basis vector of $\mathbb{R}^{n}$. Since $B$ is full-rank, we can denote $\balpha = B^{-1}\bu$. We have, 
\bas{
    \nabla f\bb{C^{-1}\bu|P} &= \sum_{i=1}^{d} g_{i}\bb{C^{-1}\bu} \\ 
                        &= \sum_{i=1}^{d} \frac{\mathbb{E}\bbb{\exp\bb{\bu^{T}\bb{C^{-1}}^{T}B_{.i}z_{i}}B_{.i}z_{i}}}{\mathbb{E}\bbb{\exp\bb{\bu^{T}\bb{C^{-1}}^{T}B_{.i}z_{i}}}} -\Var\bb{z_{i}}B_{.i}B_{.i}^{T}C^{-1}\bu \\ 
                        % &= \sum_{i=1}^{d} \frac{\mathbb{E}\bbb{\exp\bb{u^{T}\bb{B^{T}}^{-1}D_{0}^{-1}B^{-1}B_{.i}z_{i}}B_{.i}z_{i}}}{\mathbb{E}\bbb{\exp\bb{u^{T}\bb{B^{T}}^{-1}D_{0}^{-1}B^{-1}B_{.i}z_{i}}}} -\Var\bb{z_{i}}B_{.i}B_{.i}^{T}\bb{B^{T}}^{-1}D_{0}^{-1}B^{-1}u \\ 
                        &= \sum_{i=1}^{d} \frac{\mathbb{E}\bbb{\exp\bb{\bu^{T}\bb{B^{T}}^{-1}D_{0}^{-1}\be_{i}z_{i}}B_{.i}z_{i}}}{\mathbb{E}\bbb{\exp\bb{\bu^{T}\bb{B^{T}}^{-1}D_{0}^{-1}\be_{i}z_{i}}}} -\Var\bb{z_{i}}B_{.i}\be_{i}^{T}D_{0}^{-1}B^{-1}u \\ 
                        % &= \sum_{i=1}^{d} \frac{\mathbb{E}\bbb{\exp\bb{\balpha^{T}D_{0}^{-1}\be_{i}z_{i}}B_{.i}z_{i}}}{\mathbb{E}\bbb{\exp\bb{\balpha^{T}D_{0}^{-1}\be_{i}z_{i}}}} -\Var\bb{z_{i}}B_{.i}\be_{i}^{T}D_{0}^{-1}\balpha \\ 
                        % &= \sum_{i=1}^{d} \frac{\mathbb{E}\bbb{\exp\bb{\frac{\alpha_i}{\bb{D_{0}}_{ii}}z_{i}}B_{.i}z_{i}}}{\mathbb{E}\bbb{\exp\bb{\frac{\alpha_i}{\bb{D_{0}}_{ii}}z_{i}}}} -\Var\bb{z_{i}}B_{.i}\frac{\alpha_i}{\bb{D_{0}}_{ii}} \\ 
                        &= \sum_{i=1}^{d} \bb{\frac{\mathbb{E}\bbb{\exp\bb{\frac{\alpha_{i}}{\bb{D_{0}}_{ii}}z_{i}}z_{i}}}{\mathbb{E}\bbb{\exp\bb{\frac{\alpha_{i}}{\bb{D_{0}}_{ii}}z_{i}}}} -\Var\bb{z_{i}}\frac{\alpha_{i}}{\bb{D_{0}}_{ii}}}B_{.i}
}
For $t \in \mathbb{R}$, define the function $q_{i}\bb{t} : \mathbb{R} \rightarrow \mathbb{R}$ as - 
\bas{
    q_{i}\bb{t} &:= \frac{\mathbb{E}\bbb{\exp\bb{\frac{t}{\bb{D_{0}}_{ii}}z_{i}}z_{i}}}{\mathbb{E}\bbb{\exp\bb{\frac{t}{\bb{D_{0}}_{ii}}z_{i}}}} -\Var\bb{z_{i}}\frac{t}{\bb{D_{0}}_{ii}}
}
Note that $q_{i}\bb{0} = \mathbb{E}\bbb{z_{i}} = 0$.
For $\bm{t} = \bb{t_1, t_2, \cdots t_d} \in \mathbb{R}^{d}$, let $\bm{q}\bb{\bm{t}} := [q_{1}\bb{t_1}, q_{2}\bb{t_2} \cdots q_{d}\bb{t_d}]^{T} \in \mathbb{R}^{d}$. Then, 
\bas{
    \nabla f\bb{C^{-1}u|P} = B\bm{q}\bb{\balpha}
}
Therefore, if $\bu_{t+1} = B\balpha_{t+1}$, then
\bas{
    & \;\;\;\;\;\;\;\; B\balpha_{t+1} = \frac{B\bm{q}\bb{\balpha_{t}}}{\left\|B\bm{q}\bb{\balpha_{t}}\right\|} \\ 
    & \implies \balpha_{t+1} = \frac{\bm{q}\bb{\balpha_{t}}}{\left\|B\bm{q}\bb{\balpha_{t}}\right\|} \\ 
    & \implies \forall i \in [d], \; \bb{\balpha_{t+1}}_{i} = \frac{q_{i}\bb{\bb{\balpha_{t}}_{i}}}{\left\|B\bm{q}\bb{\balpha_{t}}\right\|}
}
At the fixed point, $\balpha^{*} = \dfrac{\be_{1}}{\|B\be_{1}\|}$ and $\dfrac{B\bq\bb{\balpha^{*}}}{\left\|B\bq\bb{\balpha^{*}}\right\|_{2}} = B\be_{1}$. Therefore, 
\ba{
    & \forall i \in [2,d], \bb{\balpha_{t+1}}_{i} - \alpha^{*}_{i} = \bb{\balpha_{t+1}}_{i} = \frac{q_{i}\bb{\bb{\balpha_{t}}_{i}}}{\left\|B\bm{q}\bb{\balpha_{t}}\right\|} \label{eq:i_not_1} \\ 
    & \text{ for } i = 1, \bb{\balpha_{t+1}}_{1} - \alpha^{*}_{1} = \frac{q_{1}\bb{\bb{\balpha_{t}}_{i}}}{\left\|B\bm{q}\bb{\balpha_{t}}\right\|} - \frac{1}{\|B\be_{1}\|} \label{eq:i_eq_1}
}
\noindent
Note the smoothness assumptions on $q_{i}(.)$ mentioned in Theorem \ref{theorem:local_convergence}, $\forall i \in [d]$, 
\begin{enumerate}
    \item $\sup_{t \in [-\|B^{-1}\|_{2},\|B^{-1}\|_{2}]}|q_{i}\bb{t}| \leq c_{1}$ 
    \item $\sup_{t \in [-\|B^{-1}\|_{2},\|B^{-1}\|_{2}]}|q_{i}'\bb{t}| \leq c_{2}$
    \item $\sup_{t \in [-\|B^{-1}\|_{2},\|B^{-1}\|_{2}]}|q_{i}''\bb{t}| \leq c_{3}$
\end{enumerate}
% Further, $\balpha_{0}$ is initialised such that
% \begin{enumerate}
%     % \item $|\bb{\balpha_{0}}_{1}| > |\bb{\balpha_{0}}_{2}| \geq |\bb{\balpha_{0}}_{3}| \geq |\bb{\balpha_{0}}_{4}| \cdots |\bb{\balpha_{0}}_{d}|, \label{assumption:1}$ 
%     \item $\|\balpha_{0}-\balpha^{*}\| \leq R \leq \frac{c_{2}}{c_{3}}$
%     \item $\epsilon := \frac{\left\|B\right\|_{F}}{\left\|B\be_{1}\right\|}\frac{c_{2}}{\left|q_{1}\bb{\alpha^{*}_{1}}\right|}, \;\;  \max\left\{\epsilon R, \frac{\epsilon^{2}}{\left\|B\right\|_{F}},\frac{\epsilon c_{2}}{\left\|B\be_{1}\right\|}\right\} \leq \frac{1}{10}$
% \end{enumerate}
Since $\forall k, \; \|\bu_{k}\| = 1$, therefore, $\forall k, \; \|\balpha_{t}\| \leq \|B^{-1}\|_{2}$. We seek to prove that the sequence $\{\balpha_{t}\}_{k=1}^{n}$ converges to $\balpha^{*}$.
\subsubsection{Taylor expansions}
Consider the function $g_{i}\bb{\bm{y}} := \dfrac{q_{i}\bb{y_{i}}}{\|B\bm{q}\bb{\by}\|}$ for $\by \in \mathbb{R}^{d}$. We start by computing the gradient, $\nabla_{\bm{y}} g_{i}\bb{\bm{y}}$.
\begin{lemma} \label{lemma:grad_g_i} 
    Let $g_{i}\bb{\bm{y}} := \dfrac{q_{i}\bb{\bm{y}_{i}}}{\|B\bm{q}\bb{\by}\|}$, then we have
    \bas{
        \bbb{\nabla_{\bm{y}} g_{i}\bb{\bm{y}}}_{j} = \begin{cases}
                                & \dfrac{1}{\|B\bm{q}\bb{\by}\|}\dfrac{\partial q_{i}\bb{y_{i}}}{\partial y_{i}} - \dfrac{q_{i}\bb{y_{i}}}{\|B\bm{q}\bb{\by}\|^{2}}\dfrac{\partial \|B\bm{q}\bb{\by}\|}{\partial q_{i}\bb{y_{i}}}\dfrac{\partial q_{i}\bb{y_{i}}}{\partial y_{i}}, \;\; \text{ for } j = i \\ \\
                                & -\dfrac{q'_{j}\bb{y_{j}}}{\|B\bm{q}\bb{\by}\|}\dfrac{q_{i}\bb{y_{i}}\be_{j}^{T}B^{T}B\bm{q}\bb{\by}}{\|B\bm{q}\bb{\by}\|^{2}}, \;\; \text{ for } j \neq i
                            \end{cases}
    }
\end{lemma}
\begin{proof} The derivative w.r.t $y_{i}$ is given as - 
\bas{
    \frac{\partial g_{i}\bb{\by}}{\partial y_{i}} &= \frac{1}{\|B\bm{q}\bb{\by}\|}\frac{\partial q_{i}\bb{y_{i}}}{\partial y_{i}} - \frac{q_{i}\bb{y_{i}}}{\|B\bm{q}\bb{\by}\|^{2}}\frac{\partial \|B\bm{q}\bb{\by}\|}{\partial y_{i}} \\ 
    &= \frac{1}{\|B\bm{q}\bb{\by}\|}\dfrac{\partial q_{i}\bb{y_{i}}}{\partial y_{i}} - \frac{q_{i}\bb{y_{i}}}{\|B\bm{q}\bb{\by}\|^{2}}\frac{\partial \|B\bm{q}\bb{\by}\|}{\partial q_{i}\bb{y_{i}}}\frac{\partial q_{i}\bb{y_{i}}}{\partial y_{i}}
} 
Note that 
\bas{
    \nabla_{\by} \|A\by\|_{2} = \frac{1}{\|A\by\|}A^{T}A\by
}
Therefore, 
\ba{
    \frac{\partial g_{i}\bb{\by}}{\partial y_{i}} &= \frac{1}{\|B\bm{q}\bb{\by}\|}\frac{\partial q_{i}\bb{y_{i}}}{\partial y_{i}} - \frac{q_{i}\bb{y_{i}}}{\|B\bm{q}\bb{\by}\|^{2}}\frac{1}{\|B\bm{q}\bb{\by}\|}\be_{i}^{T}B^{T}B\bm{q}\bb{\by}\frac{\partial q_{i}\bb{y_{i}}}{\partial y_{i}} \notag \\
    &= \frac{q_{i}'\bb{y_{i}}}{\|B\bm{q}\bb{\by}\|}\bbb{1 - \frac{q_{i}\bb{y_{i}}\be_{i}^{T}B^{T}B\bm{q}\bb{\by}}{\|B\bm{q}\bb{\by}\|^{2}}} \label{eq:derivative_of_gi_xi}
}
For $j \neq i$, the derivative w.r.t $y_{j}$ is given as -  
\ba{
   \frac{\partial g_{i}\bb{\by}}{\partial y_{j}} &= -\frac{q_{i}\bb{y_{i}}}{ \|B\bm{q}\bb{\by}\|^{2}}\frac{\partial \|B\bm{q}\bb{\by}\|}{\partial q_{j}\bb{y_{j}}}\frac{\partial q_{j}\bb{y_{j}}}{\partial y_{j}} \notag \\ 
   &= -\frac{q'_{j}\bb{y_{j}}}{\|B\bm{q}\bb{\by}\|}\frac{q_{i}\bb{y_{i}}\be_{j}^{T}B^{T}B\bm{q}\bb{\by}}{\|B\bm{q}\bb{\by}\|^{2}} \label{eq:derivative_of_gi_xj}
}
\end{proof}
Next, we bound $q_{i}\bb{t}$ and $q'_{i}\bb{t}$. 
\begin{lemma}
    \label{lemma:q_i_q_i'_bound} Under the smoothness assumptions on $q_{i}(.)$ mentioned in Theorem~\ref{theorem:local_convergence}, we have for $t \in [-\|B^{-1}\|_{2},\|B^{-1}\|_{2}]$, 
    \begin{enumerate}
        \item $\left|q_{1}\bb{t} - q_{1}\bb{\alpha^{*}_{1}}\right| \leq c_{2}\left|t-\alpha^{*}_{1}\right|, \;\; \left|q_{1}'\bb{t} - q_{1}'\bb{\alpha^{*}_{1}}\right| \leq c_{3}\left|t-\alpha^{*}_{1}\right|$
        \item $\forall i$, $|q_{i}\bb{t}| \leq \frac{c_{3}t^{2}}{2}, \;\; |q_{i}'\bb{t}| \leq c_{3}|t|$
    \end{enumerate}
\end{lemma}
\begin{proof}
First consider $q_{i}\bb{t}$ and $q_{i}'\bb{t}$ for $i \neq 1$. Using Taylor expansion around $t = 0$, we have for some $c \in \bb{0,t}$ and $\forall i \neq 1$,
\bas{
    & q_{i}\bb{t} = q_{i}\bb{0} + tq_{i}'\bb{0} + \frac{t^{2}}{2}q_{i}''\bb{c} \text{ and } \\
    & q_{i}'\bb{t} = q_{i}'\bb{0} + tq_{i}''\bb{0}
}
Now, we know that $q_{i}\bb{0} = q_{i}'\bb{0} = 0$. Then using Assumption \ref{assump:third}, we have, for $t \in [-\|B^{-1}\|_{2},\|B^{-1}\|_{2}]$, 
\ba{
    & |q_{i}\bb{t}| \leq \frac{c_{3}t^{2}}{2} \text{ and } \label{eq:taylor_bound_1} \\
    & |q_{i}'\bb{t}| \leq c_{3}|t| \label{eq:taylor_bound_2}
}
Similarly, using Taylor expansion around $\alpha^{*}_{1} = \frac{1}{\|B\be_{1}\|_{2}}$ for $q_{1}\bb{.}$ we have, for some $c' \in \bb{0,t}$
\bas{
    & q_{1}\bb{t} = q_{1}\bb{\alpha^{*}_{1}} + \bb{t-\alpha^{*}_{1}}q_{1}'\bb{c'} \text{ and } \\
    & q_{1}'\bb{t} = q_{1}'\bb{\alpha^{*}_{1}} + \bb{t-\alpha^{*}_{1}}q_{1}''\bb{c'}
}
Therefore, using Assumption \ref{assumption:2}, we have, for $t \in [-\|B^{-1}\|_{2},\|B^{-1}\|_{2}]$
\ba{
    & \left|q_{1}\bb{t} - q_{1}\bb{\alpha^{*}_{1}}\right| \leq c_{2}\left|t-\alpha^{*}_{1}\right| \text{ and } \label{eq:taylor_bound_3} \\
    & \left|q_{1}'\bb{t} - q_{1}'\bb{\alpha^{*}_{1}}\right| \leq c_{3}\left|t-\alpha^{*}_{1}\right| \label{eq:taylor_bound_4}
}
\end{proof}
\subsubsection{Using convergence radius}
In this section, we use the Taylor expansion results for $q_{i}\bb{.}$ and $q_{i}'\bb{.}$ in the previous section to analyze the following functions for $\by \in \mathbb{R}^{d}$,
\begin{enumerate}
    \item $w\bb{\by} := \|B\by\|$
    \item $v\bb{\by;i} := \dfrac{\be_{i}^{T}B^{T}B\bm{q}\bb{\by}}{\|B\bm{q}\bb{\by}\|^{2}}$
\end{enumerate} 

Under the constraints specified in the theorem statement we have,  
\ba{
    & \|\by-\balpha^{*}\|_{2} \leq R, \; R \leq \max\left\{c_{2}/c_{3}, 1\right\}, \notag \\
    & \epsilon := \frac{\left\|B\right\|_{F}}{\left\|B\be_{1}\right\|}\max\left\{\frac{c_{2}}{\left|q_{1}\bb{\alpha^{*}_{1}}\right|}, \frac{c_{3}}{\left|q_{1}'\bb{\alpha^{*}_{1}}\right|}\right\}, \;\; \epsilon R \leq \frac{1}{10} \label{eq:local_convergence_radius}
}
We start with $w\bb{\by}$.
\begin{lemma}
    \label{lemma:normBy_bounds}
    $\forall \; \by \in \mathbb{R}^{d},$ satisfying \eqref{eq:local_convergence_radius}, let $\delta\bb{\by} := \bm{q}\bb{\by} - q_{1}\bb{\alpha^{*}_{1}}\be_{1}$. Then, we have
    \begin{enumerate}
        \item $\bb{1-\epsilon\|\by-\balpha^{*}\|}\left|q_{1}\bb{\alpha^{*}_{1}}\right|\left\|B\be_{1}\right\| \leq \|B\bm{q}\bb{\by}\| \leq \bb{1+\epsilon\|\by-\balpha^{*}\|}\left|q_{1}\bb{\alpha^{*}_{1}}\right|\left\|B\be_{1}\right\|$
        \item $\|\delta\bb{\by}\| \leq c_{2}\|\by-\balpha^{*}\|$
    \end{enumerate}
\end{lemma}
\begin{proof}
    Consider the vector $\delta\bb{\by} := \bm{q}\bb{\by} - q_{1}\bb{\alpha^{*}_{1}}\be_{1}$. Then, 
\ba{
    |\bb{\delta\bb{\by}}_{\ell}| &\leq \begin{cases}
                                 c_{2}\left|y_{1} - \alpha^{*}_{1}\right|, \text{ for } \ell = 1, \text{ using Eq~\ref{eq:taylor_bound_3}} \\ 
                                \frac{c_{3}}{2}\bb{y}_{\ell}^{2}, \text{ for } \ell \neq 1, \text{ using Eq~\ref{eq:taylor_bound_1}}
                              \end{cases} \label{eq:delta_definition}
}
Note that 
\ba{
 \|\delta\bb{\by}\| \leq c_{2}\|\by-\balpha^{*}\| \label{eq:delta_norm}
}
Now consider $w\bb{\by}$. Using the mean-value theorem on the Euclidean norm we have, 
\ba{
    \|B\bm{q}\bb{\by}\| &= \left|q_{1}\bb{\alpha^{*}_{1}}\right|\left\|B\be_{1}\right\| + \frac{1}{\|B\gamma\bb{\by}\|}\bb{B^{T}B\gamma\bb{\by}}^{T}\delta\bb{\by}, \label{eq:Bwx_taylor_expansion}
}
where $\gamma\bb{\by} = \mu \bm{q}\bb{\by} + \bb{1-\mu}q_{1}\bb{\alpha^{*}_{1}}\be_{1}, \; \mu \in \bb{0,1}$. Then, 
\ba{
    \left\|\frac{1}{\|B\gamma\bb{\by}\|}\bb{B^{T}B\gamma\bb{\by}}^{T}\delta\bb{\by}\right\| &\leq \frac{1}{\|B\gamma\bb{\by}\|}\left\|B^{T}B\gamma\bb{\by}\right\|\left\|\delta\bb{\by}\right\| \notag \\ 
    &\leq \frac{1}{\|B\gamma\bb{\by}\|}\left\|B\right\|\left\|B\gamma\bb{\by}\right\|\left\|\delta\bb{\by}\right\| \notag \\ 
    &= \left\|B\right\|\left\|\delta\bb{\by}\right\| \notag \\ 
    &\leq c_{2}\left\|B\right\|\|\by-\balpha^{*}\|, \text{ using Eq~\ref{eq:delta_norm}} \notag \\
    &\leq c_{2}\left\|B\right\|_{F}\|\by-\balpha^{*}\|, \notag \\
    &\leq \epsilon\left|q_{1}\bb{\alpha^{*}_{1}}\right|\left\|B\be_{1}\right\|\|\by-\balpha^{*}\|, \text{ using Eq~\ref{eq:local_convergence_radius}}
}
Therefore using \ref{eq:Bwx_taylor_expansion},
\ba{
   & \bb{1-\epsilon\|\by-\balpha^{*}\|}\left|q_{1}\bb{\alpha^{*}_{1}}\right|\left\|B\be_{1}\right\| \leq \|B\bm{q}\bb{\by}\| \leq \bb{1+\epsilon\|\by-\balpha^{*}\|}\left|q_{1}\bb{\alpha^{*}_{1}}\right|\left\|B\be_{1}\right\| \label{eq:Bwnorm_approximation}
}
\end{proof}
Finally, we consider the function $v\bb{\by;i} := \dfrac{\be_{i}^{T}B^{T}B\bm{q}\bb{\by}}{\|B\bm{q}\bb{\by}\|^{2}}$.
\begin{lemma}
    \label{lemma:normvy_i_bounds} $\forall \; \by \in \mathbb{R}^{d}$ satisfying \eqref{eq:local_convergence_radius}, we have
    \begin{enumerate}
        \item $\text{For } i = 1, \; 1 - 5\epsilon\|\by-\balpha^{*}\| \leq q_{1}\bb{\alpha^{*}_{1}}v\bb{\by|1} \leq 1 + 5\epsilon\|\by-\balpha^{*}\|$
        \item $\text{For } i \neq 1, \; \left|q_{1}\bb{\alpha^{*}_{1}}v\bb{\by;i}\right| \leq \bb{1+5\epsilon\|\by-\balpha^{*}\|}\dfrac{\left\|B\be_{i}\right\|}{\|B\be_{1}\|}$
    \end{enumerate}
\end{lemma}
\begin{proof}
    Define $\theta := \epsilon\|\by-\balpha^{*}\|$ for convenience of notation. We have, 
    \ba{
        \frac{\be_{i}^{T}B^{T}B\bm{q}\bb{\by}}{\|B\bm{q}\bb{\by}\|^{2}} &= q_{1}\bb{\alpha^{*}_{1}}\frac{\be_{i}^{T}B^{T}B\be_{1}}{\|B\bm{q}\bb{\by}\|^{2}} + \frac{\be_{i}^{T}B^{T}B\delta\bb{\by}}{\|B\bm{q}\bb{\by}\|^{2}}, \text{ using definition of } \delta\bb{\by}
    }
    Therefore, if $i = 1$, then using Lemma~\ref{lemma:normBy_bounds},
    \ba{
        \frac{1}{\bb{1+\theta}^{2}} + \frac{q_{1}\bb{\alpha^{*}_{1}}\be_{1}^{T}B^{T}B\delta\bb{\by}}{\|B\bm{q}\bb{\by}\|^{2}} \leq \frac{q_{1}\bb{\alpha^{*}_{1}}\be_{1}^{T}B^{T}B\bm{q}\bb{\by}}{\|B\bm{q}\bb{\by}\|^{2}} \leq \frac{1}{\bb{1-\theta}^{2}} + \frac{q_{1}\bb{\alpha^{*}_{1}}\be_{1}^{T}B^{T}B\delta\bb{\by}}{\|B\bm{q}\bb{\by}\|^{2}}  \label{eq:R_{1}{x}_bound}
    }
    Using the following inequalities  -
    \bas{
        & \frac{1}{\bb{1+\theta}^{2}} \geq 1-2\theta, \theta \geq 0, \text{ and,  }\frac{1}{\bb{1-\theta}^{2}} \leq 1+5\theta, \theta \leq \frac{1}{5}
    }
    and observing that using Eq~\ref{eq:local_convergence_radius}, and Lemma~\ref{lemma:normBy_bounds}, 
    \bas{
    \left|\frac{q_{1}\bb{\alpha^{*}_{1}}\be_{1}^{T}B^{T}B\delta\bb{\by}}{\|B\bm{q}\bb{\by}\|^{2}}\right| &\leq \left|q_{1}\bb{\alpha^{*}_{1}}\right|\frac{\left\|B\be_{1}\right\|}{\bb{1-\theta}^{2}\left|q_{1}\bb{\alpha^{*}_{1}}\right|^{2}\left\|B\be_{1}\right\|^{2}}\left\|B\right\|\left\|\delta\bb{\by}\right\|  \\
    &\leq \frac{1}{\bb{1-\theta}^{2}}\frac{c_{2}\|B\|}{\left|q_{1}\bb{\alpha^{*}_{1}}\right|\|B\be_{1}\|}\left\|\by-\balpha^{*}\right\| \\
    &\leq \frac{\epsilon}{\bb{1-\theta}^{2}}\left\|\by-\balpha^{*}\right\|
    }
    Therefore we have from Eq~\ref{eq:R_{1}{x}_bound}, 
    \ba{
        1 - 5\theta \leq q_{1}\bb{\alpha^{*}_{1}}v\bb{\by;i} \leq 1 + 5\theta  \label{eq:R_i_bound_i_eq_1}
    }
    where we used $\theta := \epsilon\|\by-\balpha^{*}\|$. For the case of $i \neq 1$, using Lemma~\ref{lemma:normBy_bounds} we have
    \ba{
        \left|q_{1}\bb{\alpha^{*}_{1}}\right|\left|v\bb{\by;i}\right| &= \left|q_{1}\bb{\alpha^{*}_{1}}\right|\left|\frac{\be_{i}^{T}B^{T}B\bm{q}\bb{\by}}{\|B\bm{q}\bb{\by}\|^{2}}\right| \notag \\
        &\leq \left|q_{1}\bb{\alpha^{*}_{1}}\right|\left|\frac{\left\|B\be_{i}\right\|}{\|B\bm{q}\bb{\by}\|}\right| \notag \\
        &\leq \frac{1}{\bb{1-\theta}^{2}}\frac{\left\|B\be_{i}\right\|}{\|B\be_{1}\|} \notag \\
        &\leq \bb{1+5\theta}\frac{\left\|B\be_{i}\right\|}{\|B\be_{1}\|} \label{eq:R_i_bound_i_neq_1}
    }
\end{proof}
% \subsubsection{Using Mean-Value Theorem to establish convergence}
We now operate under the assumption that Eq~\ref{eq:local_convergence_radius} holds for $\by = \balpha_{t}$ and inductively show that it holds for $\by = \balpha_{t+1}$ as well. 
Recall the function $g_{i}\bb{\by} := \dfrac{q_{i}\bb{y_{i}}}{\|B\bm{q}\bb{\by}\|}$. By applying the mean-value theorem for $g_{i}\bb{.}$ for the points $\balpha_{t}$ and $\balpha^{*}$, we have from Eq~\ref{eq:i_not_1} and \ref{eq:i_eq_1} -
\ba{
    \left|\bb{\balpha_{t+1}}_{i} - \alpha^{*}_{i}\right| &= \left|g_{i}\bb{\balpha_{t}} - g_{i}\bb{\balpha^{*}}\right| \notag \\ 
    &= \left|\nabla g_{i}\bb{\bbeta_{i}}^{T}\bb{\balpha_{t} - \balpha^{*}}\right| \text{ for } \bbeta_{i} := \bb{1-\lambda_{i}}\balpha_{t} + \lambda_{i}\balpha^{*}, \; \lambda_{i} \in \bb{0,1} \notag \\
    &\leq \|\nabla g_{i}\bb{\bbeta_{i}}\|\|\balpha_{t}-\balpha^{*}\|
    \label{eq:recursion_i}
}
Note the induction hypothesis assumes that Eq~\ref{eq:local_convergence_radius} is true for $x = \balpha_{t}$. Since $\forall i, \lambda_{i} \in \bb{0,1}$, therefore Eq~\ref{eq:local_convergence_radius} holds for all $\bbeta_{i}$ as well.
Squaring and adding Eq~\ref{eq:recursion_i} for $i \in [d]$ and taking a square-root, we have
\ba{
    \|\balpha_{t+1}-\balpha^{*}\| &\leq \|\balpha_{t}-\balpha^{*}\|\sqrt{\sum_{i=1}^{k}\|\nabla g_{i}\bb{\bbeta_{i}}\|^{2}} \notag \\ 
    &\leq \|\balpha_{t}-\balpha^{*}\|\sqrt{\sum_{i=1}^{k}\sum_{j=1}^{k}\bb{\frac{\partial g_{i}\bb{\by}}{\partial y_{j}}\Bigg|_{\bbeta_{i}}}^{2}} \label{eq:vector_recursion}
}
Let us consider the expression $G_{ij} := \dfrac{\partial g_{i}\bb{\by}}{\partial y_{j}}\Bigg|_{\bbeta_{i}}, \forall i,j \in [k]$. For the purpose of the subsequent analysis, we define $\theta := \epsilon\left\|\balpha_{t}-\balpha^{*}\right\|$. We divide the analysis into the following cases -
\\ \\
\textbf{Case 1} : $i=1, j=1$ 
\\ \\
Using Lemma~\ref{lemma:grad_g_i}, 
\bas{
    \left|G_{11}\right| &= \left|\frac{q_{1}'\bb{\bb{\bbeta_{1}}_{1}}}{\|Bw\bb{\bbeta_{1}}\|}\bbb{1 - \frac{q_{1}\bb{\bb{\bbeta_{1}}_{1}}\be_{1}^{T}B^{T}B\bq\bb{\bbeta_{1}}}{\|B\bq\bb{\bbeta_{1}}\|^{2}}}\right|
}
From Lemma~\ref{lemma:q_i_q_i'_bound}, 
\bas{
    \left|q_{1}'\bb{\bb{\bbeta_{1}}_{1}}\right| &\leq \left|q_{1}'\bb{\alpha^{*}_{1}}\right| + c_{3}\left|\bb{\bbeta_{1}}_{1} - \alpha^{*}_{1}\right| \\ 
    &= \left|q_{1}'\bb{\alpha^{*}_{1}}\right|\bb{1 + c_{3}\frac{\left|\bb{\bbeta_{1}}_{1} - \alpha^{*}_{1}\right|}{\left|q_{1}'\bb{\alpha^{*}_{1}}\right|}} \\
    &\leq \left|q_{1}'\bb{\alpha^{*}_{1}}\right|\bb{1 + c_{3}\frac{\norm{\balpha_{t}-\balpha^{*}}}{\left|q_{1}'\bb{\alpha^{*}_{1}}\right|}} \\
    &\leq \left|q_{1}'\bb{\alpha^{*}_{1}}\right|\bb{1 + \theta}
}
and,
\bas{
    q_{1}\bb{\bb{\bbeta_{1}}_{1}} &\leq |q_{1}\bb{\alpha^{*}_{1}}| + c_{2}\left|\bb{\bbeta_{1}}_{1} - \alpha^{*}_{1}\right| \\ 
    &= |q_{1}\bb{\alpha^{*}_{1}}|\bb{1 + c_{2}\frac{\left|\bb{\bbeta_{1}}_{1} - \alpha^{*}_{1}\right|}{|q_{1}\bb{\alpha^{*}_{1}}}} \\
    &\leq |q_{1}\bb{\alpha^{*}_{1}}|\bb{1 + c_{2}\frac{\norm{\balpha_{t}-\balpha^{*}}}{|q_{1}\bb{\alpha^{*}_{1}}|}} \\
    &\leq |q_{1}\bb{\alpha^{*}_{1}}|\bb{1 + \theta}
}
From Lemma~\ref{lemma:normBy_bounds}, 
\bas{
    \|B \bq\bb{\bbeta_{1}}\| &\geq \bb{1-\theta}\left|q_{1}\bb{\alpha^{*}_{1}}\right|\left\|B\be_{1}\right\|
}
% Therefore, 
% \bas{
%     \frac{q_{1}\bb{\bb{\bbeta_{1}}_{1}}\be_{1}^{T}B^{T}Bw\bb{\bbeta_{1}}}{\|Bw\bb{\bbeta_{1}}\|^{2}} &\leq 
% }
From Lemma~\ref{lemma:normvy_i_bounds} and $\theta \leq \frac{1}{10}$,
\bas{
    -6\theta \leq 1 - \frac{q_{1}\bb{\bb{\bbeta_{1}}_{1}}\be_{1}^{T}B^{T}Bw\bb{\bbeta_{1}}}{\|Bw\bb{\bbeta_{1}}\|^{2}} \leq 6\theta
}
Therefore,
\bas{
    \left|G_{11}\right| &\leq 6\bb{\frac{1 + \theta}{1 - \theta}}\frac{\left|q_{1}'\bb{\alpha^{*}_{1}}\right|\theta}{\|B\be_{1}\|\left|q_{1}\bb{\alpha^{*}_{1}}\right|} \\ 
    &\leq \frac{7.5\epsilon \left|q_{1}'\bb{\alpha^{*}_{1}}\right|}{\|B\be_{1}\|\left|q_{1}\bb{\alpha^{*}_{1}}\right|}\left\|\balpha_{t}-\balpha^{*}\right\|, \text{ since } \theta \leq \frac{1}{10} 
    % &\leq 7.5c_{3}\frac{\left\|B\right\|_{F}\left\|\balpha_{t}-\balpha^{*}\right\|}{\left\|B\be_{1}\right\|^{2}\left|q_{1}\bb{\alpha^{*}_{1}}\right|}, \text{ using Eq}~\ref{eq:local_convergence_radius} 
}
\textbf{Case 2} : $i=1, j \neq 1$ \\ \\ 
Using Lemma~\ref{lemma:grad_g_i},
\bas{
    \left|G_{1j}\right| &= \left|\frac{q'_{j}\bb{\bb{\bbeta_{1}}_{j}}}{\|B\bq\bb{\bbeta_{1}}\|}\frac{q_{1}\bb{\bb{\bbeta_{1}}_{1}}\be_{j}^{T}B^{T}B\bq\bb{\bbeta_{1}}}{\|B\bq\bb{\bbeta_{1}}\|^{2}}\right|
}
From Lemma~\ref{lemma:q_i_q_i'_bound},
\bas{
    \left|q'_{j}\bb{\bb{\bbeta_{1}}_{j}}\right| \leq c_{3}\left|\bb{\bbeta_{1}}_{j} - \alpha^{*}_{j}\right| \leq c_{3}\left|\bb{\balpha_{t}}_{j} - \alpha^{*}_{j}\right|, \text{ since } \alpha^{*}_{j} = 0
}
From Lemma~\ref{lemma:normBy_bounds}, 
\bas{
    \|B\bq\bb{\bbeta_{1}}\| \geq \bb{1-\theta}\left|q_{1}\bb{\alpha^{*}_{1}}\right|\left\|B\be_{1}\right\|
}
From Lemma~\ref{lemma:q_i_q_i'_bound}, 
\bas{
    \left|q_{1}\bb{\bb{\bbeta_{1}}_{1}}\right| &\leq \left|q_{1}\bb{\alpha^{*}_{1}}\right| + c_{2}\left|\bb{\bbeta_{1}}_{1} - \alpha^{*}_{1}\right| \\ 
    &= \left|q_{1}\bb{\alpha^{*}_{1}}\right|\bb{1 + c_{2}\frac{\left|\bb{\bbeta_{1}}_{1} - \alpha^{*}_{1}\right|}{\left|q_{1}\bb{\alpha^{*}_{1}}\right|}} \\ 
    &\leq \left|q_{1}\bb{\alpha^{*}_{1}}\right|\bb{1 + \theta}
}
From Lemma~\ref{lemma:normvy_i_bounds}, 
\bas{
    \left|\frac{\be_{j}^{T}B^{T}B\bq\bb{\bbeta_{1}}}{\|B\bq\bb{\bbeta_{1}}\|^{2}}\right| \leq \frac{\bb{1+5\theta}}{|q_{1}\bb{\alpha^{*}_{1}}|}\frac{\left\|B\be_{j}\right\|}{\|B\be_{1}\|}
}
Therefore, 
\bas{
    \left|G_{1j}\right| \leq c_{3}\bb{\frac{1+\theta}{1-\theta}}\bb{1+5\theta}\frac{\left\|B\be_{j}\right\|}{|q_{1}\bb{\alpha^{*}_{1}}|\|B\be_{1}\|^{2}}\left|\bb{\balpha_{t}}_{j} - \alpha^{*}_{j}\right| \leq 2c_{3}\frac{\left\|B\be_{j}\right\|}{|q_{1}\bb{\alpha^{*}_{1}}|\|B\be_{1}\|^{2}}\left|\bb{\balpha_{t}}_{j} - \alpha^{*}_{j}\right|
}
\textbf{Case 3} : $i \neq 1, j=1$ \\ \\
Using Lemma~\ref{lemma:grad_g_i},
\bas{
    \left|G_{i1}\right| &= \left|\frac{q'_{1}\bb{\bb{\bbeta_{i}}_{1}}}{\|B\bq\bb{\bbeta_{i}}\|}\frac{q_{i}\bb{\bb{\bbeta_{i}}_{i}}\be_{1}^{T}B^{T}B\bq\bb{\bbeta_{i}}}{\|B\bq\bb{\bbeta_{i}}\|^{2}}\right|
}
From Lemma~\ref{lemma:q_i_q_i'_bound},
\bas{
    \left|q'_{1}\bb{\bb{\bbeta_{i}}_{1}}\right| &\leq \left|q_{1}'\bb{\alpha^{*}_{1}}\right| + c_{3}\left|\bb{\bbeta_{i}}_{1} - \alpha^{*}_{1}\right| \\ 
    &= \left|q_{1}'\bb{\alpha^{*}_{1}}\right|\bb{1 + c_{3}\frac{\left|\bb{\bbeta_{i}}_{1} - \alpha^{*}_{1}\right|}{\left|q_{1}'\bb{\alpha^{*}_{1}}\right|}} \\ 
    &\leq \left|q_{1}'\bb{\alpha^{*}_{1}}\right|\bb{1 + \theta}
}
From Lemma~\ref{lemma:normBy_bounds}, 
\bas{
    \|B\bq\bb{\bbeta_{i}}\| \geq \bb{1-\theta}\left|q_{1}\bb{\alpha^{*}_{1}}\right|\left\|B\be_{1}\right\|
}
From Lemma~\ref{lemma:q_i_q_i'_bound}, 
\bas{
    \left|q_{i}\bb{\bb{\bbeta_{i}}_{i}}\right| &\leq \frac{c_{3}}{2}\bb{\bb{\bbeta_{i}}_{i} - \alpha^{*}_{i}}^{2} \leq \frac{c_{3}}{2}\bb{\bb{\balpha_{t}}_{i} - \alpha^{*}_{i}}^{2}
}
From Lemma~\ref{lemma:normvy_i_bounds}, 
\bas{
    \left|\frac{\be_{1}^{T}B^{T}B\bq\bb{\bbeta_{i}}}{\|B\bq\bb{\bbeta_{i}}\|^{2}}\right| \leq \frac{1+5\theta}{\left|q_{1}\bb{\alpha^{*}_{1}}\right|}
}
Therefore, 
\bas{
    \left|G_{i1}\right| &\leq \frac{c_{3}}{2}\bb{\frac{\bb{1+\theta}\bb{1+5\theta}}{1-\theta}}\frac{\left|q_{1}'\bb{\alpha^{*}_{1}}\right|}{\left|q_{1}\bb{\alpha^{*}_{1}}\right|^{2}\left\|B\be_{1}\right\|}\bb{\bb{\balpha_{t}}_{i} - \alpha^{*}_{i}}^{2} \\
    &\leq c_{3}\frac{\left|q_{1}'\bb{\alpha^{*}_{1}}\right|}{\left|q_{1}\bb{\alpha^{*}_{1}}\right|^{2}\left\|B\be_{1}\right\|}\bb{\bb{\balpha_{t}}_{i} - \alpha^{*}_{i}}^{2} \\
    &\leq c_{3}R\frac{\left|q_{1}'\bb{\alpha^{*}_{1}}\right|}{\left|q_{1}\bb{\alpha^{*}_{1}}\right|^{2}\left\|B\be_{1}\right\|}\left|\bb{\balpha_{t}}_{i} - \alpha^{*}_{i}\right| \\
    &\leq c_{2}\frac{\left|q_{1}'\bb{\alpha^{*}_{1}}\right|}{\left|q_{1}\bb{\alpha^{*}_{1}}\right|^{2}\left\|B\be_{1}\right\|}\left|\bb{\balpha_{t}}_{i} - \alpha^{*}_{i}\right|
    % &\leq \frac{c_{2}c_{3}R}{\left|q_{1}\bb{\alpha^{*}_{1}}\right|^{2}\left\|B\be_{1}\right\|}\left|\bb{\balpha_{t}}_{i} - \alpha^{*}_{i}\right|
}
\textbf{Case 4} : $i \neq 1, j \neq 1, i = j$ \\ \\
Using Lemma~\ref{lemma:grad_g_i},
\bas{
    \left|G_{ii}\right| &= \left|\frac{q_{i}'\bb{\bb{\bbeta_{i}}_{i}}}{\|B\bq\bb{\bbeta_{i}}\|}\bbb{1 - \frac{q_{i}\bb{\bb{\bbeta_{i}}_{i}}\be_{i}^{T}B^{T}B\bq\bb{\bbeta_{i}}}{\|B\bq\bb{\bbeta_{i}}\|^{2}}}\right|
}
From Lemma~\ref{lemma:q_i_q_i'_bound},
\bas{
    \left|q'_{i}\bb{\bb{\bbeta_{i}}_{i}}\right| &\leq c_{3}\left|\bb{\bbeta_{i}}_{i} - \alpha^{*}_{j}\right| \leq c_{3}\left|\bb{\balpha_{t}}_{i} - \alpha^{*}_{i}\right|
    % \left|q_{i}\bb{\bb{\bbeta_{i}}_{i}}\right| &\leq c_{3}\frac{\left|\bb{\bbeta_{i}}_{i} - \alpha^{*}_{j}\right|^{2}}{2} \leq c_{3}\frac{\left|\bb{\balpha_{t}}_{i} - \alpha^{*}_{i}\right|^{2}}{2} \leq c_{3}\frac{\norm{\balpha_{t}-\balpha^{*}}^{2}}{2} \leq c_{3}R\norm{\balpha_{t}-\balpha^{*}} \leq c_{2}\norm{\balpha_{t}-\balpha^{*}},
}
From Lemma~\ref{lemma:normBy_bounds}, 
\bas{
    \|B\bq\bb{\bbeta_{i}}\| \geq \bb{1-\theta}\left|q_{1}\bb{\alpha^{*}_{1}}\right|\left\|B\be_{1}\right\|
}
From Lemma~\ref{lemma:q_i_q_i'_bound}, 
\bas{
    \left|q_{i}\bb{\bb{\bbeta_{i}}_{i}}\right| &\leq \frac{c_{3}}{2}\bb{\bb{\bbeta_{i}}_{i} - \alpha^{*}_{i}}^{2} \leq \frac{c_{3}}{2}\bb{\bb{\balpha_{t}}_{i} - \alpha^{*}_{i}}^{2} \leq c_{3}R\norm{\balpha_{t}-\balpha^{*}} \leq c_{2}\norm{\balpha_{t}-\balpha^{*}}
}
From Lemma~\ref{lemma:normvy_i_bounds}, 
\bas{
    \left|\frac{\be_{i}^{T}B^{T}B\bq\bb{\bbeta_{i}}}{\|B\bq\bb{\bbeta_{i}}\|^{2}}\right| \leq \bb{1+5\theta}\frac{\left\|B\be_{i}\right\|}{|q_{1}\bb{\alpha^{*}_{1}}|\|B\be_{1}\|}
}
Then, 
\bas{
    \frac{q_{i}\bb{\bb{\bbeta_{i}}_{i}}\be_{i}^{T}B^{T}B\bq\bb{\bbeta_{i}}}{\|B\bq\bb{\bbeta_{i}}\|^{2}} &\leq \bb{1+5\theta}\norm{\balpha_{t}-\balpha^{*}}\frac{\left\|B\be_{i}\right\|}{\|B\be_{1}\|}\frac{c_{2}}{|q_{1}\bb{\alpha^{*}_{1}}|} \\
    &\leq \bb{1+5\theta}\norm{\balpha_{t}-\balpha^{*}}\frac{\left\|B\right\|_{F}}{\|B\be_{1}\|}\frac{c_{2}}{|q_{1}\bb{\alpha^{*}_{1}}|} \\
    &\leq \bb{1+5\theta}\epsilon\norm{\balpha_{t}-\balpha^{*}} \\
    &\leq 2\theta
}
Therefore, 
\bas{
    \left|G_{ii}\right| \leq \frac{c_{3}\left|\bb{\balpha_{t}}_{i} - \alpha^{*}_{i}\right|}{\bb{1-\theta}\left|q_{1}\bb{\alpha^{*}_{1}}\right|\left\|B\be_{1}\right\|} \leq \frac{2c_{3}}{\left|q_{1}\bb{\alpha^{*}_{1}}\right|\left\|B\be_{1}\right\|}\left|\bb{\balpha_{t}}_{i} - \alpha^{*}_{i}\right|
}
\textbf{Case 5} : $i \neq 1, j \neq 1, i \neq j$ \\ \\ 
Using Lemma~\ref{lemma:grad_g_i},
\bas{
    \left|G_{ij}\right| = \left|\frac{q'_{j}\bb{\bb{\bbeta_{i}}_{j}}}{\|B\bq\bb{\bb{\bbeta_{i}}}\|}\frac{q_{i}\bb{\bb{\bbeta_{i}}_{i}}\be_{j}^{T}B^{T}B\bq\bb{\bb{\bbeta_{i}}}}{\|B\bq\bb{\bb{\bbeta_{i}}}\|^{2}}\right|
}
From Lemma~\ref{lemma:q_i_q_i'_bound},
\bas{
    \left|q'_{j}\bb{\bb{\bbeta_{i}}_{j}}\right| &\leq c_{3}\left|\bb{\bbeta_{i}}_{j} - \alpha^{*}_{j}\right| \\ &\leq c_{3}\left|\bb{\balpha_{t}}_{j} - \alpha^{*}_{j}\right|
}
From Lemma~\ref{lemma:normBy_bounds}, 
\bas{
    \|B\bq\bb{\bbeta_{i}}\| \geq \bb{1-\theta}\left|q_{1}\bb{\alpha^{*}_{1}}\right|\left\|B\be_{1}\right\|
}
From Lemma~\ref{lemma:q_i_q_i'_bound}, 
\bas{
    \left|q_{i}\bb{\bb{\bbeta_{i}}_{i}}\right| &\leq \frac{c_{3}}{2}\bb{\bb{\bbeta_{i}}_{i} - \alpha^{*}_{i}}^{2} \leq \frac{c_{3}}{2}\bb{\bb{\balpha_{t}}_{i} - \alpha^{*}_{i}}^{2} \leq \frac{c_{3}R}{2}|\bb{\balpha_{t}}_{i} - \alpha^{*}_{i}|
}
From Lemma~\ref{lemma:normvy_i_bounds}, 
\bas{
    \left|\frac{\be_{j}^{T}B^{T}B\bq\bb{\bbeta_{i}}}{\|B\bq\bb{\bbeta_{i}}\|^{2}}\right| \leq \frac{\bb{1+5\theta}}{|q_{1}\bb{\alpha^{*}_{1}}|}\frac{\left\|B\be_{j}\right\|}{\|B\be_{1}\|}
}
Therefore, 
\bas{
    \left|G_{ij}\right| \leq \frac{c_{3}^{2}R}{2}\bb{\frac{1+5\theta}{1-\theta}}\frac{\left\|B\be_{j}\right\|}{\left|q_{1}\bb{\alpha^{*}_{1}}\right|^{2}\|B\be_{1}\|^{2}}\left|\bb{\balpha_{t}}_{j} - \alpha^{*}_{j}\right| \left|\bb{\balpha_{t}}_{i} - \alpha^{*}_{i}\right| \leq \frac{c_{3}^{2}R\left\|B\be_{j}\right\|}{\left|q_{1}\bb{\alpha^{*}_{1}}\right|^{2}\|B\be_{1}\|^{2}}\left|\bb{\balpha_{t}}_{j} - \alpha^{*}_{j}\right| \left|\bb{\balpha_{t}}_{i} - \alpha^{*}_{i}\right|
}
\subsubsection{Putting everything together}
Putting all the cases together in Eq~\ref{eq:vector_recursion}, we have $\frac{\left\|\balpha_{t+1}-\balpha^{*}\right\|}{\left\|\balpha_{t}-\balpha^{*}\right\|} \leq C\left\|\balpha_{t}-\balpha^{*}\right\|$, where 
\[
C := \sqrt{
    \begin{aligned}
        &\bb{\frac{7.5\epsilon \left|q_{1}'\bb{\alpha^{*}_{1}}\right|}{\|B\be_{1}\|\left|q_{1}\bb{\alpha^{*}_{1}}\right|}}^{2} + 
        \bb{\frac{2c_{3}\left\|B\right\|_{F}}{|q_{1}\bb{\alpha^{*}_{1}}|\|B\be_{1}\|^{2}}}^{2} + \bb{\frac{c_{2}\left|q_{1}'\bb{\alpha^{*}_{1}}\right|}{\left|q_{1}\bb{\alpha^{*}_{1}}\right|^{2}\left\|B\be_{1}\right\|}}^{2} + \bb{\frac{2c_{3}}{\left|q_{1}\bb{\alpha^{*}_{1}}\right|\left\|B\be_{1}\right\|}}^{2} \\ 
        & + \bb{\frac{c_{3}^{2}R^{2}\left\|B\right\|_{F}}{\left|q_{1}\bb{\alpha^{*}_{1}}\right|^{2}\|B\be_{1}\|^{2}}}^{2}
    \end{aligned}
}
\]
Then we have, $C \leq \frac{5\norm{B}_{F}}{\norm{B\mathbf{e}_{1}}^{2}}\max\left\{\frac{c_{3}}{\left|q_{1}\bb{\alpha^{*}_{1}}\right|}, \frac{c_{2}^{2}}{\left|q_{1}\bb{\alpha^{*}_{1}}\right|^{2}}\right\}$.
For linear convergence, we require $CR < 1$, which is ensured by the condition mentioned in Theorem~\ref{theorem:local_convergence}.
% \[
% C := \sqrt{
%     \begin{aligned}
%         &\bb{7.5\frac{\left\|B\right\|c_{2}}{\left\|B\be_{1}\right\|^{2}\left|q_{1}\bb{\alpha^{*}_{1}}\right|^{2}}}^{2} + 
%         \bb{\frac{2c_{3}\left\|B\right\|_{F}}{\|B\be_{1}\|^{2}}}^{2} + \bb{c_{3}R\frac{\left|q_{1}'\bb{\alpha^{*}_{1}}\right|}{\left|q_{1}\bb{\alpha^{*}_{1}}\right|^{2}\left\|B\be_{1}\right\|}}^{2} + \bb{\frac{2c_{3}}{\left|q_{1}\bb{\alpha^{*}_{1}}\right|\left\|B\be_{1}\right\|}}^{2} \\ 
%         & + \bb{\frac{c_{3}^{2}R\left\|B\right\|_{F}}{\left|q_{1}\bb{\alpha^{*}_{1}}\right|\|B\be_{1}\|^{2}}}^{2}
%     \end{aligned}
% }
% \]
% Therefore, simplifying, we have
% \bas{
%     C &\leq 10\max\left\{c_{2}, c_{3}, c_{2}c_{3}\right\}\frac{\|B\|_{F}}{\|B\be_{1}\|^{2}}\max\left\{\frac{1}{\|B\|_{F}} \frac{1}{\left|q_{1}\bb{\alpha^{*}_{1}}\right|}, \frac{1}{\left|q_{1}\bb{\alpha^{*}_{1}}\right|^{2}}\right\} \\
%     &\leq 10\bb{c_{2}+1}\bb{c_{3}+1}\frac{\|B\|_{F}}{\|B\be_{1}\|^{2}}\max\left\{\frac{1}{\|B\|_{F}} \frac{1}{\left|q_{1}\bb{\alpha^{*}_{1}}\right|}, \frac{1}{\left|q_{1}\bb{\alpha^{*}_{1}}\right|^{2}}\right\}
% }
% Recall that $\epsilon := \frac{\left\|B\right\|_{F}}{\left\|B\be_{1}\right\|}\frac{c_{2}}{\left|q_{1}\bb{\alpha^{*}_{1}}\right|}$. Therefore, rewriting we have
% \bas{
%     C \leq 10\epsilon\frac{\bb{c_{2}+1}\bb{c_{3}+1}}{c_{2}\|B\be_{1}\|}\max\left\{1, \frac{\left|q_{1}\bb{\alpha^{*}_{1}}\right|}{\|B\|_{F}}, \frac{1}{\left|q_{1}\bb{\alpha^{*}_{1}}\right|}\right\}
% }
% which is smaller than $\frac{1}{2}$ based on the conditions mentioned in Theorem~\ref{theorem:local_convergence}.
% \bas{
% &:= \sqrt{\bb{7.5\frac{\left\|B\right\|c_{2}\max\left\{c_{2}, Rc_{3}\right\}}{\left\|B\be_{1}\right\|^{2}\left|q_{1}\bb{\alpha^{*}_{1}}\right|^{2}}}^{2} + 
% \bb{\frac{2c_{3}\left\|B\right\|_{F}}{\|B\be_{1}\|^{2}}}^{2} + \bb{c_{3}R\frac{\left|q_{1}'\bb{\alpha^{*}_{1}}\right|}{\left|q_{1}\bb{\alpha^{*}_{1}}\right|^{2}\left\|B\be_{1}\right\|}}^{2} + \bb{\frac{2c_{3}}{\left|q_{1}\bb{\alpha^{*}_{1}}\right|\left\|B\be_{1}\right\|}}^{2}  \\ 
% & \;\;\;\;\;\;\; + \bb{\frac{c_{3}^{2}R\left\|B\right\|_{F}}{\left|q_{1}\bb{\alpha^{*}_{1}}\right|\|B\be_{1}\|^{2}}}^{2}} \\
% &\leq 10\max\left\{c_{2}^{2},c_{3}^{2}R, c_{2}c_{3}R\right\}^{2}\frac{\left\|B\right\|_{F}}{\left\|B\be_{1}\right\|^{2}}\max\left\{\frac{1}{\left|q_{1}\bb{\alpha^{*}_{1}}\right|}, \frac{1}{\left|q_{1}\bb{\alpha^{*}_{1}}\right|^{2}}\right\}
% }
% We have a quadratic convergence if $10\max\left\{c_{2}^{2},c_{3}^{2}R, c_{2}c_{3}R\right\}\frac{\left\|B\right\|_{F}}{\left\|B\be_{1}\right\|^{2}}\max\left\{\frac{1}{\left|q_{1}\bb{\alpha^{*}_{1}}\right|}, \frac{1}{\left|q_{1}\bb{\alpha^{*}_{1}}\right|^{2}}\right\} \leq 1$.

\section{Additional experiments for noisy ICA}
\label{section:ica_additional_experiments}

\subsection{Experiments with super-gaussian source signals}
\label{subsection:supergaussian}
Table~\ref{tab:super} shows the Amari error in the presence of super-Gaussian signals. The signals are 1) a Uniform distribution $U(-\sqrt{3}, \sqrt{3})$, 2) Bernoulli$\left(\frac{1}{2} + \frac{1}{\sqrt{12}}\right)$, 3) Laplace(0, 0.05) (mean 0 and standard deviation 0.05), 4) Exponential(5), and 5) and 6) a Student's $t$-distribution with 3 and 5 degrees of freedom, respectively. The Meta algorithm closely follows the best candidate algorithm even in the presence of many super-Gaussian signals.
\begin{table}[h!]
    \centering
    \caption{\label{tab:super}Variation of Amari error with Sample Size for Heavy-Tailed distributions, averaged over 100 random runs. Noise power $\rho = 0.001$, number of sources $k=6$. }
    \begin{tabular}{|l|ccccc|}
        \toprule
        \diagbox{\textbf{Algorithm}}{$n$} & \textbf{200} & \textbf{500} & \textbf{1000} & \textbf{5000} & \textbf{10000} \\
        \midrule
        Meta &\rd {\bf  0.44376} & \rd{\bf 0.25215} & \rd{\bf 0.20222} & \rd{\bf 0.11435} & \rd{\bf 0.0838}  \\
        CHF & 1.10103 & 0.70823 & 0.52529 & 0.21344 & {\bf 0.09194} \\
        CGF & 1.84266 & 1.51216 & 1.3702  & 0.84351 & 0.58753 \\
        PEGI & 2.27873 & 1.86561 & 1.71474 & 1.33709 & 1.23322 \\
        PFICA & {\bf 0.39237} & {\bf 0.25468} & {\bf 0.21222} & {\bf 0.12878} & 0.12347 \\
        JADE & 0.70174 & 0.39246 & 0.28199 & 0.12652 & 0.09686 \\
        FastICA & 0.66441 & 0.3869  & 0.28419 & 0.13215 & 0.09946 \\
        \bottomrule
    \end{tabular}
\end{table}
\subsection{Image demixing experiments}
\label{section:additional_image_demixing_experiments}

In this section, we provide additional experiments for image-demixing using ICA. We mix images using flattening and linearly mixing them with a $4 \times 4$ matrix $B$ (i.i.d entries $\sim \mathcal{N}(0,1)$) and Wishart noise ($\rho = 0.001$). Demixing is performed using the SINR-optimal demixing matrix (see Section~\ref{sec:background}) and the results are shown in Figure~\ref{figure:image_demixing_experiment_appendix} along with their corresponding independence scores. The CHF-based method recovers the original sources well, upto sign. The Kurtosis and CGF-based method fails to recover the second source. This is consistent with their higher independence score. The Meta algorithm selects CHF from candidates CHF, CGF, Kurtosis, FastICA, and JADE.

% \begin{figure}
%     \centering
%     \begin{subfigure}[b]{\textwidth}
%         \centering
%         \begin{tabular}{cccc}
%             \includegraphics[width=0.2\textwidth]{demixing_exp/Original_Images/source1.png} &
%             \includegraphics[width=0.2\textwidth]{demixing_exp/Original_Images/source2} &
%             \includegraphics[width=0.2\textwidth]{demixing_exp/Original_Images/source3} &
%             \includegraphics[width=0.2\textwidth]{demixing_exp/Original_Images/source4} \\
%         \end{tabular}
%         \caption{Source Images}
%         \vspace{5pt}
%     \end{subfigure}
%     \begin{subfigure}[b]{\textwidth}
%         \centering
%         \begin{tabular}{cccc}
%             \includegraphics[width=0.2\textwidth]{demixing_exp/Mixed_Images/mixed_image_1.png} &
%             \includegraphics[width=0.2\textwidth]{demixing_exp/Mixed_Images/mixed_image_2} &
%             \includegraphics[width=0.2\textwidth]{demixing_exp/Mixed_Images/mixed_image_3} &
%             \includegraphics[width=0.2\textwidth]{demixing_exp/Mixed_Images/mixed_image_4} \\
%         \end{tabular}
%         \caption{Mixed Images}
%     \end{subfigure} 
%     \vspace{0.1pt}
%     \caption{\label{figure:original_mixed_images}Source and Mixed Images}
% \end{figure}

\begin{figure}
    \centering
        \begin{subfigure}[b]{\textwidth}
        \centering
        \begin{tabular}{cccc}
            \includegraphics[width=0.15\textwidth]{demixing_exp/Original_Images/source1.png} &
            \includegraphics[width=0.15\textwidth]{demixing_exp/Original_Images/source2} &
            \includegraphics[width=0.15\textwidth]{demixing_exp/Original_Images/source3} &
            \includegraphics[width=0.15\textwidth]{demixing_exp/Original_Images/source4} \\
        \end{tabular}
        \caption{Source Images}
        \vspace{3pt}
    \end{subfigure}
    \begin{subfigure}[b]{\textwidth}
        \centering
        \begin{tabular}{cccc}
            \includegraphics[width=0.15\textwidth]{demixing_exp/Mixed_Images/mixed_image_1.png} &
            \includegraphics[width=0.15\textwidth]{demixing_exp/Mixed_Images/mixed_image_2} &
            \includegraphics[width=0.15\textwidth]{demixing_exp/Mixed_Images/mixed_image_3} &
            \includegraphics[width=0.15\textwidth]{demixing_exp/Mixed_Images/mixed_image_4} \\
        \end{tabular}
        \caption{Mixed Images}
        \vspace{3pt}
    \end{subfigure}
        \begin{subfigure}[b]{\textwidth}
        \centering
        \begin{tabular}{cccc}
            \includegraphics[width=0.15\textwidth]{demixing_exp/CHF/source1_chf.png} &
            \includegraphics[width=0.15\textwidth]{demixing_exp/CHF/source2_chf} &
            \includegraphics[width=0.15\textwidth]{demixing_exp/CHF/source3_chf} &
            \includegraphics[width=0.15\textwidth]{demixing_exp/CHF/source4_chf} \\
        \end{tabular}
        \caption{Demixed Images using CHF-based contrast function ($\Delta(\bt,F|P) = 5.26 \times 10^{-3}$).}
        \vspace{3pt}
    \end{subfigure}
        \begin{subfigure}[b]{\textwidth}
            \centering
            \begin{tabular}{cccc}
                \includegraphics[width=0.15\textwidth]{demixing_exp/Kurtosis/source1_kurtosis.png} &
                \includegraphics[width=0.15\textwidth]{demixing_exp/Kurtosis/source2_kurtosis} &
                \includegraphics[width=0.15\textwidth]{demixing_exp/Kurtosis/source3_kurtosis} &
                \includegraphics[width=0.15\textwidth]{demixing_exp/Kurtosis/source4_kurtosis} \\
            \end{tabular}
            \caption{Demixed Images using Kurtosis-based contrast function ($\Delta(\bt,F|P) = 2.48 \times 10^{-2}$)}
            \vspace{3pt}
        \end{subfigure}
        \begin{subfigure}[b]{\textwidth}
            \centering
            \begin{tabular}{cccc}
                \includegraphics[width=0.15\textwidth]{demixing_exp/CGF/source1_cgf.png} &
                \includegraphics[width=0.15\textwidth]{demixing_exp/CGF/source2_cgf} &
                \includegraphics[width=0.15\textwidth]{demixing_exp/CGF/source3_cgf} &
                \includegraphics[width=0.15\textwidth]{demixing_exp/CGF/source4_cgf} \\
            \end{tabular}
            \caption{Demixed Images using CGF-based contrast function ($\Delta(\bt,F|P) = 4 \times 10^{-2}$).}
            \vspace{3pt}
        \end{subfigure}
        \begin{subfigure}[b]{\textwidth}
        \centering
        \begin{tabular}{cccc}
            \includegraphics[width=0.15\textwidth]{demixing_exp/FastICA/source1_fastica.png} &
            \includegraphics[width=0.15\textwidth]{demixing_exp/FastICA/source2_fastica.png} &
            \includegraphics[width=0.15\textwidth]{demixing_exp/FastICA/source3_fastica.png} &
            \includegraphics[width=0.15\textwidth]{demixing_exp/FastICA/source4_fastica.png} \\
        \end{tabular}
        \caption{Demixed Images using FastICA ($\Delta(\bt,F|P) = 7.1 \times 10^{-3}$).}
        \vspace{3pt}
    \end{subfigure}
    \begin{subfigure}[b]{\textwidth}
        \centering
        \begin{tabular}{cccc}
            \includegraphics[width=0.15\textwidth]{demixing_exp/JADE/source1_jade.png} &
            \includegraphics[width=0.15\textwidth]{demixing_exp/JADE/source2_jade.png} &
            \includegraphics[width=0.15\textwidth]{demixing_exp/JADE/source3_jade.png} &
            \includegraphics[width=0.15\textwidth]{demixing_exp/JADE/source4_jade.png} \\
        \end{tabular}
        \caption{Demixed Images using JADE ($\Delta(\bt,F|P) = 6.8 \times 10^{-3}$).}
    \end{subfigure}
    \vspace{1pt}    \caption{\label{figure:image_demixing_experiment_appendix} Image-Demixing using ICA}
\end{figure}

\subsection{Image denoising experiments}
\label{section:image_denoising_experiments}

In this experiment, we use the ICA-based denoising technique proposed in \cite{oja1999image} to compare candidate Noisy ICA algorithms and show that the Meta algorithm can pick the best-denoised image based on the independence score proposed in our work.
 
We use the noisy MNIST dataset and further add entrywise Gaussian noise with variance proportional to $\sin^2(c_1\pi(i+j)) (c_1 = 50)$ to the $(i,j)^{\text{th}}$ pixel. Training images are flattened to create a 784-dimensional vector, and PCA is performed to reduce dimensionality to 25, on which subsequently ICA is performed. The original and denoised images along with their independence score are shown in Figure~\ref{fig:image_denoising_ica}. We note that CHF-based denoising provides qualitatively better results, while CGF-based denoising provides the worst results. This is consistent with their corresponding independence scores.

\begin{figure}
    \centering
    \begin{subfigure}[b]{\textwidth}
    \includegraphics[width=\textwidth]{ICA_MNIST_Images/original_images.png}
    \caption{Original Images}
    \vspace{5pt}
    \end{subfigure}
    \begin{subfigure}[b]{\textwidth}
    \includegraphics[width=\textwidth]{ICA_MNIST_Images/chf_denoising.png}
    \caption{CHF-based denoising ($\Delta(\bt,F|P) = 4.4 \times 10^{-6}$) }
    \vspace{5pt}
    \end{subfigure}
    \begin{subfigure}[b]{\textwidth}
    \includegraphics[width=\textwidth]{ICA_MNIST_Images/kurtosis_denoising.png}
    \caption{Kurtosis-based denoising ($\Delta(\bt,F|P) = 1.7 \times 10^{-4}$) }
    \vspace{5pt}
    \end{subfigure}
    \begin{subfigure}[b]{\textwidth}
    \includegraphics[width=\textwidth]{ICA_MNIST_Images/cgf_denoising.png}
    \caption{CGF-based denoising ($\Delta(\bt,F|P) = 9.2 \times 10^{-6}$) }
    \vspace{5pt}
    \end{subfigure}
    \begin{subfigure}[b]{\textwidth}
    \includegraphics[width=\textwidth]{ICA_MNIST_Images/fastica_denoising.png}
    \caption{FastICA-based denoising ($\Delta(\bt,F|P) = 5.9 \times 10^{-6}$) }
    \vspace{5pt}
    \end{subfigure}
    \begin{subfigure}[b]{\textwidth}
    \includegraphics[width=\textwidth]{ICA_MNIST_Images/jade_denoising.png}
    \caption{JADE-based denoising ($\Delta(\bt,F|P) = 4.6 \times 10^{-6}$) }
    \vspace{5pt}
    \end{subfigure}
    \caption{\label{fig:image_denoising_ica} Image Denoising using ICA}
\end{figure}


\section{Algorithm for sequential calculation of independence scores}
\label{section:sequential_algorithm}
In this section, we provide an algorithm for the computation of the independence score proposed in the manuscript, but in a sequential manner. The detailed algorithm is provided as Algorithm \ref{alg:independence_column_wise}. We assume that we have access to a matrix of the form $C=BDB^T$. 
\\ \\
The power method (see Eq~\ref{eq:power}) essentially extracts one column of $B$ (up to scaling) at every step.~\cite{NIPS2015_89f03f7d} provides an elegant way to use the pseudo-Euclidean space to successively project the data onto columns orthogonal to the ones extracted. This enables one to extract each column of the mixing matrix. For completeness, we present the steps of this projection.
 After extracting the $i^{th}$ column $\bu^{(i)}$, the algorithm maintains two matrices. The first, denoted by $U$, estimates the mixing matrix $B$ one column at a time (up to scaling). The second, denoted by $V$ estimates $B^{-1}$ \textit{one row} at a time.
% \ba{\label{eq:projection}
% U(:,i)=\bu^{(i)}\qquad V(i,:)=\frac{C^{\dag}\bu^{(i)}}{{\bu^{(i)}}^TC^{\dag}\bu^{(i)}}
% }
It is possible to extend the independence score in Eq~\ref{eq:score} to the sequential setting as follows. %by using a similar update scheme as in Eq~\ref{eq:projection}.
Let $\bu^{(j)}$ be the $j^{th}$ vector extracted by the power method (Eq~\ref{eq:power}) until convergence. 
After extracting $\ell$ columns, we project the data using $\min(\ell+1,k)$ projection matrices. These would be mutually independent if we indeed extracted different columns of $B$.  Let $C=BDB^T$ for some diagonal matrix $D$. For convenience of notation, denote $B_{i} \equiv B(:,i)$. Set -
\ba{\label{eq:projectionfirst}
\bv^{(j)}=\frac{C^{\dag}\bu^{(j)}}{{\bu^{(j)}}^T C^{\dag}\bu^{(j)}}
}
When $\bu^{(j)}=B_j/\|B_j\|$, 
$\begin{aligned}
\bv^{(j)}=\frac{(B^T)^{-1}D^{\dag}\be_j}{\be_j^TD^{\dag}\be_j}\|B_j\|=(B^T)^{-1}\be_j\|B_j\|.
\end{aligned}$
Let $\bx$ denote an arbitrary datapoint.
Thus 
\ba{\label{eq:proj}
U(:,j)V(j,:)\bx=B_jz_j+U(:,j)V(j,:)\bg
}
Thus the projection on all other columns $j>\ell$ is given by:
\bas{
(I-UV)\bx%=(I-\sum_{j=1}^\ell U(:,j)V(j,:))\bx
=\sum_{i=1}^k B_i z_i-\sum_{j=1}^\ell B_j z_j+ \tilde{\bg}=\sum_{j=\ell+1}^k B_j z_j+\tilde{\bg}
}
where $\tilde{\bg}=F\bg$, where $F$ is some $k\times k$ matrix.
So we have $\min(\ell,k)$ vectors of the form $z_i B_i +\tilde{\bg}_i$, and when $\ell<k$ an additional vector which contains all independent random variables $z_j$, $j>\ell$ along with a mean-zero Gaussian vector. Then we can project each vector to a scalar using unit random vectors and check for independence using the Independence Score defined in~\ref{eq:score}. When $\ell=k$, then all we need are the $j=1,\dots,k$ projections on the $k$ directions identified via ICA in Eq~\ref{eq:proj}.
\begin{figure}
    \centering
    \includegraphics[width=0.5\textwidth]{images/iscore_all.png}
    \caption{Mean independence score with errorbars from 50 random runs}
         \label{fig:iscore_all_exp}
    \label{fig:enter-label1}
    \vspace{20pt}
\end{figure}
\\ \\
As an example, we conduct an experiment (Figure \ref{fig:iscore_all_exp}), where we fix a mixing matrix $B$ using the same generating mechanism in Section~\ref{section:experiments}. Now we create vectors $\mathbf{q}$ which interpolate between $B(:,1)$ and a fixed arbitrary vector orthogonal to $B(:,1)$. As the interpolation changes, we plot the independence score for direction $\mathbf{q}$. The dataset is the 9-dimensional dataset, which has independent components from many different distributions (see Section~\ref{section:experiments}).

This plot clearly shows that there is a clear negative correlation between the score and the dot product of a vector with the column $B(:,1)$. To be concrete, when $\mathbf{q}$ has a small angle with $B(:,1)$, the independence score is small, and the error bars are also very small. However, as the dot product decreases, the score grows and the error bars become larger.
% \begin{algorithm}[H]
% \caption{\label{alg:independence}Independence Score after extracting $\ell$ columns. $U$ and $V$ are running estimates of $B$ and $B^{-1}$ upto $\ell$ columns and rows respectively.} \label{alg:oja}
% \begin{algorithmic}
% \STATE \textbf{Input}
% \bindent 
% \STATE $X\in \mathbb{R}^{n\times k}$, $U\in\mathbb{R}^{k\times \ell}$, $V\in\mathbb{R}^{\ell\times k}$
% \eindent
% \STATE $k_0\leftarrow \min(\ell+1,k)$
% %\STATE Initialize $Y\in\mathbb{R}^{n\times \min(i+1,k)}$
% %\STATE $Y\leftarrow XU$
%  \FOR{$j$ in range[1, $\ell$]}
%  \STATE $Y_j \leftarrow X\bb{U(:,j)V(j,:)}^T$
%  \ENDFOR
% \IF{$\ell<k$}
% \STATE $Y_{\ell+1}\leftarrow X\bb{I-V^TU^T}$
% \ENDIF
% \FOR{$j$ in range[1, $M$]}
% \STATE $\br\leftarrow$ random unit vector in $\mathbb{R}^k$
% \FOR{$a$ in range[1, $k_0$]}
% \STATE $W(:,j)\leftarrow Y_j \br$
% \ENDFOR
% \STATE Let $\balpha\in \mathbb{R}^{1,k_0}$ represent a row of $W$
% \STATE $\tilde{S} \leftarrow \cov(W)$
% \STATE $\gamma\leftarrow \sum_{i,j}\tilde{S}_{ij}$
% \STATE $\bbeta\leftarrow \sum_{i}\tilde{S}_{ii}$
% \STATE $s(j) = |\hat{\E}\exp(\sum_j i\alpha_j)\exp(-\bbeta)-\prod_{j=1}^{k_0}\hat{\E}\exp(i\alpha_j)\exp(-\gamma)|$
% \ENDFOR
% \STATE Return $\text{mean}(s),\text{stdev}(s)$
% \end{algorithmic}
% \end{algorithm}
% We refer to the function sketched in this subsection as $\Delta\bb{X,U,V,\ell}$ which takes as input, the data $X$, the number of columns $\ell$ and matrices $U$, $V$ which are running estimates of $B$ and $B^{-1}$ upto $\ell$ columns and rows respectively. 

\begin{algorithm}[!hbt]
\caption{\label{alg:independence_column_wise}Independence Score after extracting $\ell$ columns. $U$ and $V$ are running estimates of $B$ and $B^{-1}$ upto $\ell$ columns and rows respectively.} \label{alg:oja}
\begin{algorithmic}
\STATE \textbf{Input}
\bindent 
\STATE $X\in \mathbb{R}^{n\times k}$, $U\in\mathbb{R}^{k\times \ell}$, $V\in\mathbb{R}^{\ell\times k}$, Number of random projections $M$
\eindent
\STATE $k_0\leftarrow \min(\ell+1,k)$
%\STATE Initialize $Y\in\mathbb{R}^{n\times \min(i+1,k)}$
%\STATE $Y\leftarrow XU$
 \FOR{$j$ in range[1, $\ell$]}
 \STATE $Y_j \leftarrow X\bb{U(:,j)V(j,:)}^T$
 \ENDFOR
\IF{$\ell<k$}
\STATE $Y_{\ell+1}\leftarrow X\bb{I-V^TU^T}$
\ENDIF
\FOR{$j$ in range[1, $M$]}
\STATE $\bt \leftarrow$ random unit vector in $\mathbb{R}^k$
\FOR{$a$ in range[1, $k_0$]}
\STATE $W(:,j)\leftarrow Y_j \bt$
\ENDFOR
\STATE Let $\balpha\in \mathbb{R}^{1,k_0}$ represent a row of $W$
\STATE $\tilde{S} \leftarrow \cov(W)$
\STATE $\gamma\leftarrow \sum_{i,j}\tilde{S}_{ij}$
\STATE $\bbeta\leftarrow \sum_{i}\tilde{S}_{ii}$
\STATE $s(j) = |\hat{\E}\exp(\sum_j i\alpha_j)\exp(-\bbeta)-\prod_{j=1}^{k_0}\hat{\E}\exp(i\alpha_j)\exp(-\gamma)|$
\ENDFOR
\STATE Return $\text{mean}(s),\text{stdev}(s)$
\end{algorithmic}
\end{algorithm}
We refer to this function as $\Delta\bb{X,U,V,\ell}$ which takes as input, the data $X$, the number of columns $\ell$ and matrices $U$, $V$ which are running estimates of $B$ and $B^{-1}$ upto $\ell$ columns and rows respectively.

% \begin{algorithm}[H]
%     \caption{\label{alg:meta_sequential} Meta-algorithm for choosing best candidate algorithm sequentially .}
%     \begin{algorithmic}
%         \STATE \textbf{Input}
%         \bindent
%         $\text{Algorithm list } \mathcal{L}$, $\text{Data } X \in \mathbb{R}^{n \times k}$
%         \eindent
%         \FOR{$j$ in range[1, $\text{size}\bb{\mathcal{L}}$]}
%             \STATE{$B_{j} \leftarrow \mathcal{L}_{j}\bb{X}$}
%             \COMMENT{\text{Extract mixing matrix} $B_{j}$ \text{using } $j^{th}$ \text{candidate algorithm}}
%             \STATE{$\delta_{j} \leftarrow \widehat{\E}_{\bt \sim \mathcal{N}\bb{0,\bm{I}_{k}}}\bbb{\Delta\bb{\mathbf{t}, B_{j}^{-1}|\hat{P}}}$}
%         \ENDFOR
%         \STATE{$i_{*} \leftarrow \argmin_{j \in [\text{size}\bb{\mathcal{L}}]}\bbb{\delta_{j}}$}
%         \RETURN $B_{i_{*}}$
%     \end{algorithmic}
% \end{algorithm}

\section{More details for third derivative condition in theorem~\ref{theorem:global_convergence}}
\label{section:Assumption_1d}
Theorem~\ref{theorem:global_convergence} requires the condition ``The third derivative of $h_X(u)$ does not change the sign in the half line $[0,\infty)$ for the non-Gaussian random variable considered in the ICA problem.". In this section, we provide sufficient conditions with respect to contrast functions and datasets where this holds. Further, we also provide an interesting example to demonstrate that our Assumption \ref{assump:third}-(d) might not be too far from being necessary.

\textbf{CGF-based contrast function.} Consider the cumulant generating function of a random variable $X$, i.e. the logarithm of the moment generating function. Now consider the contrast function 
\bas{g_X(t)=CGF(t)-\var(X)t^2/2.}  
 We first note that it satisfies Assumption~\ref{assump:third}(a)-(d). Next, we observe that it is enough for a distribution to have all cumulants of the same sign to satisfy Assumption 1(d) for the CGF. For example, the Poisson distribution has all positive cumulants. Figure \ref{fig:cgf_examples} depicts the third derivative of the contrast function for a Bernoulli(1/2), Uniform, Poisson, and Exponential. 
\begin{figure}
    \centering
    \begin{subfigure}[b]{0.48\textwidth}
        \centering
        \includegraphics[width=\textwidth]{third_derivative/cgf_bernoulli.png}
        \caption{\label{fig:cgf_bernoulli}}
    \end{subfigure}
    \begin{subfigure}[b]{0.48\textwidth}
        \centering
        \includegraphics[width=\textwidth]{third_derivative/cgf_uniform.png}
        \caption{\label{fig:cgf_uniform}}
    \end{subfigure}
    \begin{subfigure}[b]{0.48\textwidth}
        \centering
        \includegraphics[width=\textwidth]{third_derivative/cgf_poisson.png}
        \caption{\label{fig:cgf_poisson}}
    \end{subfigure}
    \begin{subfigure}[b]{0.48\textwidth}
        \centering
        \includegraphics[width=\textwidth]{third_derivative/cgf_exponential.png}
        \caption{\label{fig:cgf_exp}}
    \end{subfigure}
    \vspace{30pt}
    \caption{Plots of the third derivative of the CGF-based contrast function for different datasets - Bernoulli($p = 0.5$) \ref{fig:cgf_bernoulli}, Uniform($\mathcal{U}\bb{-\sqrt{3},\sqrt{3}}$) \ref{fig:cgf_uniform}, Poisson($\lambda = 1$) \ref{fig:cgf_poisson} and Exponential($\lambda = 1$) \ref{fig:cgf_exp}. Note that the sign stays the same in each half-line}
    \label{fig:cgf_examples}
\end{figure}

\textbf{Logarithm of the symmetrized characteristic function.} Figure \ref{fig:chf_examples} depicts the third derivative of this contrast function for the logistic distribution where Assumption 1(d) holds. Although the Assumption does not hold for a large class of distributions here, we believe that the notion of having the same sign can be relaxed up to some bounded parameter values instead of the entire half line, so that the global convergence results of Theorem \ref{theorem:global_convergence} still hold. This belief is further reinforced by Figure \ref{fig:surface_plots} which shows that the loss landscape of functions where Assumption 1(d) doesn't hold is still smooth.

\begin{figure}
    \centering
    \includegraphics[width=0.48\textwidth]{third_derivative/chf_logistic.png}
    \caption{Plots of the third derivative of the CHF-based contrast function for Logistic$(0,1)$ dataset.}
    \label{fig:chf_examples}
\end{figure}

\subsection{Towards a necessary condition}
\label{subsection:necessary_condition}
Consider the probability distribution function (see~\cite{doi:10.1080/00031305.1994.10476058}) $f(x) = ke^{-\frac{x^{2}}{2}}(a + b\cos(c\pi x))$ for constants $a,b,c > 0$. For $f(x)$ to be a valid pdf, we require that $f(x) \geq 0$, and $\int_{-\infty}^{\infty} f(x) \,dx = 1$. Therefore, we have,
\begin{align}
    \int_{-\infty}^{\infty} ke^{-\frac{x^{2}}{2}}(a + b\cos(c\pi x)) \,dx = 1
\end{align}

Noting standard integration results, we have that, 
\begin{align}
    \int_{-\infty}^{\infty} e^{-\frac{x^{2}}{2}} \,dx = \sqrt{2\pi} \;\;\;\; \text{ and } \;\;\;\; \int_{-\infty}^{\infty} e^{-\frac{x^{2}}{2}}\cos(c\pi x) \,dx = \sqrt{2\pi}e^{-\frac{c^{2}\pi^{2}}{2}} 
\end{align}
Therefore, $k = \frac{1}{\sqrt{2\pi}(a + be^{-\frac{c^{2}\pi^{2}}{2}})}$.
Some algebraic manipulations yield the following:
\begin{align*}
    MGF_{f(x)}(t) &= %\mathbb{E}_{f(x)}[e^{tx}] \\
    % &= \int_{-\infty}^{\infty} f(x)e^{tx} \,dx \\
    % &= \int_{-\infty}^{\infty} ke^{-\frac{x^{2}}{2}}(a + b\cos(c\pi x))e^{tx} \,dx \\
    % &= \int_{-\infty}^{\infty} ke^{-\frac{1}{2}(x^{2} - 2tx)}(a + b\cos(c\pi x)) \,dx \\
    % &= \int_{-\infty}^{\infty} ke^{-\frac{1}{2}((x-t)^{2} - t^{2})}(a + b\cos(c\pi x)) \,dx \\
    % &= e^{\frac{t^{2}}{2}}\int_{-\infty}^{\infty} ke^{-\frac{1}{2}(x-t)^{2}}(a + b\cos(c\pi x)) \,dx \\
    % & \text{ Let }  u = (x-t), \text{ therefore, } \\
    % &= e^{\frac{t^{2}}{2}}\int_{-\infty}^{\infty} ke^{-\frac{1}{2}u^{2}}(a + b\cos(c\pi u + c\pi t)) \,du \\
    % &= ke^{\frac{t^{2}}{2}}\int_{-\infty}^{\infty} e^{-\frac{1}{2}u^{2}}(a + b\cos(c\pi u + c\pi t)) \,du \\
    % &= ke^{\frac{t^{2}}{2}}\int_{-\infty}^{\infty} e^{-\frac{1}{2}u^{2}}\left(a + b\cos(c\pi u)\cos(c\pi t) - b\sin(c\pi u)\sin(c\pi t)\right) \,du \\
    % &= ke^{\frac{t^{2}}{2}}\left[\int_{-\infty}^{\infty} e^{-\frac{1}{2}u^{2}}\left(a + b\cos(c\pi u)\cos(c\pi t)\right) \,du - \int_{-\infty}^{\infty} e^{-\frac{1}{2}u^{2}}b\sin(c\pi u)\sin(c\pi t) \,du\right] \\
    % &= ke^{\frac{t^{2}}{2}}\left[\int_{-\infty}^{\infty} e^{-\frac{1}{2}u^{2}}\left(a + b\cos(c\pi u)\cos(c\pi t)\right) \,du\right] \\
    % &= ke^{\frac{t^{2}}{2}}\left(\sqrt{2\pi}(a + be^{-\frac{c^{2}\pi^{2}}{2}}\cos(c\pi t))\right) \\
     \frac{1}{(a + be^{-\frac{c^{2}\pi^{2}}{2}})}e^{\frac{t^{2}}{2}}\left(a + be^{-\frac{c^{2}\pi^{2}}{2}}\cos(c\pi t)\right)
\end{align*}
% \begin{align*}
%     MGF_{f(x)}(t) &:= \mathbb{E}_{f(x)}[e^{tx}] \\
%     &= \int_{-\infty}^{\infty} f(x)e^{tx} \,dx \\
%     &= \int_{-\infty}^{\infty} ke^{-\frac{x^{2}}{2}}(a + b\cos(c\pi x))e^{tx} \,dx \\
%     &= \int_{-\infty}^{\infty} ke^{-\frac{1}{2}(x^{2} - 2tx)}(a + b\cos(c\pi x)) \,dx \\
%     &= \int_{-\infty}^{\infty} ke^{-\frac{1}{2}((x-t)^{2} - t^{2})}(a + b\cos(c\pi x)) \,dx \\
%     &= e^{\frac{t^{2}}{2}}\int_{-\infty}^{\infty} ke^{-\frac{1}{2}(x-t)^{2}}(a + b\cos(c\pi x)) \,dx \\
%     & \text{ Let }  u = (x-t), \text{ therefore, } \\
%     &= e^{\frac{t^{2}}{2}}\int_{-\infty}^{\infty} ke^{-\frac{1}{2}u^{2}}(a + b\cos(c\pi u + c\pi t)) \,du \\
%     &= ke^{\frac{t^{2}}{2}}\int_{-\infty}^{\infty} e^{-\frac{1}{2}u^{2}}(a + b\cos(c\pi u + c\pi t)) \,du \\
%     &= ke^{\frac{t^{2}}{2}}\int_{-\infty}^{\infty} e^{-\frac{1}{2}u^{2}}\left(a + b\cos(c\pi u)\cos(c\pi t) - b\sin(c\pi u)\sin(c\pi t)\right) \,du \\
%     &= ke^{\frac{t^{2}}{2}}\left[\int_{-\infty}^{\infty} e^{-\frac{1}{2}u^{2}}\left(a + b\cos(c\pi u)\cos(c\pi t)\right) \,du - \int_{-\infty}^{\infty} e^{-\frac{1}{2}u^{2}}b\sin(c\pi u)\sin(c\pi t) \,du\right] \\
%     &= ke^{\frac{t^{2}}{2}}\left[\int_{-\infty}^{\infty} e^{-\frac{1}{2}u^{2}}\left(a + b\cos(c\pi u)\cos(c\pi t)\right) \,du\right] \\
%     &= ke^{\frac{t^{2}}{2}}\left(\sqrt{2\pi}(a + be^{-\frac{c^{2}\pi^{2}}{2}}\cos(c\pi t))\right) \\
%     &= \frac{1}{(a + be^{-\frac{c^{2}\pi^{2}}{2}})}e^{\frac{t^{2}}{2}}\left(a + be^{-\frac{c^{2}\pi^{2}}{2}}\cos(c\pi t)\right)
% \end{align*}

Therefore, $\ln(MGF_{f(x)}(t)) = -\ln(a + be^{-\frac{c^{2}\pi^{2}}{2}}) + \frac{t^2}{2} + \ln(a + be^{-\frac{c^{2}\pi^{2}}{2}}\cos(c\pi t))$.

\noindent
We now compute the mean, $\mu$, and variance $\sigma^{2}$. \\ \\
\noindent
The mean, $\mu$ can be given as :
\begin{align*}
    \mu = \int_{-\infty}^{\infty} f(x)x \,dx = 0 \text{  since } f(x) \text{ is an even function. }
\end{align*} 
The variance, $\sigma^{2}$ can therefore be given as :
\begin{align*}
    \sigma^{2} %&= \int_{-\infty}^{\infty} f(x)x^{2} \,dx \\
   % &= \int_{-\infty}^{\infty} ke^{-\frac{x^{2}}{2}}(a + b\cos(c\pi x))x^{2} \,dx \\
    &= k\left(a\int_{-\infty}^{\infty} e^{-\frac{x^{2}}{2}}x^{2}\,dx + b\int_{-\infty}^{\infty} e^{-\frac{x^{2}}{2}}x^{2}\cos(c\pi x) \,dx\right)
\end{align*}
Now, we note that $\int_{-\infty}^{\infty} e^{-\frac{x^{2}}{2}}x^{2}\,dx = \sqrt{2\pi}$ and $\int_{-\infty}^{\infty} e^{-\frac{x^{2}}{2}}x^{2}\cos(c\pi x) \,dx = e^{-\frac{c^{2}\pi^{2}}{2}}\sqrt{2\pi}(1 - c^{2}\pi^{2})$
Therefore, 
\begin{align*}
    \sigma^{2} &= k\left(a\sqrt{2\pi} + be^{-\frac{c^{2}\pi^{2}}{2}}\sqrt{2\pi}(1 - c^{2}\pi^{2})\right) 
   % &= \frac{1}{\sqrt{2\pi}(a + be^{-\frac{c^{2}\pi^{2}}{2}})}\left(a\sqrt{2\pi} + be^{-\frac{c^{2}\pi^{2}}{2}}\sqrt{2\pi}(1 - c^{2}\pi^{2})\right) \\
   % &= \frac{1}{(a + be^{-\frac{c^{2}\pi^{2}}{2}})}\left((a + be^{-\frac{c^{2}\pi^{2}}{2}}) - be^{-\frac{c^{2}\pi^{2}}{2}}c^{2}\pi^{2}\right) \\
    = 1 - \frac{bc^{2}\pi^{2}e^{-\frac{c^{2}\pi^{2}}{2}}}{a + be^{-\frac{c^{2}\pi^{2}}{2}}}
\end{align*}

The symmetrized CGF can therefore be written as 
\begin{align*}
    \sym CGF(t) &:= \ln(MGF_{f(x)}(t)) + \ln(MGF_{f(x)}(-t)) \\ 
    &= -2\ln(a + be^{-\frac{c^{2}\pi^{2}}{2}}) + t^2 + 2\ln(a + be^{-\frac{c^{2}\pi^{2}}{2}}\cos(c\pi t))
\end{align*}

Therefore, $g(t) = \sym CGF(t) - (1 - \frac{bc^{2}\pi^{2}e^{-\frac{c^{2}\pi^{2}}{2}}}{a + be^{-\frac{c^{2}\pi^{2}}{2}}})t^{2}$ can be written as : \begin{align*} g(t) &= -2\ln(a + be^{-\frac{c^{2}\pi^{2}}{2}}) + 2\ln(a + be^{-\frac{c^{2}\pi^{2}}{2}}\cos(c\pi t)) + \frac{bc^{2}\pi^{2}e^{-\frac{c^{2}\pi^{2}}{2}}}{a + be^{-\frac{c^{2}\pi^{2}}{2}}}t^2
\end{align*}

Finally, $g'(t)$ can be written as :
\begin{align*}
    g'(t) &= \frac{d}{dt}\left(2\ln(a + be^{-\frac{c^{2}\pi^{2}}{2}}\cos(c\pi t)) + \frac{bc^{2}\pi^{2}e^{-\frac{c^{2}\pi^{2}}{2}}}{a + be^{-\frac{c^{2}\pi^{2}}{2}}}t^2\right) \\
    % &= -\frac{2bc\pi}{a + be^{-\frac{c^{2}\pi^{2}}{2}}\cos(c\pi t)}.e^{-\frac{c^{2}\pi^{2}}{2}}.\sin(c\pi t) + \frac{2bc^{2}\pi^{2}e^{-\frac{c^{2}\pi^{2}}{2}}}{a + be^{-\frac{c^{2}\pi^{2}}{2}}}t \\
    &= -\frac{2bc\pi e^{-\frac{c^{2}\pi^{2}}{2}} \sin(c\pi t)}{a + be^{-\frac{c^{2}\pi^{2}}{2}}\cos(c\pi t)} + \frac{2bc^{2}\pi^{2}e^{-\frac{c^{2}\pi^{2}}{2}}}{a + be^{-\frac{c^{2}\pi^{2}}{2}}}t
\end{align*}

$g''(t)$ can be written as : 
\begin{align*}
    g''(t) % &=  -2bc\pi e^{-\frac{c^{2}\pi^{2}}{2}}\frac{(a + be^{-\frac{c^{2}\pi^{2}}{2}}\cos(c\pi t))c\pi \cos(c\pi t) - \sin(c\pi t)(-bc\pi e^{-\frac{c^{2}\pi^{2}}{2}}\sin(c\pi t))}{(a + be^{-\frac{c^{2}\pi^{2}}{2}}\cos(c\pi t))^{2}} + \frac{2bc^{2}\pi^{2}e^{-\frac{c^{2}\pi^{2}}{2}}}{a + be^{-\frac{c^{2}\pi^{2}}{2}}} \\
  %  &= -2bc\pi e^{-\frac{c^{2}\pi^{2}}{2}}\frac{ac\pi \cos(c\pi t) +  bc\pi e^{-\frac{c^{2}\pi^{2}}{2}} \left(\cos^{2}(c\pi t) + \sin^{2}(c\pi t)\right)}{(a + be^{-\frac{c^{2}\pi^{2}}{2}}\cos(c\pi t))^{2}} + \frac{2bc^{2}\pi^{2}e^{-\frac{c^{2}\pi^{2}}{2}}}{a + be^{-\frac{c^{2}\pi^{2}}{2}}} \\
    &= -2bc^{2}\pi^{2} e^{-\frac{c^{2}\pi^{2}}{2}}\frac{a \cos(c\pi t) +  b e^{-\frac{c^{2}\pi^{2}}{2}}}{(a + be^{-\frac{c^{2}\pi^{2}}{2}}\cos(c\pi t))^{2}} + \frac{2bc^{2}\pi^{2}e^{-\frac{c^{2}\pi^{2}}{2}}}{a + be^{-\frac{c^{2}\pi^{2}}{2}}}
\end{align*}

and $g'''(t)$ can be evaluated as:
\begin{align*}
% &= -2bc^{2}\pi^{2} e^{-\frac{c^{2}\pi^{2}}{2}}\frac{(a + be^{-\frac{c^{2}\pi^{2}}{2}}\cos(c\pi t))^{2}(-ac\pi \sin(c\pi t)) -  2(a \cos(c\pi t) +  b e^{-\frac{c^{2}\pi^{2}}{2}})(a + be^{-\frac{c^{2}\pi^{2}}{2}}\cos(c\pi t))(-bc\pi e^{-\frac{c^{2}\pi^{2}}{2}}\sin(c\pi t))}{(a + be^{-\frac{c^{2}\pi^{2}}{2}}\cos(c\pi t))^{4}} \\
%&= -2bc^{2}\pi^{2} e^{-\frac{c^{2}\pi^{2}}{2}}\frac{(a + be^{-\frac{c^{2}\pi^{2}}{2}}\cos(c\pi t))(-ac\pi \sin(c\pi t)) -  2(a \cos(c\pi t) +  b e^{-\frac{c^{2}\pi^{2}}{2}})(-bc\pi e^{-\frac{c^{2}\pi^{2}}{2}}\sin(c\pi t))}{(a + be^{-\frac{c^{2}\pi^{2}}{2}}\cos(c\pi t))^{3}} \\
% &= -2bc^{2}\pi^{2} e^{-\frac{c^{2}\pi^{2}}{2}}\frac{-a^{2}c\pi \sin(c\pi t) - abc\pi e^{-\frac{c^{2}\pi^{2}}{2}}\cos(c\pi t)\sin(c\pi t) + 2abc\pi e^{-\frac{c^{2}\pi^{2}}{2}}\cos(c\pi t)\sin(c\pi t) + 2b^{2}c\pi e^{-c^{2}\pi^{2}}\sin(c\pi t)  }{(a + be^{-\frac{c^{2}\pi^{2}}{2}}\cos(c\pi t))^{3}} \\
%&= 2bc^{2}\pi^{2} e^{-\frac{c^{2}\pi^{2}}{2}}\sin(c\pi t)\frac{\left(a^{2}c\pi - 2b^{2}c\pi e^{-c^{2}\pi^{2}} - abc\pi e^{-\frac{c^{2}\pi^{2}}{2}}\cos(c\pi t)\right) }{(a + be^{-\frac{c^{2}\pi^{2}}{2}}\cos(c\pi t))^{3}} \\
g'''(t)&= 2bc^{3}\pi^{3} e^{-\frac{c^{2}\pi^{2}}{2}}\sin(c\pi t)\frac{\left(a^{2} - 2b^{2} e^{-c^{2}\pi^{2}} - ab e^{-\frac{c^{2}\pi^{2}}{2}}\cos(c\pi t)\right) }{(a + be^{-\frac{c^{2}\pi^{2}}{2}}\cos(c\pi t))^{3}}
\end{align*}

Lets set $a = 2, b = -1, c = \frac{4}{\pi}$. Then, we have that the pdf, $f(x) = ke^{-\frac{x^{2}}{2}}(2 - \cos(4x)) = ke^{-\frac{x^{2}}{2}}(1 + 2\sin^{2}(2x))$. The corresponding functions are : 
\begin{align*}
    E[\exp(t X)] &= \frac{1}{(2 - e^{-8})}e^{\frac{t^{2}}{2}}\left(2 - e^{-8}\cos(4t)\right)\\
    g(t) &:=\sym CGF(t)-\var(X)t^2=  -2\ln(2 - e^{-8}) + 2\ln(2 - e^{-8}\cos(4t)) - \frac{16e^{-8}}{2 - e^{-8}}t^2\label{eq:contrastcos}\\
    g'(t) &= \frac{8 e^{-8} \sin(4t)}{2 - e^{-8}\cos(4t)} - \frac{32e^{-8}}{2 - e^{-8}}t\\
    g''(t) &= 32 e^{-8}\frac{2 \cos(4 t) -  e^{-8}}{(2 - e^{-8}\cos(4t))^{2}} - \frac{32e^{-8}}{2 - e^{-8}}\\
    g'''(t) &= -128 e^{-8}\sin(4t)\frac{\left(4 - 2e^{-16} + 2 e^{-8}\cos(4t)\right) }{(2 - e^{-8}\cos(4t))^{3}}
\end{align*}
% \begin{enumerate}
%     \item $MGF = \frac{1}{(2 - e^{-8})}e^{\frac{t^{2}}{2}}\left(2 - e^{-8}\cos(4t)\right)$ \\
%     \item $g(t) = -2\ln(2 - e^{-8}) + 2\ln(2 - e^{-8}\cos(4t)) - \frac{16e^{-8}}{2 - e^{-8}}t^2$ \\
%     \item $g'(t) = \frac{8 e^{-8} \sin(4t)}{2 - e^{-8}\cos(4t)} - \frac{32e^{-8}}{2 - e^{-8}}t$ \\
%     \item $g''(t) = 32 e^{-8}\frac{2 \cos(4 t) -  e^{-8}}{(2 - e^{-8}\cos(4t))^{2}} - \frac{32e^{-8}}{2 - e^{-8}}$ \\
%     \item $g'''(t) = -128 e^{-8}\sin(4t)\frac{\left(4 - 2e^{-16} + 2 e^{-8}\cos(4t)\right) }{(2 - e^{-8}\cos(4t))^{3}}$
% \end{enumerate}

 \textbf{Closeness to a Gaussian MGF:}
It can be seen there that $g'''(t)$ changes sign in the half-line. However, the key is to note that the MGF of this distribution is very close to that of a Gaussian and can be made arbitrarily small by varying the parameters $a$ and $b$. This renders the CGF-based contrast function ineffectual. This leads us to believe that Assumption 1(d) is not far from being necessary to ensure global optimality of the corresponding objective function.

% \textbf{Closeness to a Gaussian MGF}
% Note that contrast function $g(t)=0$ for a Gaussian distribution. Eq~\ref{eq:contrastcos} shows that the contrast function of the above random variable is also very cl


\section{Surface plots of the CHF-based contrast functions}
\label{section:surface_plots}
\begin{figure}
    \centering
    \begin{subfigure}[b]{0.4\textwidth}
        \centering
        \includegraphics[width=\textwidth]{surface_plots/surface_plot_uniform_1.png}
        \caption{\label{fig:uniform_1}}
    \end{subfigure}
    \begin{subfigure}[b]{0.4\textwidth}
        \centering
        \includegraphics[width=\textwidth]{surface_plots/surface_plot_uniform_2.png}
        \caption{\label{fig:uniform_2}}
    \end{subfigure}
    \begin{subfigure}[b]{0.4\textwidth}
        \centering
        \includegraphics[width=\textwidth]{surface_plots/surface_plot_zerok_1.png}
        \caption{\label{fig:zerok_1}}
    \end{subfigure}
    \begin{subfigure}[b]{0.4\textwidth}
        \centering
        \includegraphics[width=\textwidth]{surface_plots/surface_plot_zerok_2.png}
        \caption{\label{fig:zerok_2}}
    \end{subfigure}
    \vspace{30pt}
    \caption{Surface plots for zero-kurtosis (\ref{fig:zerok_1} and \ref{fig:zerok_2}) and Uniform (\ref{fig:uniform_1}, \ref{fig:uniform_2}) data with $n=10000$ points, noise-power $\rho = 0.1$ and number of source signals, $k = 2$.}
    \label{fig:surface_plots}
\end{figure}

Figure \ref{fig:surface_plots} depicts the loss landscape of the Characteristic function (CHF) based contrast function described in Section \ref{section:chf_cgf_fn} of the manuscript. We plot the value of the contrast function evaluated at $B^{-1}\bu$, $\bu = \frac{1}{\sqrt{x^2 + y^2}}\bb{\begin{matrix}
    x \\ y 
\end{matrix}}$ for $x,y \in [-1,1]$. As shown in the figure. 
We have rotated the data so that the columns of $B$ align with the $X$ and $Y$ axes. The global maxima occur at $\bu$ aligned with the columns of $B$.

%%%%%%%%%%%%%%%%%%%%%%%%%%%%%%%%%%%%%%%%%%%%%%%%%%%%%%%%%%%%



% \newpage
% \section*{NeurIPS Paper Checklist}

% %%% BEGIN INSTRUCTIONS %%%
% % The checklist is designed to encourage best practices for responsible machine learning research, addressing issues of reproducibility, transparency, research ethics, and societal impact. Do not remove the checklist: {\bf The papers not including the checklist will be desk rejected.} The checklist should follow the references and precede the (optional) supplemental material.  The checklist does NOT count towards the page
% % limit. 

% % Please read the checklist guidelines carefully for information on how to answer these questions. For each question in the checklist:
% % \begin{itemize}
% %     \item You should answer \answerYes{}, \answerNo{}, or \answerNA{}.
% %     \item \answerNA{} means either that the question is Not Applicable for that particular paper or the relevant information is Not Available.
% %     \item Please provide a short (1–2 sentence) justification right after your answer (even for NA). 
% %    % \item {\bf The papers not including the checklist will be desk rejected.}
% % \end{itemize}

% % {\bf The checklist answers are an integral part of your paper submission.} They are visible to the reviewers, area chairs, senior area chairs, and ethics reviewers. You will be asked to also include it (after eventual revisions) with the final version of your paper, and its final version will be published with the paper.

% % The reviewers of your paper will be asked to use the checklist as one of the factors in their evaluation. While "\answerYes{}" is generally preferable to "\answerNo{}", it is perfectly acceptable to answer "\answerNo{}" provided a proper justification is given (e.g., "error bars are not reported because it would be too computationally expensive" or "we were unable to find the license for the dataset we used"). In general, answering "\answerNo{}" or "\answerNA{}" is not grounds for rejection. While the questions are phrased in a binary way, we acknowledge that the true answer is often more nuanced, so please just use your best judgment and write a justification to elaborate. All supporting evidence can appear either in the main paper or the supplemental material, provided in appendix. If you answer \answerYes{} to a question, in the justification please point to the section(s) where related material for the question can be found.

% % IMPORTANT, please:
% % \begin{itemize}
% %     \item {\bf Delete this instruction block, but keep the section heading ``NeurIPS paper checklist"},
% %     \item  {\bf Keep the checklist subsection headings, questions/answers and guidelines below.}
% %     \item {\bf Do not modify the questions and only use the provided macros for your answers}.
% % \end{itemize} 
 

% %%% END INSTRUCTIONS %%%


% \begin{enumerate}

% \item {\bf Claims}
%     \item[] Question: Do the main claims made in the abstract and introduction accurately reflect the paper's contributions and scope?
%     \item[] Answer: \answerYes{} % Replace by \answerYes{}, \answerNo{}, or \answerNA{}.
%     \item[] Justification: Our main contributions are the Meta algorithm for non-parametric selection of the best noisy ICA algorithm for a given distribution, along with 2 new contrast functions of independent interest. 
%     \item[] Guidelines:
%     \begin{itemize}
%         \item The answer NA means that the abstract and introduction do not include the claims made in the paper.
%         \item The abstract and/or introduction should clearly state the claims made, including the contributions made in the paper and important assumptions and limitations. A No or NA answer to this question will not be perceived well by the reviewers. 
%         \item The claims made should match theoretical and experimental results, and reflect how much the results can be expected to generalize to other settings. 
%         \item It is fine to include aspirational goals as motivation as long as it is clear that these goals are not attained by the paper. 
%     \end{itemize}

% \item {\bf Limitations}
%     \item[] Question: Does the paper discuss the limitations of the work performed by the authors?
%     \item[] Answer: \answerYes{} % Replace by \answerYes{}, \answerNo{}, or \answerNA{}.
%     \item[] Justification: In the Experiments (Section~\ref{section:experiments}, there are some noise/sample-size regimes where other algorithms may perform slightly better than our proposed contrast functions, but we show that even then our diagnostic can pick out the best algorithm.
%     \item[] Guidelines:
%     \begin{itemize}
%         \item The answer NA means that the paper has no limitation while the answer No means that the paper has limitations, but those are not discussed in the paper. 
%         \item The authors are encouraged to create a separate "Limitations" section in their paper.
%         \item The paper should point out any strong assumptions and how robust the results are to violations of these assumptions (e.g., independence assumptions, noiseless settings, model well-specification, asymptotic approximations only holding locally). The authors should reflect on how these assumptions might be violated in practice and what the implications would be.
%         \item The authors should reflect on the scope of the claims made, e.g., if the approach was only tested on a few datasets or with a few runs. In general, empirical results often depend on implicit assumptions, which should be articulated.
%         \item The authors should reflect on the factors that influence the performance of the approach. For example, a facial recognition algorithm may perform poorly when image resolution is low or images are taken in low lighting. Or a speech-to-text system might not be used reliably to provide closed captions for online lectures because it fails to handle technical jargon.
%         \item The authors should discuss the computational efficiency of the proposed algorithms and how they scale with dataset size.
%         \item If applicable, the authors should discuss possible limitations of their approach to address problems of privacy and fairness.
%         \item While the authors might fear that complete honesty about limitations might be used by reviewers as grounds for rejection, a worse outcome might be that reviewers discover limitations that aren't acknowledged in the paper. The authors should use their best judgment and recognize that individual actions in favor of transparency play an important role in developing norms that preserve the integrity of the community. Reviewers will be specifically instructed to not penalize honesty concerning limitations.
%     \end{itemize}

% \item {\bf Theory Assumptions and Proofs}
%     \item[] Question: For each theoretical result, does the paper provide the full set of assumptions and a complete (and correct) proof?
%     \item[] Answer: \answerYes{} % Replace by \answerYes{}, \answerNo{}, or \answerNA{}.
%     \item[] Justification: Every theorem statement mentions the assumptions and the location of the proof in the Appendix.
%     \item[] Guidelines:
%     \begin{itemize}
%         \item The answer NA means that the paper does not include theoretical results. 
%         \item All the theorems, formulas, and proofs in the paper should be numbered and cross-referenced.
%         \item All assumptions should be clearly stated or referenced in the statement of any theorems.
%         \item The proofs can either appear in the main paper or the supplemental material, but if they appear in the supplemental material, the authors are encouraged to provide a short proof sketch to provide intuition. 
%         \item Inversely, any informal proof provided in the core of the paper should be complemented by formal proofs provided in appendix or supplemental material.
%         \item Theorems and Lemmas that the proof relies upon should be properly referenced. 
%     \end{itemize}

%     \item {\bf Experimental Result Reproducibility}
%     \item[] Question: Does the paper fully disclose all the information needed to reproduce the main experimental results of the paper to the extent that it affects the main claims and/or conclusions of the paper (regardless of whether the code and data are provided or not)?
%     \item[] Answer: \answerYes{} % Replace by \answerYes{}, \answerNo{}, or \answerNA{}.
%     \item[] Justification: Section~\ref{section:experiments} provides extensive details of the experimental setup.
%     \item[] Guidelines:
%     \begin{itemize}
%         \item The answer NA means that the paper does not include experiments.
%         \item If the paper includes experiments, a No answer to this question will not be perceived well by the reviewers: Making the paper reproducible is important, regardless of whether the code and data are provided or not.
%         \item If the contribution is a dataset and/or model, the authors should describe the steps taken to make their results reproducible or verifiable. 
%         \item Depending on the contribution, reproducibility can be accomplished in various ways. For example, if the contribution is a novel architecture, describing the architecture fully might suffice, or if the contribution is a specific model and empirical evaluation, it may be necessary to either make it possible for others to replicate the model with the same dataset, or provide access to the model. In general. releasing code and data is often one good way to accomplish this, but reproducibility can also be provided via detailed instructions for how to replicate the results, access to a hosted model (e.g., in the case of a large language model), releasing of a model checkpoint, or other means that are appropriate to the research performed.
%         \item While NeurIPS does not require releasing code, the conference does require all submissions to provide some reasonable avenue for reproducibility, which may depend on the nature of the contribution. For example
%         \begin{enumerate}
%             \item If the contribution is primarily a new algorithm, the paper should make it clear how to reproduce that algorithm.
%             \item If the contribution is primarily a new model architecture, the paper should describe the architecture clearly and fully.
%             \item If the contribution is a new model (e.g., a large language model), then there should either be a way to access this model for reproducing the results or a way to reproduce the model (e.g., with an open-source dataset or instructions for how to construct the dataset).
%             \item We recognize that reproducibility may be tricky in some cases, in which case authors are welcome to describe the particular way they provide for reproducibility. In the case of closed-source models, it may be that access to the model is limited in some way (e.g., to registered users), but it should be possible for other researchers to have some path to reproducing or verifying the results.
%         \end{enumerate}
%     \end{itemize}


% \item {\bf Open access to data and code}
%     \item[] Question: Does the paper provide open access to the data and code, with sufficient instructions to faithfully reproduce the main experimental results, as described in supplemental material?
%     \item[] Answer: \answerYes{} % Replace by \answerYes{}, \answerNo{}, or \answerNA{}.
%     \item[] Justification: We make the code for our experiments available.
%     \item[] Guidelines:
%     \begin{itemize}
%         \item The answer NA means that paper does not include experiments requiring code.
%         \item Please see the NeurIPS code and data submission guidelines (\url{https://nips.cc/public/guides/CodeSubmissionPolicy}) for more details.
%         \item While we encourage the release of code and data, we understand that this might not be possible, so “No” is an acceptable answer. Papers cannot be rejected simply for not including code, unless this is central to the contribution (e.g., for a new open-source benchmark).
%         \item The instructions should contain the exact command and environment needed to run to reproduce the results. See the NeurIPS code and data submission guidelines (\url{https://nips.cc/public/guides/CodeSubmissionPolicy}) for more details.
%         \item The authors should provide instructions on data access and preparation, including how to access the raw data, preprocessed data, intermediate data, and generated data, etc.
%         \item The authors should provide scripts to reproduce all experimental results for the new proposed method and baselines. If only a subset of experiments are reproducible, they should state which ones are omitted from the script and why.
%         \item At submission time, to preserve anonymity, the authors should release anonymized versions (if applicable).
%         \item Providing as much information as possible in supplemental material (appended to the paper) is recommended, but including URLs to data and code is permitted.
%     \end{itemize}


% \item {\bf Experimental Setting/Details}
%     \item[] Question: Does the paper specify all the training and test details (e.g., data splits, hyperparameters, how they were chosen, type of optimizer, etc.) necessary to understand the results?
%     \item[] Answer: \answerYes{} % Replace by \answerYes{}, \answerNo{}, or \answerNA{}.
%     \item[] Justification: Section~\ref{section:experiments} provides extensive details of the experimental setup.
%     \item[] Guidelines:
%     \begin{itemize}
%         \item The answer NA means that the paper does not include experiments.
%         \item The experimental setting should be presented in the core of the paper to a level of detail that is necessary to appreciate the results and make sense of them.
%         \item The full details can be provided either with the code, in appendix, or as supplemental material.
%     \end{itemize}

% \item {\bf Experiment Statistical Significance}
%     \item[] Question: Does the paper report error bars suitably and correctly defined or other appropriate information about the statistical significance of the experiments?
%     \item[] Answer: \answerYes{} % Replace by \answerYes{}, \answerNo{}, or \answerNA{}.
%     \item[] Justification: Figure~\ref{fig:amari_ind} and \ref{fig:plots_exp}(a) show error bars and error distribution respectively. Figure~\ref{fig:plots_exp}(b),(c) show the mean over 100 runs, similar to \cite{NIPS2015_89f03f7d}
%     \item[] Guidelines:
%     \begin{itemize}
%         \item The answer NA means that the paper does not include experiments.
%         \item The authors should answer "Yes" if the results are accompanied by error bars, confidence intervals, or statistical significance tests, at least for the experiments that support the main claims of the paper.
%         \item The factors of variability that the error bars are capturing should be clearly stated (for example, train/test split, initialization, random drawing of some parameter, or overall run with given experimental conditions).
%         \item The method for calculating the error bars should be explained (closed form formula, call to a library function, bootstrap, etc.)
%         \item The assumptions made should be given (e.g., Normally distributed errors).
%         \item It should be clear whether the error bar is the standard deviation or the standard error of the mean.
%         \item It is OK to report 1-sigma error bars, but one should state it. The authors should preferably report a 2-sigma error bar than state that they have a 96\% CI, if the hypothesis of Normality of errors is not verified.
%         \item For asymmetric distributions, the authors should be careful not to show in tables or figures symmetric error bars that would yield results that are out of range (e.g. negative error rates).
%         \item If error bars are reported in tables or plots, The authors should explain in the text how they were calculated and reference the corresponding figures or tables in the text.
%     \end{itemize}

% \item {\bf Experiments Compute Resources}
%     \item[] Question: For each experiment, does the paper provide sufficient information on the computer resources (type of compute workers, memory, time of execution) needed to reproduce the experiments?
%     \item[] Answer: \answerYes{} % Replace by \answerYes{}, \answerNo{}, or \answerNA{}.
%     \item[] Justification: Our experiments were performed on a single Macbook Pro M2 2022 CPU with 8 GB RAM.
%     \item[] Guidelines:
%     \begin{itemize}
%         \item The answer NA means that the paper does not include experiments.
%         \item The paper should indicate the type of compute workers CPU or GPU, internal cluster, or cloud provider, including relevant memory and storage.
%         \item The paper should provide the amount of compute required for each of the individual experimental runs as well as estimate the total compute. 
%         \item The paper should disclose whether the full research project required more compute than the experiments reported in the paper (e.g., preliminary or failed experiments that didn't make it into the paper). 
%     \end{itemize}
    
% \item {\bf Code Of Ethics}
%     \item[] Question: Does the research conducted in the paper conform, in every respect, with the NeurIPS Code of Ethics \url{https://neurips.cc/public/EthicsGuidelines}?
%     \item[] Answer: \answerYes{} % Replace by \answerYes{}, \answerNo{}, or \answerNA{}.
%     \item[] Justification: We abide by the NeurIPS Code of Ethics in our work.
%     \item[] Guidelines:
%     \begin{itemize}
%         \item The answer NA means that the authors have not reviewed the NeurIPS Code of Ethics.
%         \item If the authors answer No, they should explain the special circumstances that require a deviation from the Code of Ethics.
%         \item The authors should make sure to preserve anonymity (e.g., if there is a special consideration due to laws or regulations in their jurisdiction).
%     \end{itemize}


% \item {\bf Broader Impacts}
%     \item[] Question: Does the paper discuss both potential positive societal impacts and negative societal impacts of the work performed?
%     \item[] Answer: \answerNo{}% Replace by \answerYes{}, \answerNo{}, or \answerNA{}.
%     \item[] Justification: This paper presents work whose goal is to advance the field of Machine Learning. There are many potential societal consequences of our work, none of which we feel must be specifically highlighted here.
%     \item[] Guidelines:
%     \begin{itemize}
%         \item The answer NA means that there is no societal impact of the work performed.
%         \item If the authors answer NA or No, they should explain why their work has no societal impact or why the paper does not address societal impact.
%         \item Examples of negative societal impacts include potential malicious or unintended uses (e.g., disinformation, generating fake profiles, surveillance), fairness considerations (e.g., deployment of technologies that could make decisions that unfairly impact specific groups), privacy considerations, and security considerations.
%         \item The conference expects that many papers will be foundational research and not tied to particular applications, let alone deployments. However, if there is a direct path to any negative applications, the authors should point it out. For example, it is legitimate to point out that an improvement in the quality of generative models could be used to generate deepfakes for disinformation. On the other hand, it is not needed to point out that a generic algorithm for optimizing neural networks could enable people to train models that generate Deepfakes faster.
%         \item The authors should consider possible harms that could arise when the technology is being used as intended and functioning correctly, harms that could arise when the technology is being used as intended but gives incorrect results, and harms following from (intentional or unintentional) misuse of the technology.
%         \item If there are negative societal impacts, the authors could also discuss possible mitigation strategies (e.g., gated release of models, providing defenses in addition to attacks, mechanisms for monitoring misuse, mechanisms to monitor how a system learns from feedback over time, improving the efficiency and accessibility of ML).
%     \end{itemize}
    
% \item {\bf Safeguards}
%     \item[] Question: Does the paper describe safeguards that have been put in place for responsible release of data or models that have a high risk for misuse (e.g., pretrained language models, image generators, or scraped datasets)?
%     \item[] Answer: \answerNA{} % Replace by \answerYes{}, \answerNo{}, or \answerNA{}.
%     \item[] Justification: We do not release any models with a risk for misuse.
%     \item[] Guidelines:
%     \begin{itemize}
%         \item The answer NA means that the paper poses no such risks.
%         \item Released models that have a high risk for misuse or dual-use should be released with necessary safeguards to allow for controlled use of the model, for example by requiring that users adhere to usage guidelines or restrictions to access the model or implementing safety filters. 
%         \item Datasets that have been scraped from the Internet could pose safety risks. The authors should describe how they avoided releasing unsafe images.
%         \item We recognize that providing effective safeguards is challenging, and many papers do not require this, but we encourage authors to take this into account and make a best faith effort.
%     \end{itemize}

% \item {\bf Licenses for existing assets}
%     \item[] Question: Are the creators or original owners of assets (e.g., code, data, models), used in the paper, properly credited and are the license and terms of use explicitly mentioned and properly respected?
%     \item[] Answer: \answerYes{} % Replace by \answerYes{}, \answerNo{}, or \answerNA{}.
%     \item[] Justification: MATLAB implementations (under the GNU General Public License) can be found at - \href{https://research.ics.aalto.fi/ica/fastica/}{FastICA} and \href{https://github.com/ruohoruotsi/Riddim/blob/master/MATLAB/jade.m}{JADE}. The code for PFICA was provided on request by the authors.
%     \item[] Guidelines:
%     \begin{itemize}
%         \item The answer NA means that the paper does not use existing assets.
%         \item The authors should cite the original paper that produced the code package or dataset.
%         \item The authors should state which version of the asset is used and, if possible, include a URL.
%         \item The name of the license (e.g., CC-BY 4.0) should be included for each asset.
%         \item For scraped data from a particular source (e.g., website), the copyright and terms of service of that source should be provided.
%         \item If assets are released, the license, copyright information, and terms of use in the package should be provided. For popular datasets, \url{paperswithcode.com/datasets} has curated licenses for some datasets. Their licensing guide can help determine the license of a dataset.
%         \item For existing datasets that are re-packaged, both the original license and the license of the derived asset (if it has changed) should be provided.
%         \item If this information is not available online, the authors are encouraged to reach out to the asset's creators.
%     \end{itemize}

% \item {\bf New Assets}
%     \item[] Question: Are new assets introduced in the paper well documented and is the documentation provided alongside the assets?
%     \item[] Answer: \answerNA{} % Replace by \answerYes{}, \answerNo{}, or \answerNA{}.
%     \item[] Justification: The paper does not release new assets to require documentation or licensing.
%     \item[] Guidelines:
%     \begin{itemize}
%         \item The answer NA means that the paper does not release new assets.
%         \item Researchers should communicate the details of the dataset/code/model as part of their submissions via structured templates. This includes details about training, license, limitations, etc. 
%         \item The paper should discuss whether and how consent was obtained from people whose asset is used.
%         \item At submission time, remember to anonymize your assets (if applicable). You can either create an anonymized URL or include an anonymized zip file.
%     \end{itemize}

% \item {\bf Crowdsourcing and Research with Human Subjects}
%     \item[] Question: For crowdsourcing experiments and research with human subjects, does the paper include the full text of instructions given to participants and screenshots, if applicable, as well as details about compensation (if any)? 
%     \item[] Answer: \answerNA{}%\answerTODO{} % Replace by \answerYes{}, \answerNo{}, or \answerNA{}.
%     \item[] Justification:  We only use publicly available datasets in Section~\ref{section:experiments}.%\justificationTODO{}
%     \item[] Guidelines:
%     \begin{itemize}
%         \item The answer NA means that the paper does not involve crowdsourcing nor research with human subjects.
%         \item Including this information in the supplemental material is fine, but if the main contribution of the paper involves human subjects, then as much detail as possible should be included in the main paper. 
%         \item According to the NeurIPS Code of Ethics, workers involved in data collection, curation, or other labor should be paid at least the minimum wage in the country of the data collector. 
%     \end{itemize}

% \item {\bf Institutional Review Board (IRB) Approvals or Equivalent for Research with Human Subjects}
%     \item[] Question: Does the paper describe potential risks incurred by study participants, whether such risks were disclosed to the subjects, and whether Institutional Review Board (IRB) approvals (or an equivalent approval/review based on the requirements of your country or institution) were obtained?
%     \item[] Answer: \answerNA{} % Replace by \answerYes{}, \answerNo{}, or \answerNA{}.
%     \item[] Justification: We do not have experiments with crowdsourcing or human subjects.
%     \item[] Guidelines:
%     \begin{itemize}
%         \item The answer NA means that the paper does not involve crowdsourcing nor research with human subjects.
%         \item Depending on the country in which research is conducted, IRB approval (or equivalent) may be required for any human subjects research. If you obtained IRB approval, you should clearly state this in the paper. 
%         \item We recognize that the procedures for this may vary significantly between institutions and locations, and we expect authors to adhere to the NeurIPS Code of Ethics and the guidelines for their institution. 
%         \item For initial submissions, do not include any information that would break anonymity (if applicable), such as the institution conducting the review.
%     \end{itemize}

% \end{enumerate}


\end{document}
